\section{Twins, ladders and what remains beyond the reflection axis}
\label{sec:outlook}

The classification has a simple mechanism. At the ancestor midpoint, the Gram reflection
becomes a conjugacy that fixes the terminal column; the finite word is then recoverable from
a rigid succession of paired blocks, and palindromicity turns the paired blocks into an odd
common divisor of the completed counts. \Cref{prop:any-centre} shows that the constructive
half of this mechanism works around every palindrome. This section asks whether the block
exchange is the only source of equal labels among palindromes, records what an exact
computation says about it, and isolates the step that a proof would need.

\subsection{A conjecture on twins}

Call two distinct palindromes with the same Gram label \emph{twins}. By
\cref{prop:thue-morse}, twins are exactly the pairs of palindromes $av\neq av'$ without the
factor $bb$ for which the lengths $\abs{\operatorname{Pal}(\theta(av))}$ and
$\abs{\operatorname{Pal}(\theta(av'))}$ coincide. Reutenauer and Vuillon proved that the
Frobenius--Markov uniqueness conjecture is equivalent to the injectivity of this length on
the central words $\operatorname{Pal}(av)$~\cite[Theorem~1]{ReutenauerVuillon2017}. A search for
coincidences among all palindromes without $bb$ finds only words built around a palindromic
centre, and it suggests the following conjecture.

\begin{conjecture}[Structure of twins]
\label{conj:structure}
\textup{(i)}~A palindrome has a twin if and only if it is structured around some palindrome.
\textup{(ii)}~Two palindromes are twins if and only if a chain of block exchanges, each around
its own centre, leads from one to the other.
\end{conjecture}

Part (ii) implies part (i), and it is the form that \cref{rem:quadruple} imposes, since the
twins of a palindrome need not be its block exchanges. In both parts the \emph{if} direction is \cref{prop:any-centre}. Either part would
imply the Frobenius--Markov uniqueness conjecture. Indeed, two distinct Markov occurrences with
the same Markov number are twin central words, while a structured palindrome has completed
counts divisible by $2k+1\geqslant3$ and is therefore not central.

An exact enumeration by computer of all $16\,743\,538$ palindromes whose label does not exceed
$10^{26}$ supports the conjecture in its stronger form. Exactly $1\,719$ of these labels are
carried by more than one palindrome, $1\,695$ by two and $24$ by four. Every palindrome
involved is structured, every class of twins is connected by exchanges, and every exchange
pair has its root midpoint at its centre. Every class of four is a square of
\cref{cor:nested-square}, so that every pair of twins whose midpoint is shifted off its centres is
a diagonal of such a square, and the two palindromes of every class have the same numbers of
letters $a$ and $b$. The enumeration is part of the supplementary material.

\begin{figure}[htbp]
\centering
% Author: Dr. Denys Dutykh (Khalifa University of Science and Technology, Abu Dhabi, UAE)
% Generated by supplement/code/make_figures.py from exact data; do not edit by hand.

\begin{tikzpicture}[font=\footnotesize]
\draw[ink, line width=0.5pt] (0,7) -- (0,0) -- (12.4,0);
\draw[ink, line width=0.5pt] (0,1.489) -- (-0.08,1.489);
\node[ink, left] at (-0.08,1.489) {$10^{3}$};
\draw[ink, line width=0.5pt] (0,3.277) -- (-0.08,3.277);
\node[ink, left] at (-0.08,3.277) {$10^{6}$};
\draw[ink, line width=0.5pt] (0,5.064) -- (-0.08,5.064);
\node[ink, left] at (-0.08,5.064) {$10^{9}$};
\draw[ink, line width=0.5pt] (0,6.851) -- (-0.08,6.851);
\node[ink, left] at (-0.08,6.851) {$10^{12}$};
\draw[ink, line width=0.5pt] (1.167,0) -- (1.167,-0.08);
\node[ink, below] at (1.167,-0.08) {$0.385$};
\draw[ink, line width=0.5pt] (3.089,0) -- (3.089,-0.08);
\node[ink, below] at (3.089,-0.08) {$0.390$};
\draw[ink, line width=0.5pt] (5.012,0) -- (5.012,-0.08);
\node[ink, below] at (5.012,-0.08) {$0.395$};
\draw[ink, line width=0.5pt] (6.935,0) -- (6.935,-0.08);
\node[ink, below] at (6.935,-0.08) {$0.400$};
\draw[ink, line width=0.5pt] (8.857,0) -- (8.857,-0.08);
\node[ink, below] at (8.857,-0.08) {$0.405$};
\draw[ink, line width=0.5pt] (10.78,0) -- (10.78,-0.08);
\node[ink, below] at (10.78,-0.08) {$0.410$};
\node[ink, below] at (6.2,-0.48) {$\rho(z)/M(z)$, the frequency of $b$ in $\Omega_z$};
\node[ink, rotate=90, above] at (-0.75,3.5) {Gram label $M(z)$};
\fill[slate] (12.4,6.082) circle (0.28pt);
\fill[slate] (12.4,6.33) circle (0.28pt);
\fill[slate] (12.4,6.533) circle (0.28pt);
\fill[slate] (12.4,6.579) circle (0.28pt);
\fill[slate] (12.4,6.828) circle (0.28pt);
\fill[slate] (12.4,5.169) circle (0.28pt);
\fill[slate] (12.4,5.418) circle (0.28pt);
\fill[slate] (12.4,5.621) circle (0.28pt);
\fill[slate] (12.4,6.077) circle (0.28pt);
\fill[slate] (12.4,6.32) circle (0.28pt);
\fill[slate] (12.4,6.533) circle (0.28pt);
\fill[slate] (12.4,6.568) circle (0.28pt);
\fill[slate] (12.4,6.817) circle (0.28pt);
\fill[slate] (12.4,5.667) circle (0.28pt);
\fill[slate] (12.4,5.916) circle (0.28pt);
\fill[slate] (12.4,6.117) circle (0.28pt);
\fill[slate] (12.4,6.573) circle (0.28pt);
\fill[slate] (12.4,6.817) circle (0.28pt);
\fill[slate] (12.4,6.165) circle (0.28pt);
\fill[slate] (12.4,6.414) circle (0.28pt);
\fill[slate] (12.4,6.615) circle (0.28pt);
\fill[slate] (12.4,6.663) circle (0.28pt);
\fill[slate] (12.4,4.257) circle (0.28pt);
\fill[slate] (12.4,4.506) circle (0.28pt);
\fill[slate] (12.4,4.709) circle (0.28pt);
\fill[slate] (12.4,5.165) circle (0.28pt);
\fill[slate] (12.4,5.408) circle (0.28pt);
\fill[slate] (12.4,6.32) circle (0.28pt);
\fill[slate] (12.4,6.568) circle (0.28pt);
\fill[slate] (12.4,6.817) circle (0.28pt);
\fill[slate] (12.4,5.656) circle (0.28pt);
\fill[slate] (12.4,5.905) circle (0.28pt);
\fill[slate] (12.4,6.108) circle (0.28pt);
\fill[slate] (12.4,6.563) circle (0.28pt);
\fill[slate] (12.4,6.807) circle (0.28pt);
\fill[slate] (12.4,6.154) circle (0.28pt);
\fill[slate] (12.4,6.403) circle (0.28pt);
\fill[slate] (12.4,6.604) circle (0.28pt);
\fill[slate] (12.4,6.652) circle (0.28pt);
\fill[slate] (12.4,4.755) circle (0.28pt);
\fill[slate] (12.4,5.004) circle (0.28pt);
\fill[slate] (12.4,5.205) circle (0.28pt);
\fill[slate] (12.4,5.661) circle (0.28pt);
\fill[slate] (12.4,5.904) circle (0.28pt);
\fill[slate] (12.4,6.117) circle (0.28pt);
\fill[slate] (12.4,6.573) circle (0.28pt);
\fill[slate] (12.4,6.816) circle (0.28pt);
\fill[slate] (12.4,6.153) circle (0.28pt);
\fill[slate] (12.4,6.402) circle (0.28pt);
\fill[slate] (12.4,6.604) circle (0.28pt);
\fill[slate] (12.4,6.651) circle (0.28pt);
\fill[slate] (12.4,5.253) circle (0.28pt);
\fill[slate] (12.4,5.502) circle (0.28pt);
\fill[slate] (12.4,5.703) circle (0.28pt);
\fill[slate] (12.4,6.159) circle (0.28pt);
\fill[slate] (12.4,6.402) circle (0.28pt);
\fill[slate] (12.4,6.615) circle (0.28pt);
\fill[slate] (12.4,6.65) circle (0.28pt);
\fill[slate] (12.4,5.751) circle (0.28pt);
\fill[slate] (12.4,6) circle (0.28pt);
\fill[slate] (12.4,6.201) circle (0.28pt);
\fill[slate] (12.4,6.657) circle (0.28pt);
\fill[slate] (12.4,6.249) circle (0.28pt);
\fill[slate] (12.4,6.498) circle (0.28pt);
\fill[slate] (12.4,6.699) circle (0.28pt);
\fill[slate] (12.4,6.747) circle (0.28pt);
\fill[slate] (12.4,3.345) circle (0.28pt);
\fill[slate] (12.4,3.594) circle (0.28pt);
\fill[slate] (12.4,3.796) circle (0.28pt);
\fill[slate] (12.4,4.496) circle (0.28pt);
\fill[slate] (12.4,5.408) circle (0.28pt);
\fill[slate] (12.4,6.077) circle (0.28pt);
\fill[slate] (12.4,6.533) circle (0.28pt);
\fill[slate] (12.4,6.568) circle (0.28pt);
\fill[slate] (12.4,6.817) circle (0.28pt);
\fill[slate] (12.4,5.656) circle (0.28pt);
\fill[slate] (12.4,5.905) circle (0.28pt);
\fill[slate] (12.4,6.107) circle (0.28pt);
\fill[slate] (12.4,6.563) circle (0.28pt);
\fill[slate] (12.4,6.807) circle (0.28pt);
\fill[slate] (12.4,6.154) circle (0.28pt);
\fill[slate] (12.4,6.403) circle (0.28pt);
\fill[slate] (12.4,6.604) circle (0.28pt);
\fill[slate] (12.4,6.652) circle (0.28pt);
\fill[slate] (12.4,4.744) circle (0.28pt);
\fill[slate] (12.4,4.993) circle (0.28pt);
\fill[slate] (12.4,5.195) circle (0.28pt);
\fill[slate] (12.4,5.651) circle (0.28pt);
\fill[slate] (12.4,5.895) circle (0.28pt);
\fill[slate] (12.4,6.107) circle (0.28pt);
\fill[slate] (12.4,6.563) circle (0.28pt);
\fill[slate] (12.4,6.807) circle (0.28pt);
\fill[slate] (12.4,6.143) circle (0.28pt);
\fill[slate] (12.4,6.392) circle (0.28pt);
\fill[slate] (12.4,6.594) circle (0.28pt);
\fill[slate] (12.4,6.641) circle (0.28pt);
\fill[slate] (12.4,5.242) circle (0.28pt);
\fill[slate] (12.4,5.491) circle (0.28pt);
\fill[slate] (12.4,5.692) circle (0.28pt);
\fill[slate] (12.4,6.148) circle (0.28pt);
\fill[slate] (12.4,6.391) circle (0.28pt);
\fill[slate] (12.4,6.604) circle (0.28pt);
\fill[slate] (12.4,6.64) circle (0.28pt);
\fill[slate] (12.4,5.74) circle (0.28pt);
\fill[slate] (12.4,5.989) circle (0.28pt);
\fill[slate] (12.4,6.19) circle (0.28pt);
\fill[slate] (12.4,6.645) circle (0.28pt);
\fill[slate] (12.4,6.238) circle (0.28pt);
\fill[slate] (12.4,6.487) circle (0.28pt);
\fill[slate] (12.4,6.688) circle (0.28pt);
\fill[slate] (12.4,6.736) circle (0.28pt);
\fill[slate] (12.4,3.843) circle (0.28pt);
\fill[slate] (12.4,4.092) circle (0.28pt);
\fill[slate] (12.4,4.293) circle (0.28pt);
\fill[slate] (12.4,4.749) circle (0.28pt);
\fill[slate] (12.4,4.992) circle (0.28pt);
\fill[slate] (12.4,5.205) circle (0.28pt);
\fill[slate] (12.4,5.661) circle (0.28pt);
\fill[slate] (12.4,5.904) circle (0.28pt);
\fill[slate] (12.4,6.117) circle (0.28pt);
\fill[slate] (12.4,6.573) circle (0.28pt);
\fill[slate] (12.4,6.816) circle (0.28pt);
\fill[slate] (12.4,6.152) circle (0.28pt);
\fill[slate] (12.4,6.401) circle (0.28pt);
\fill[slate] (12.4,6.604) circle (0.28pt);
\fill[slate] (12.4,6.65) circle (0.28pt);
\fill[slate] (12.4,5.24) circle (0.28pt);
\fill[slate] (12.4,5.489) circle (0.28pt);
\fill[slate] (12.4,5.692) circle (0.28pt);
\fill[slate] (12.4,6.148) circle (0.28pt);
\fill[slate] (12.4,6.391) circle (0.28pt);
\fill[slate] (12.4,6.604) circle (0.28pt);
\fill[slate] (12.4,6.639) circle (0.28pt);
\fill[slate] (12.4,5.738) circle (0.28pt);
\fill[slate] (12.4,5.987) circle (0.28pt);
\fill[slate] (12.4,6.188) circle (0.28pt);
\fill[slate] (12.4,6.644) circle (0.28pt);
\fill[slate] (12.4,6.236) circle (0.28pt);
\fill[slate] (12.4,6.485) circle (0.28pt);
\fill[slate] (12.4,6.686) circle (0.28pt);
\fill[slate] (12.4,6.734) circle (0.28pt);
\fill[slate] (12.4,4.341) circle (0.28pt);
\fill[slate] (12.4,4.59) circle (0.28pt);
\fill[slate] (12.4,4.791) circle (0.28pt);
\fill[slate] (12.4,5.246) circle (0.28pt);
\fill[slate] (12.4,5.49) circle (0.28pt);
\fill[slate] (12.4,5.703) circle (0.28pt);
\fill[slate] (12.4,6.159) circle (0.28pt);
\fill[slate] (12.4,6.402) circle (0.28pt);
\fill[slate] (12.4,6.615) circle (0.28pt);
\fill[slate] (12.4,6.65) circle (0.28pt);
\fill[slate] (12.4,5.738) circle (0.28pt);
\fill[slate] (12.4,5.987) circle (0.28pt);
\fill[slate] (12.4,6.19) circle (0.28pt);
\fill[slate] (12.4,6.645) circle (0.28pt);
\fill[slate] (12.4,6.236) circle (0.28pt);
\fill[slate] (12.4,6.485) circle (0.28pt);
\fill[slate] (12.4,6.686) circle (0.28pt);
\fill[slate] (12.4,6.734) circle (0.28pt);
\fill[slate] (12.4,4.839) circle (0.28pt);
\fill[slate] (12.4,5.088) circle (0.28pt);
\fill[slate] (12.4,5.289) circle (0.28pt);
\fill[slate] (12.4,5.744) circle (0.28pt);
\fill[slate] (12.4,5.988) circle (0.28pt);
\fill[slate] (12.4,6.2) circle (0.28pt);
\fill[slate] (12.4,6.657) circle (0.28pt);
\fill[slate] (12.4,6.236) circle (0.28pt);
\fill[slate] (12.4,6.485) circle (0.28pt);
\fill[slate] (12.4,6.688) circle (0.28pt);
\fill[slate] (12.4,6.734) circle (0.28pt);
\fill[slate] (12.4,5.337) circle (0.28pt);
\fill[slate] (12.4,5.586) circle (0.28pt);
\fill[slate] (12.4,5.787) circle (0.28pt);
\fill[slate] (12.4,6.242) circle (0.28pt);
\fill[slate] (12.4,6.486) circle (0.28pt);
\fill[slate] (12.4,6.699) circle (0.28pt);
\fill[slate] (12.4,6.734) circle (0.28pt);
\fill[slate] (12.4,5.835) circle (0.28pt);
\fill[slate] (12.4,6.084) circle (0.28pt);
\fill[slate] (12.4,6.285) circle (0.28pt);
\fill[slate] (12.4,6.74) circle (0.28pt);
\fill[slate] (12.4,6.333) circle (0.28pt);
\fill[slate] (12.4,6.582) circle (0.28pt);
\fill[slate] (12.4,6.783) circle (0.28pt);
\fill[slate] (12.4,6.831) circle (0.28pt);
\fill[slate] (12.389,2.433) circle (0.28pt);
\fill[slate] (12.389,2.682) circle (0.28pt);
\fill[slate] (12.391,3.584) circle (0.28pt);
\fill[slate] (12.391,4.708) circle (0.28pt);
\fill[slate] (12.391,5.165) circle (0.28pt);
\fill[slate] (12.391,5.621) circle (0.28pt);
\fill[slate] (12.391,6.077) circle (0.28pt);
\fill[slate] (12.391,6.32) circle (0.28pt);
\fill[slate] (12.391,6.533) circle (0.28pt);
\fill[slate] (12.391,6.568) circle (0.28pt);
\fill[slate] (12.391,6.817) circle (0.28pt);
\fill[slate] (12.391,5.656) circle (0.28pt);
\fill[slate] (12.391,5.905) circle (0.28pt);
\fill[slate] (12.391,6.107) circle (0.28pt);
\fill[slate] (12.391,6.807) circle (0.28pt);
\fill[slate] (12.391,6.154) circle (0.28pt);
\fill[slate] (12.391,6.403) circle (0.28pt);
\fill[slate] (12.391,6.604) circle (0.28pt);
\fill[slate] (12.391,6.652) circle (0.28pt);
\fill[slate] (12.391,4.744) circle (0.28pt);
\fill[slate] (12.391,4.993) circle (0.28pt);
\fill[slate] (12.391,5.895) circle (0.28pt);
\fill[slate] (12.391,6.563) circle (0.28pt);
\fill[slate] (12.391,6.143) circle (0.28pt);
\fill[slate] (12.391,6.392) circle (0.28pt);
\fill[slate] (12.391,6.594) circle (0.28pt);
\fill[slate] (12.391,6.641) circle (0.28pt);
\fill[slate] (12.391,5.242) circle (0.28pt);
\fill[slate] (12.391,5.491) circle (0.28pt);
\fill[slate] (12.391,5.692) circle (0.28pt);
\fill[slate] (12.391,6.147) circle (0.28pt);
\fill[slate] (12.391,6.391) circle (0.28pt);
\fill[slate] (12.391,6.603) circle (0.28pt);
\fill[slate] (12.391,6.639) circle (0.28pt);
\fill[slate] (12.391,5.74) circle (0.28pt);
\fill[slate] (12.391,5.989) circle (0.28pt);
\fill[slate] (12.391,6.189) circle (0.28pt);
\fill[slate] (12.391,6.645) circle (0.28pt);
\fill[slate] (12.391,6.238) circle (0.28pt);
\fill[slate] (12.391,6.487) circle (0.28pt);
\fill[slate] (12.391,6.687) circle (0.28pt);
\fill[slate] (12.391,6.736) circle (0.28pt);
\fill[slate] (12.391,3.832) circle (0.28pt);
\fill[slate] (12.391,4.081) circle (0.28pt);
\fill[slate] (12.391,4.283) circle (0.28pt);
\fill[slate] (12.391,4.739) circle (0.28pt);
\fill[slate] (12.391,4.983) circle (0.28pt);
\fill[slate] (12.391,5.651) circle (0.28pt);
\fill[slate] (12.391,5.895) circle (0.28pt);
\fill[slate] (12.391,6.107) circle (0.28pt);
\fill[slate] (12.391,6.563) circle (0.28pt);
\fill[slate] (12.391,6.807) circle (0.28pt);
\fill[slate] (12.391,6.143) circle (0.28pt);
\fill[slate] (12.391,6.392) circle (0.28pt);
\fill[slate] (12.391,6.594) circle (0.28pt);
\fill[slate] (12.391,6.641) circle (0.28pt);
\fill[slate] (12.391,5.231) circle (0.28pt);
\fill[slate] (12.391,5.48) circle (0.28pt);
\fill[slate] (12.391,5.682) circle (0.28pt);
\fill[slate] (12.391,6.138) circle (0.28pt);
\fill[slate] (12.391,6.382) circle (0.28pt);
\fill[slate] (12.391,6.594) circle (0.28pt);
\fill[slate] (12.391,6.63) circle (0.28pt);
\fill[slate] (12.391,5.729) circle (0.28pt);
\fill[slate] (12.391,5.978) circle (0.28pt);
\fill[slate] (12.391,6.179) circle (0.28pt);
\fill[slate] (12.391,6.635) circle (0.28pt);
\fill[slate] (12.391,6.227) circle (0.28pt);
\fill[slate] (12.391,6.476) circle (0.28pt);
\fill[slate] (12.391,6.676) circle (0.28pt);
\fill[slate] (12.391,6.725) circle (0.28pt);
\fill[slate] (12.391,4.33) circle (0.28pt);
\fill[slate] (12.391,4.579) circle (0.28pt);
\fill[slate] (12.391,4.78) circle (0.28pt);
\fill[slate] (12.391,5.236) circle (0.28pt);
\fill[slate] (12.391,5.479) circle (0.28pt);
\fill[slate] (12.391,5.692) circle (0.28pt);
\fill[slate] (12.391,6.148) circle (0.28pt);
\fill[slate] (12.391,6.391) circle (0.28pt);
\fill[slate] (12.391,6.604) circle (0.28pt);
\fill[slate] (12.391,6.639) circle (0.28pt);
\fill[slate] (12.391,5.727) circle (0.28pt);
\fill[slate] (12.391,5.976) circle (0.28pt);
\fill[slate] (12.391,6.179) circle (0.28pt);
\fill[slate] (12.391,6.635) circle (0.28pt);
\fill[slate] (12.391,6.225) circle (0.28pt);
\fill[slate] (12.391,6.474) circle (0.28pt);
\fill[slate] (12.391,6.675) circle (0.28pt);
\fill[slate] (12.391,6.723) circle (0.28pt);
\fill[slate] (12.391,4.828) circle (0.28pt);
\fill[slate] (12.391,5.077) circle (0.28pt);
\fill[slate] (12.391,5.277) circle (0.28pt);
\fill[slate] (12.391,5.733) circle (0.28pt);
\fill[slate] (12.391,5.977) circle (0.28pt);
\fill[slate] (12.391,6.189) circle (0.28pt);
\fill[slate] (12.391,6.645) circle (0.28pt);
\fill[slate] (12.391,6.225) circle (0.28pt);
\fill[slate] (12.391,6.474) circle (0.28pt);
\fill[slate] (12.391,6.676) circle (0.28pt);
\fill[slate] (12.391,6.723) circle (0.28pt);
\fill[slate] (12.391,5.326) circle (0.28pt);
\fill[slate] (12.391,5.575) circle (0.28pt);
\fill[slate] (12.391,5.775) circle (0.28pt);
\fill[slate] (12.391,6.231) circle (0.28pt);
\fill[slate] (12.391,6.475) circle (0.28pt);
\fill[slate] (12.391,6.687) circle (0.28pt);
\fill[slate] (12.391,6.723) circle (0.28pt);
\fill[slate] (12.391,5.824) circle (0.28pt);
\fill[slate] (12.391,6.073) circle (0.28pt);
\fill[slate] (12.391,6.273) circle (0.28pt);
\fill[slate] (12.391,6.729) circle (0.28pt);
\fill[slate] (12.391,6.322) circle (0.28pt);
\fill[slate] (12.391,6.571) circle (0.28pt);
\fill[slate] (12.391,6.771) circle (0.28pt);
\fill[slate] (12.391,6.82) circle (0.28pt);
\fill[slate] (12.389,2.931) circle (0.28pt);
\fill[slate] (12.389,3.18) circle (0.28pt);
\fill[slate] (12.389,3.381) circle (0.28pt);
\fill[slate] (12.389,3.837) circle (0.28pt);
\fill[slate] (12.389,4.08) circle (0.28pt);
\fill[slate] (12.389,4.293) circle (0.28pt);
\fill[slate] (12.389,4.749) circle (0.28pt);
\fill[slate] (12.389,4.992) circle (0.28pt);
\fill[slate] (12.389,5.205) circle (0.28pt);
\fill[slate] (12.389,5.661) circle (0.28pt);
\fill[slate] (12.389,5.904) circle (0.28pt);
\fill[slate] (12.389,6.117) circle (0.28pt);
\fill[slate] (12.389,6.573) circle (0.28pt);
\fill[slate] (12.389,6.816) circle (0.28pt);
\fill[slate] (12.389,6.152) circle (0.28pt);
\fill[slate] (12.389,6.401) circle (0.28pt);
\fill[slate] (12.389,6.604) circle (0.28pt);
\fill[slate] (12.389,6.65) circle (0.28pt);
\fill[slate] (12.389,5.24) circle (0.28pt);
\fill[slate] (12.389,5.489) circle (0.28pt);
\fill[slate] (12.389,5.691) circle (0.28pt);
\fill[slate] (12.389,6.147) circle (0.28pt);
\fill[slate] (12.389,6.391) circle (0.28pt);
\fill[slate] (12.389,6.603) circle (0.28pt);
\fill[slate] (12.389,6.639) circle (0.28pt);
\fill[slate] (12.389,5.738) circle (0.28pt);
\fill[slate] (12.389,5.987) circle (0.28pt);
\fill[slate] (12.389,6.188) circle (0.28pt);
\fill[slate] (12.389,6.644) circle (0.28pt);
\fill[slate] (12.389,6.236) circle (0.28pt);
\fill[slate] (12.389,6.485) circle (0.28pt);
\fill[slate] (12.389,6.686) circle (0.28pt);
\fill[slate] (12.389,6.734) circle (0.28pt);
\fill[slate] (12.389,4.328) circle (0.28pt);
\fill[slate] (12.389,4.577) circle (0.28pt);
\fill[slate] (12.389,4.78) circle (0.28pt);
\fill[slate] (12.389,5.236) circle (0.28pt);
\fill[slate] (12.389,5.479) circle (0.28pt);
\fill[slate] (12.389,5.692) circle (0.28pt);
\fill[slate] (12.389,6.148) circle (0.28pt);
\fill[slate] (12.389,6.391) circle (0.28pt);
\fill[slate] (12.389,6.604) circle (0.28pt);
\fill[slate] (12.389,6.639) circle (0.28pt);
\fill[slate] (12.389,5.727) circle (0.28pt);
\fill[slate] (12.389,5.976) circle (0.28pt);
\fill[slate] (12.389,6.179) circle (0.28pt);
\fill[slate] (12.389,6.635) circle (0.28pt);
\fill[slate] (12.389,6.225) circle (0.28pt);
\fill[slate] (12.389,6.474) circle (0.28pt);
\fill[slate] (12.389,6.675) circle (0.28pt);
\fill[slate] (12.389,6.723) circle (0.28pt);
\fill[slate] (12.389,4.826) circle (0.28pt);
\fill[slate] (12.389,5.075) circle (0.28pt);
\fill[slate] (12.389,5.276) circle (0.28pt);
\fill[slate] (12.389,5.732) circle (0.28pt);
\fill[slate] (12.389,5.976) circle (0.28pt);
\fill[slate] (12.389,6.188) circle (0.28pt);
\fill[slate] (12.389,6.644) circle (0.28pt);
\fill[slate] (12.389,6.224) circle (0.28pt);
\fill[slate] (12.389,6.473) circle (0.28pt);
\fill[slate] (12.389,6.675) circle (0.28pt);
\fill[slate] (12.389,6.722) circle (0.28pt);
\fill[slate] (12.389,5.324) circle (0.28pt);
\fill[slate] (12.389,5.573) circle (0.28pt);
\fill[slate] (12.389,5.774) circle (0.28pt);
\fill[slate] (12.389,6.23) circle (0.28pt);
\fill[slate] (12.389,6.473) circle (0.28pt);
\fill[slate] (12.389,6.686) circle (0.28pt);
\fill[slate] (12.389,6.722) circle (0.28pt);
\fill[slate] (12.389,5.822) circle (0.28pt);
\fill[slate] (12.389,6.071) circle (0.28pt);
\fill[slate] (12.389,6.272) circle (0.28pt);
\fill[slate] (12.389,6.728) circle (0.28pt);
\fill[slate] (12.389,6.32) circle (0.28pt);
\fill[slate] (12.389,6.569) circle (0.28pt);
\fill[slate] (12.389,6.77) circle (0.28pt);
\fill[slate] (12.389,6.818) circle (0.28pt);
\fill[slate] (12.389,3.429) circle (0.28pt);
\fill[slate] (12.389,3.678) circle (0.28pt);
\fill[slate] (12.389,3.878) circle (0.28pt);
\fill[slate] (12.389,4.334) circle (0.28pt);
\fill[slate] (12.389,4.578) circle (0.28pt);
\fill[slate] (12.389,4.79) circle (0.28pt);
\fill[slate] (12.389,5.246) circle (0.28pt);
\fill[slate] (12.389,5.49) circle (0.28pt);
\fill[slate] (12.389,5.703) circle (0.28pt);
\fill[slate] (12.389,6.159) circle (0.28pt);
\fill[slate] (12.389,6.402) circle (0.28pt);
\fill[slate] (12.389,6.615) circle (0.28pt);
\fill[slate] (12.389,6.65) circle (0.28pt);
\fill[slate] (12.389,5.738) circle (0.28pt);
\fill[slate] (12.389,5.987) circle (0.28pt);
\fill[slate] (12.389,6.189) circle (0.28pt);
\fill[slate] (12.389,6.645) circle (0.28pt);
\fill[slate] (12.389,6.236) circle (0.28pt);
\fill[slate] (12.389,6.485) circle (0.28pt);
\fill[slate] (12.389,6.686) circle (0.28pt);
\fill[slate] (12.389,6.734) circle (0.28pt);
\fill[slate] (12.389,4.826) circle (0.28pt);
\fill[slate] (12.389,5.075) circle (0.28pt);
\fill[slate] (12.389,5.277) circle (0.28pt);
\fill[slate] (12.389,5.733) circle (0.28pt);
\fill[slate] (12.389,5.977) circle (0.28pt);
\fill[slate] (12.389,6.189) circle (0.28pt);
\fill[slate] (12.389,6.645) circle (0.28pt);
\fill[slate] (12.389,6.225) circle (0.28pt);
\fill[slate] (12.389,6.474) circle (0.28pt);
\fill[slate] (12.389,6.676) circle (0.28pt);
\fill[slate] (12.389,6.723) circle (0.28pt);
\fill[slate] (12.389,5.324) circle (0.28pt);
\fill[slate] (12.389,5.573) circle (0.28pt);
\fill[slate] (12.389,5.774) circle (0.28pt);
\fill[slate] (12.389,6.23) circle (0.28pt);
\fill[slate] (12.389,6.473) circle (0.28pt);
\fill[slate] (12.389,6.686) circle (0.28pt);
\fill[slate] (12.389,6.722) circle (0.28pt);
\fill[slate] (12.389,5.822) circle (0.28pt);
\fill[slate] (12.389,6.071) circle (0.28pt);
\fill[slate] (12.389,6.272) circle (0.28pt);
\fill[slate] (12.389,6.727) circle (0.28pt);
\fill[slate] (12.389,6.32) circle (0.28pt);
\fill[slate] (12.389,6.569) circle (0.28pt);
\fill[slate] (12.389,6.77) circle (0.28pt);
\fill[slate] (12.389,6.818) circle (0.28pt);
\fill[slate] (12.389,3.927) circle (0.28pt);
\fill[slate] (12.389,4.176) circle (0.28pt);
\fill[slate] (12.389,4.376) circle (0.28pt);
\fill[slate] (12.389,4.832) circle (0.28pt);
\fill[slate] (12.389,5.076) circle (0.28pt);
\fill[slate] (12.389,5.288) circle (0.28pt);
\fill[slate] (12.389,5.744) circle (0.28pt);
\fill[slate] (12.389,5.988) circle (0.28pt);
\fill[slate] (12.389,6.201) circle (0.28pt);
\fill[slate] (12.389,6.657) circle (0.28pt);
\fill[slate] (12.389,6.236) circle (0.28pt);
\fill[slate] (12.389,6.485) circle (0.28pt);
\fill[slate] (12.389,6.687) circle (0.28pt);
\fill[slate] (12.389,6.734) circle (0.28pt);
\fill[slate] (12.389,5.324) circle (0.28pt);
\fill[slate] (12.389,5.573) circle (0.28pt);
\fill[slate] (12.389,5.775) circle (0.28pt);
\fill[slate] (12.389,6.231) circle (0.28pt);
\fill[slate] (12.389,6.475) circle (0.28pt);
\fill[slate] (12.389,6.687) circle (0.28pt);
\fill[slate] (12.389,6.723) circle (0.28pt);
\fill[slate] (12.389,5.822) circle (0.28pt);
\fill[slate] (12.389,6.071) circle (0.28pt);
\fill[slate] (12.389,6.272) circle (0.28pt);
\fill[slate] (12.389,6.728) circle (0.28pt);
\fill[slate] (12.389,6.32) circle (0.28pt);
\fill[slate] (12.389,6.569) circle (0.28pt);
\fill[slate] (12.389,6.77) circle (0.28pt);
\fill[slate] (12.389,6.818) circle (0.28pt);
\fill[slate] (12.389,4.425) circle (0.28pt);
\fill[slate] (12.389,4.674) circle (0.28pt);
\fill[slate] (12.389,4.874) circle (0.28pt);
\fill[slate] (12.389,5.33) circle (0.28pt);
\fill[slate] (12.389,5.574) circle (0.28pt);
\fill[slate] (12.389,5.786) circle (0.28pt);
\fill[slate] (12.389,6.242) circle (0.28pt);
\fill[slate] (12.389,6.486) circle (0.28pt);
\fill[slate] (12.389,6.699) circle (0.28pt);
\fill[slate] (12.389,6.734) circle (0.28pt);
\fill[slate] (12.389,5.822) circle (0.28pt);
\fill[slate] (12.389,6.071) circle (0.28pt);
\fill[slate] (12.389,6.273) circle (0.28pt);
\fill[slate] (12.389,6.729) circle (0.28pt);
\fill[slate] (12.389,6.32) circle (0.28pt);
\fill[slate] (12.389,6.569) circle (0.28pt);
\fill[slate] (12.389,6.77) circle (0.28pt);
\fill[slate] (12.389,6.818) circle (0.28pt);
\fill[slate] (12.389,4.923) circle (0.28pt);
\fill[slate] (12.389,5.172) circle (0.28pt);
\fill[slate] (12.389,5.372) circle (0.28pt);
\fill[slate] (12.389,5.828) circle (0.28pt);
\fill[slate] (12.389,6.072) circle (0.28pt);
\fill[slate] (12.389,6.284) circle (0.28pt);
\fill[slate] (12.389,6.74) circle (0.28pt);
\fill[slate] (12.389,6.32) circle (0.28pt);
\fill[slate] (12.389,6.569) circle (0.28pt);
\fill[slate] (12.389,6.771) circle (0.28pt);
\fill[slate] (12.389,6.818) circle (0.28pt);
\fill[slate] (12.389,5.421) circle (0.28pt);
\fill[slate] (12.389,5.67) circle (0.28pt);
\fill[slate] (12.389,5.87) circle (0.28pt);
\fill[slate] (12.389,6.326) circle (0.28pt);
\fill[slate] (12.389,6.57) circle (0.28pt);
\fill[slate] (12.389,6.782) circle (0.28pt);
\fill[slate] (12.389,6.818) circle (0.28pt);
\fill[slate] (12.389,5.919) circle (0.28pt);
\fill[slate] (12.389,6.168) circle (0.28pt);
\fill[slate] (12.389,6.368) circle (0.28pt);
\fill[slate] (12.389,6.824) circle (0.28pt);
\fill[slate] (12.389,6.417) circle (0.28pt);
\fill[slate] (12.389,6.666) circle (0.28pt);
\fill[slate] (12.032,1.77) circle (0.28pt);
\fill[slate] (12.086,3.34) circle (0.28pt);
\fill[slate] (12.086,3.796) circle (0.28pt);
\fill[slate] (12.086,4.253) circle (0.28pt);
\fill[slate] (12.086,4.496) circle (0.28pt);
\fill[slate] (12.086,4.709) circle (0.28pt);
\fill[slate] (12.086,5.165) circle (0.28pt);
\fill[slate] (12.086,5.408) circle (0.28pt);
\fill[slate] (12.086,5.621) circle (0.28pt);
\fill[slate] (12.086,6.077) circle (0.28pt);
\fill[slate] (12.086,6.32) circle (0.28pt);
\fill[slate] (12.086,6.533) circle (0.28pt);
\fill[slate] (12.086,6.568) circle (0.28pt);
\fill[slate] (12.086,6.817) circle (0.28pt);
\fill[slate] (12.086,5.656) circle (0.28pt);
\fill[slate] (12.086,5.905) circle (0.28pt);
\fill[slate] (12.086,6.107) circle (0.28pt);
\fill[slate] (12.086,6.563) circle (0.28pt);
\fill[slate] (12.086,6.807) circle (0.28pt);
\fill[slate] (12.086,6.154) circle (0.28pt);
\fill[slate] (12.086,6.403) circle (0.28pt);
\fill[slate] (12.086,6.604) circle (0.28pt);
\fill[slate] (12.086,6.652) circle (0.28pt);
\fill[slate] (12.086,4.744) circle (0.28pt);
\fill[slate] (12.086,4.993) circle (0.28pt);
\fill[slate] (12.086,5.651) circle (0.28pt);
\fill[slate] (12.086,5.895) circle (0.28pt);
\fill[slate] (12.086,6.107) circle (0.28pt);
\fill[slate] (12.086,6.563) circle (0.28pt);
\fill[slate] (12.086,6.807) circle (0.28pt);
\fill[slate] (12.086,6.143) circle (0.28pt);
\fill[slate] (12.086,6.392) circle (0.28pt);
\fill[slate] (12.086,6.594) circle (0.28pt);
\fill[slate] (12.086,6.641) circle (0.28pt);
\fill[slate] (12.086,5.242) circle (0.28pt);
\fill[slate] (12.086,5.491) circle (0.28pt);
\fill[slate] (12.086,5.692) circle (0.28pt);
\fill[slate] (12.086,6.148) circle (0.28pt);
\fill[slate] (12.086,6.391) circle (0.28pt);
\fill[slate] (12.086,6.604) circle (0.28pt);
\fill[slate] (12.086,6.639) circle (0.28pt);
\fill[slate] (12.086,5.74) circle (0.28pt);
\fill[slate] (12.086,5.989) circle (0.28pt);
\fill[slate] (12.086,6.189) circle (0.28pt);
\fill[slate] (12.086,6.645) circle (0.28pt);
\fill[slate] (12.086,6.238) circle (0.28pt);
\fill[slate] (12.086,6.487) circle (0.28pt);
\fill[slate] (12.086,6.687) circle (0.28pt);
\fill[slate] (12.086,6.736) circle (0.28pt);
\fill[slate] (12.086,3.832) circle (0.28pt);
\fill[slate] (12.086,4.081) circle (0.28pt);
\fill[slate] (12.086,5.651) circle (0.28pt);
\fill[slate] (12.086,6.107) circle (0.28pt);
\fill[slate] (12.086,6.563) circle (0.28pt);
\fill[slate] (12.086,6.807) circle (0.28pt);
\fill[slate] (12.086,6.143) circle (0.28pt);
\fill[slate] (12.086,6.392) circle (0.28pt);
\fill[slate] (12.086,6.641) circle (0.28pt);
\fill[slate] (12.086,5.231) circle (0.28pt);
\fill[slate] (12.086,5.48) circle (0.28pt);
\fill[slate] (12.086,5.682) circle (0.28pt);
\fill[slate] (12.086,6.382) circle (0.28pt);
\fill[slate] (12.086,6.594) circle (0.28pt);
\fill[slate] (12.086,6.63) circle (0.28pt);
\fill[slate] (12.086,5.729) circle (0.28pt);
\fill[slate] (12.086,5.978) circle (0.28pt);
\fill[slate] (12.086,6.179) circle (0.28pt);
\fill[slate] (12.086,6.634) circle (0.28pt);
\fill[slate] (12.086,6.227) circle (0.28pt);
\fill[slate] (12.086,6.476) circle (0.28pt);
\fill[slate] (12.086,6.676) circle (0.28pt);
\fill[slate] (12.086,6.725) circle (0.28pt);
\fill[slate] (12.086,4.33) circle (0.28pt);
\fill[slate] (12.086,4.579) circle (0.28pt);
\fill[slate] (12.086,4.78) circle (0.28pt);
\fill[slate] (12.086,5.235) circle (0.28pt);
\fill[slate] (12.086,5.479) circle (0.28pt);
\fill[slate] (12.086,5.691) circle (0.28pt);
\fill[slate] (12.086,6.148) circle (0.28pt);
\fill[slate] (12.086,6.391) circle (0.28pt);
\fill[slate] (12.086,6.604) circle (0.28pt);
\fill[slate] (12.086,6.639) circle (0.28pt);
\fill[slate] (12.086,5.727) circle (0.28pt);
\fill[slate] (12.086,5.976) circle (0.28pt);
\fill[slate] (12.086,6.178) circle (0.28pt);
\fill[slate] (12.086,6.634) circle (0.28pt);
\fill[slate] (12.086,6.225) circle (0.28pt);
\fill[slate] (12.086,6.474) circle (0.28pt);
\fill[slate] (12.086,6.675) circle (0.28pt);
\fill[slate] (12.086,6.723) circle (0.28pt);
\fill[slate] (12.086,4.828) circle (0.28pt);
\fill[slate] (12.086,5.077) circle (0.28pt);
\fill[slate] (12.086,5.277) circle (0.28pt);
\fill[slate] (12.086,5.733) circle (0.28pt);
\fill[slate] (12.086,5.977) circle (0.28pt);
\fill[slate] (12.086,6.189) circle (0.28pt);
\fill[slate] (12.086,6.645) circle (0.28pt);
\fill[slate] (12.086,6.225) circle (0.28pt);
\fill[slate] (12.086,6.474) circle (0.28pt);
\fill[slate] (12.086,6.676) circle (0.28pt);
\fill[slate] (12.086,6.723) circle (0.28pt);
\fill[slate] (12.086,5.326) circle (0.28pt);
\fill[slate] (12.086,5.575) circle (0.28pt);
\fill[slate] (12.086,5.775) circle (0.28pt);
\fill[slate] (12.086,6.231) circle (0.28pt);
\fill[slate] (12.086,6.475) circle (0.28pt);
\fill[slate] (12.086,6.687) circle (0.28pt);
\fill[slate] (12.086,6.723) circle (0.28pt);
\fill[slate] (12.086,5.824) circle (0.28pt);
\fill[slate] (12.086,6.073) circle (0.28pt);
\fill[slate] (12.086,6.273) circle (0.28pt);
\fill[slate] (12.086,6.729) circle (0.28pt);
\fill[slate] (12.086,6.322) circle (0.28pt);
\fill[slate] (12.086,6.571) circle (0.28pt);
\fill[slate] (12.086,6.771) circle (0.28pt);
\fill[slate] (12.086,6.82) circle (0.28pt);
\fill[slate] (12.084,3.169) circle (0.28pt);
\fill[slate] (12.085,4.283) circle (0.28pt);
\fill[slate] (12.085,4.739) circle (0.28pt);
\fill[slate] (12.085,5.195) circle (0.28pt);
\fill[slate] (12.085,5.651) circle (0.28pt);
\fill[slate] (12.085,5.895) circle (0.28pt);
\fill[slate] (12.085,6.108) circle (0.28pt);
\fill[slate] (12.085,6.564) circle (0.28pt);
\fill[slate] (12.085,6.807) circle (0.28pt);
\fill[slate] (12.085,6.143) circle (0.28pt);
\fill[slate] (12.085,6.392) circle (0.28pt);
\fill[slate] (12.085,6.594) circle (0.28pt);
\fill[slate] (12.085,6.641) circle (0.28pt);
\fill[slate] (12.085,5.231) circle (0.28pt);
\fill[slate] (12.085,5.48) circle (0.28pt);
\fill[slate] (12.085,6.138) circle (0.28pt);
\fill[slate] (12.085,6.594) circle (0.28pt);
\fill[slate] (12.085,6.63) circle (0.28pt);
\fill[slate] (12.085,5.729) circle (0.28pt);
\fill[slate] (12.085,5.978) circle (0.28pt);
\fill[slate] (12.085,6.178) circle (0.28pt);
\fill[slate] (12.085,6.634) circle (0.28pt);
\fill[slate] (12.085,6.227) circle (0.28pt);
\fill[slate] (12.085,6.476) circle (0.28pt);
\fill[slate] (12.085,6.676) circle (0.28pt);
\fill[slate] (12.085,6.725) circle (0.28pt);
\fill[slate] (12.085,4.568) circle (0.28pt);
\fill[slate] (12.085,5.682) circle (0.28pt);
\fill[slate] (12.085,6.138) circle (0.28pt);
\fill[slate] (12.085,6.382) circle (0.28pt);
\fill[slate] (12.085,6.594) circle (0.28pt);
\fill[slate] (12.085,6.63) circle (0.28pt);
\fill[slate] (12.085,5.967) circle (0.28pt);
\fill[slate] (12.085,6.216) circle (0.28pt);
\fill[slate] (12.085,6.465) circle (0.28pt);
\fill[slate] (12.085,6.666) circle (0.28pt);
\fill[slate] (12.085,6.714) circle (0.28pt);
\fill[slate] (12.085,4.817) circle (0.28pt);
\fill[slate] (12.085,5.066) circle (0.28pt);
\fill[slate] (12.085,5.267) circle (0.28pt);
\fill[slate] (12.085,5.966) circle (0.28pt);
\fill[slate] (12.085,6.179) circle (0.28pt);
\fill[slate] (12.085,6.635) circle (0.28pt);
\fill[slate] (12.085,6.214) circle (0.28pt);
\fill[slate] (12.085,6.463) circle (0.28pt);
\fill[slate] (12.085,6.666) circle (0.28pt);
\fill[slate] (12.085,6.712) circle (0.28pt);
\fill[slate] (12.085,5.315) circle (0.28pt);
\fill[slate] (12.085,5.564) circle (0.28pt);
\fill[slate] (12.085,5.764) circle (0.28pt);
\fill[slate] (12.085,6.22) circle (0.28pt);
\fill[slate] (12.085,6.464) circle (0.28pt);
\fill[slate] (12.085,6.676) circle (0.28pt);
\fill[slate] (12.085,6.712) circle (0.28pt);
\fill[slate] (12.085,5.813) circle (0.28pt);
\fill[slate] (12.085,6.062) circle (0.28pt);
\fill[slate] (12.085,6.262) circle (0.28pt);
\fill[slate] (12.085,6.718) circle (0.28pt);
\fill[slate] (12.085,6.311) circle (0.28pt);
\fill[slate] (12.085,6.56) circle (0.28pt);
\fill[slate] (12.085,6.76) circle (0.28pt);
\fill[slate] (12.085,6.809) circle (0.28pt);
\fill[slate] (12.084,3.418) circle (0.28pt);
\fill[slate] (12.084,3.667) circle (0.28pt);
\fill[slate] (12.084,3.868) circle (0.28pt);
\fill[slate] (12.084,4.567) circle (0.28pt);
\fill[slate] (12.084,4.78) circle (0.28pt);
\fill[slate] (12.084,5.236) circle (0.28pt);
\fill[slate] (12.084,5.479) circle (0.28pt);
\fill[slate] (12.084,5.692) circle (0.28pt);
\fill[slate] (12.084,6.148) circle (0.28pt);
\fill[slate] (12.084,6.391) circle (0.28pt);
\fill[slate] (12.084,6.604) circle (0.28pt);
\fill[slate] (12.084,6.639) circle (0.28pt);
\fill[slate] (12.084,5.976) circle (0.28pt);
\fill[slate] (12.084,6.178) circle (0.28pt);
\fill[slate] (12.084,6.634) circle (0.28pt);
\fill[slate] (12.084,6.225) circle (0.28pt);
\fill[slate] (12.084,6.474) circle (0.28pt);
\fill[slate] (12.084,6.675) circle (0.28pt);
\fill[slate] (12.084,6.723) circle (0.28pt);
\fill[slate] (12.084,4.815) circle (0.28pt);
\fill[slate] (12.084,5.064) circle (0.28pt);
\fill[slate] (12.084,5.267) circle (0.28pt);
\fill[slate] (12.084,5.966) circle (0.28pt);
\fill[slate] (12.084,6.179) circle (0.28pt);
\fill[slate] (12.084,6.635) circle (0.28pt);
\fill[slate] (12.084,6.214) circle (0.28pt);
\fill[slate] (12.084,6.463) circle (0.28pt);
\fill[slate] (12.084,6.666) circle (0.28pt);
\fill[slate] (12.084,6.712) circle (0.28pt);
\fill[slate] (12.084,5.313) circle (0.28pt);
\fill[slate] (12.084,5.562) circle (0.28pt);
\fill[slate] (12.084,5.763) circle (0.28pt);
\fill[slate] (12.084,6.219) circle (0.28pt);
\fill[slate] (12.084,6.463) circle (0.28pt);
\fill[slate] (12.084,6.675) circle (0.28pt);
\fill[slate] (12.084,6.711) circle (0.28pt);
\fill[slate] (12.084,5.811) circle (0.28pt);
\fill[slate] (12.084,6.06) circle (0.28pt);
\fill[slate] (12.084,6.261) circle (0.28pt);
\fill[slate] (12.084,6.717) circle (0.28pt);
\fill[slate] (12.084,6.309) circle (0.28pt);
\fill[slate] (12.084,6.558) circle (0.28pt);
\fill[slate] (12.084,6.759) circle (0.28pt);
\fill[slate] (12.084,6.807) circle (0.28pt);
\fill[slate] (12.084,3.916) circle (0.28pt);
\fill[slate] (12.084,4.165) circle (0.28pt);
\fill[slate] (12.084,4.365) circle (0.28pt);
\fill[slate] (12.084,4.821) circle (0.28pt);
\fill[slate] (12.084,5.065) circle (0.28pt);
\fill[slate] (12.084,5.277) circle (0.28pt);
\fill[slate] (12.084,5.733) circle (0.28pt);
\fill[slate] (12.084,5.977) circle (0.28pt);
\fill[slate] (12.084,6.19) circle (0.28pt);
\fill[slate] (12.084,6.646) circle (0.28pt);
\fill[slate] (12.084,6.225) circle (0.28pt);
\fill[slate] (12.084,6.474) circle (0.28pt);
\fill[slate] (12.084,6.676) circle (0.28pt);
\fill[slate] (12.084,6.723) circle (0.28pt);
\fill[slate] (12.084,5.313) circle (0.28pt);
\fill[slate] (12.084,5.562) circle (0.28pt);
\fill[slate] (12.084,5.764) circle (0.28pt);
\fill[slate] (12.084,6.22) circle (0.28pt);
\fill[slate] (12.084,6.464) circle (0.28pt);
\fill[slate] (12.084,6.676) circle (0.28pt);
\fill[slate] (12.084,6.712) circle (0.28pt);
\fill[slate] (12.084,5.811) circle (0.28pt);
\fill[slate] (12.084,6.06) circle (0.28pt);
\fill[slate] (12.084,6.261) circle (0.28pt);
\fill[slate] (12.084,6.717) circle (0.28pt);
\fill[slate] (12.084,6.309) circle (0.28pt);
\fill[slate] (12.084,6.558) circle (0.28pt);
\fill[slate] (12.084,6.759) circle (0.28pt);
\fill[slate] (12.084,6.807) circle (0.28pt);
\fill[slate] (12.084,4.414) circle (0.28pt);
\fill[slate] (12.084,4.663) circle (0.28pt);
\fill[slate] (12.084,4.863) circle (0.28pt);
\fill[slate] (12.084,5.319) circle (0.28pt);
\fill[slate] (12.084,5.563) circle (0.28pt);
\fill[slate] (12.084,5.775) circle (0.28pt);
\fill[slate] (12.084,6.231) circle (0.28pt);
\fill[slate] (12.084,6.475) circle (0.28pt);
\fill[slate] (12.084,6.688) circle (0.28pt);
\fill[slate] (12.084,6.723) circle (0.28pt);
\fill[slate] (12.084,5.811) circle (0.28pt);
\fill[slate] (12.084,6.06) circle (0.28pt);
\fill[slate] (12.084,6.262) circle (0.28pt);
\fill[slate] (12.084,6.718) circle (0.28pt);
\fill[slate] (12.084,6.309) circle (0.28pt);
\fill[slate] (12.084,6.558) circle (0.28pt);
\fill[slate] (12.084,6.759) circle (0.28pt);
\fill[slate] (12.084,6.807) circle (0.28pt);
\fill[slate] (12.084,4.912) circle (0.28pt);
\fill[slate] (12.084,5.161) circle (0.28pt);
\fill[slate] (12.084,5.361) circle (0.28pt);
\fill[slate] (12.084,5.817) circle (0.28pt);
\fill[slate] (12.084,6.061) circle (0.28pt);
\fill[slate] (12.084,6.273) circle (0.28pt);
\fill[slate] (12.084,6.729) circle (0.28pt);
\fill[slate] (12.084,6.309) circle (0.28pt);
\fill[slate] (12.084,6.558) circle (0.28pt);
\fill[slate] (12.084,6.76) circle (0.28pt);
\fill[slate] (12.084,6.807) circle (0.28pt);
\fill[slate] (12.084,5.41) circle (0.28pt);
\fill[slate] (12.084,5.659) circle (0.28pt);
\fill[slate] (12.084,5.859) circle (0.28pt);
\fill[slate] (12.084,6.315) circle (0.28pt);
\fill[slate] (12.084,6.559) circle (0.28pt);
\fill[slate] (12.084,6.771) circle (0.28pt);
\fill[slate] (12.084,6.807) circle (0.28pt);
\fill[slate] (12.084,5.908) circle (0.28pt);
\fill[slate] (12.084,6.157) circle (0.28pt);
\fill[slate] (12.084,6.357) circle (0.28pt);
\fill[slate] (12.084,6.813) circle (0.28pt);
\fill[slate] (12.084,6.406) circle (0.28pt);
\fill[slate] (12.084,6.655) circle (0.28pt);
\fill[slate] (12.031,2.019) circle (0.28pt);
\fill[slate] (12.031,2.268) circle (0.28pt);
\fill[slate] (12.039,2.469) circle (0.28pt);
\fill[slate] (12.039,3.168) circle (0.28pt);
\fill[slate] (12.039,3.381) circle (0.28pt);
\fill[slate] (12.039,3.837) circle (0.28pt);
\fill[slate] (12.039,4.08) circle (0.28pt);
\fill[slate] (12.039,4.293) circle (0.28pt);
\fill[slate] (12.039,4.749) circle (0.28pt);
\fill[slate] (12.039,4.992) circle (0.28pt);
\fill[slate] (12.039,5.205) circle (0.28pt);
\fill[slate] (12.039,5.661) circle (0.28pt);
\fill[slate] (12.039,5.904) circle (0.28pt);
\fill[slate] (12.039,6.117) circle (0.28pt);
\fill[slate] (12.039,6.573) circle (0.28pt);
\fill[slate] (12.039,6.816) circle (0.28pt);
\fill[slate] (12.039,6.152) circle (0.28pt);
\fill[slate] (12.039,6.401) circle (0.28pt);
\fill[slate] (12.039,6.604) circle (0.28pt);
\fill[slate] (12.039,6.65) circle (0.28pt);
\fill[slate] (12.039,5.24) circle (0.28pt);
\fill[slate] (12.039,5.489) circle (0.28pt);
\fill[slate] (12.039,5.692) circle (0.28pt);
\fill[slate] (12.039,6.148) circle (0.28pt);
\fill[slate] (12.039,6.391) circle (0.28pt);
\fill[slate] (12.039,6.604) circle (0.28pt);
\fill[slate] (12.039,6.639) circle (0.28pt);
\fill[slate] (12.039,5.738) circle (0.28pt);
\fill[slate] (12.039,5.987) circle (0.28pt);
\fill[slate] (12.039,6.188) circle (0.28pt);
\fill[slate] (12.039,6.644) circle (0.28pt);
\fill[slate] (12.039,6.236) circle (0.28pt);
\fill[slate] (12.039,6.485) circle (0.28pt);
\fill[slate] (12.039,6.686) circle (0.28pt);
\fill[slate] (12.039,6.734) circle (0.28pt);
\fill[slate] (12.039,4.577) circle (0.28pt);
\fill[slate] (12.039,4.779) circle (0.28pt);
\fill[slate] (12.039,5.235) circle (0.28pt);
\fill[slate] (12.039,5.479) circle (0.28pt);
\fill[slate] (12.039,5.691) circle (0.28pt);
\fill[slate] (12.039,6.148) circle (0.28pt);
\fill[slate] (12.039,6.391) circle (0.28pt);
\fill[slate] (12.039,6.604) circle (0.28pt);
\fill[slate] (12.039,6.639) circle (0.28pt);
\fill[slate] (12.039,5.976) circle (0.28pt);
\fill[slate] (12.039,6.178) circle (0.28pt);
\fill[slate] (12.039,6.634) circle (0.28pt);
\fill[slate] (12.039,6.225) circle (0.28pt);
\fill[slate] (12.039,6.474) circle (0.28pt);
\fill[slate] (12.039,6.675) circle (0.28pt);
\fill[slate] (12.039,6.723) circle (0.28pt);
\fill[slate] (12.039,4.826) circle (0.28pt);
\fill[slate] (12.039,5.075) circle (0.28pt);
\fill[slate] (12.039,5.276) circle (0.28pt);
\fill[slate] (12.039,5.975) circle (0.28pt);
\fill[slate] (12.039,6.188) circle (0.28pt);
\fill[slate] (12.039,6.644) circle (0.28pt);
\fill[slate] (12.039,6.224) circle (0.28pt);
\fill[slate] (12.039,6.472) circle (0.28pt);
\fill[slate] (12.039,6.675) circle (0.28pt);
\fill[slate] (12.039,6.721) circle (0.28pt);
\fill[slate] (12.039,5.324) circle (0.28pt);
\fill[slate] (12.039,5.573) circle (0.28pt);
\fill[slate] (12.039,5.774) circle (0.28pt);
\fill[slate] (12.039,6.23) circle (0.28pt);
\fill[slate] (12.039,6.473) circle (0.28pt);
\fill[slate] (12.039,6.686) circle (0.28pt);
\fill[slate] (12.039,6.721) circle (0.28pt);
\fill[slate] (12.039,5.822) circle (0.28pt);
\fill[slate] (12.039,6.071) circle (0.28pt);
\fill[slate] (12.039,6.272) circle (0.28pt);
\fill[slate] (12.039,6.727) circle (0.28pt);
\fill[slate] (12.039,6.32) circle (0.28pt);
\fill[slate] (12.039,6.569) circle (0.28pt);
\fill[slate] (12.039,6.77) circle (0.28pt);
\fill[slate] (12.039,6.818) circle (0.28pt);
\fill[slate] (12.039,3.416) circle (0.28pt);
\fill[slate] (12.039,3.665) circle (0.28pt);
\fill[slate] (12.039,3.868) circle (0.28pt);
\fill[slate] (12.039,4.567) circle (0.28pt);
\fill[slate] (12.039,4.78) circle (0.28pt);
\fill[slate] (12.039,5.236) circle (0.28pt);
\fill[slate] (12.039,5.479) circle (0.28pt);
\fill[slate] (12.039,5.692) circle (0.28pt);
\fill[slate] (12.039,6.148) circle (0.28pt);
\fill[slate] (12.039,6.391) circle (0.28pt);
\fill[slate] (12.039,6.604) circle (0.28pt);
\fill[slate] (12.039,6.639) circle (0.28pt);
\fill[slate] (12.039,5.976) circle (0.28pt);
\fill[slate] (12.039,6.179) circle (0.28pt);
\fill[slate] (12.039,6.634) circle (0.28pt);
\fill[slate] (12.039,6.225) circle (0.28pt);
\fill[slate] (12.039,6.474) circle (0.28pt);
\fill[slate] (12.039,6.675) circle (0.28pt);
\fill[slate] (12.039,6.723) circle (0.28pt);
\fill[slate] (12.039,4.815) circle (0.28pt);
\fill[slate] (12.039,5.064) circle (0.28pt);
\fill[slate] (12.039,5.267) circle (0.28pt);
\fill[slate] (12.039,5.966) circle (0.28pt);
\fill[slate] (12.039,6.179) circle (0.28pt);
\fill[slate] (12.039,6.635) circle (0.28pt);
\fill[slate] (12.039,6.214) circle (0.28pt);
\fill[slate] (12.039,6.463) circle (0.28pt);
\fill[slate] (12.039,6.666) circle (0.28pt);
\fill[slate] (12.039,6.712) circle (0.28pt);
\fill[slate] (12.039,5.313) circle (0.28pt);
\fill[slate] (12.039,5.562) circle (0.28pt);
\fill[slate] (12.039,5.763) circle (0.28pt);
\fill[slate] (12.039,6.219) circle (0.28pt);
\fill[slate] (12.039,6.463) circle (0.28pt);
\fill[slate] (12.039,6.675) circle (0.28pt);
\fill[slate] (12.039,6.711) circle (0.28pt);
\fill[slate] (12.039,5.811) circle (0.28pt);
\fill[slate] (12.039,6.06) circle (0.28pt);
\fill[slate] (12.039,6.261) circle (0.28pt);
\fill[slate] (12.039,6.717) circle (0.28pt);
\fill[slate] (12.039,6.309) circle (0.28pt);
\fill[slate] (12.039,6.558) circle (0.28pt);
\fill[slate] (12.039,6.759) circle (0.28pt);
\fill[slate] (12.039,6.807) circle (0.28pt);
\fill[slate] (12.039,3.914) circle (0.28pt);
\fill[slate] (12.039,4.163) circle (0.28pt);
\fill[slate] (12.039,4.364) circle (0.28pt);
\fill[slate] (12.039,4.82) circle (0.28pt);
\fill[slate] (12.039,5.064) circle (0.28pt);
\fill[slate] (12.039,5.276) circle (0.28pt);
\fill[slate] (12.039,5.732) circle (0.28pt);
\fill[slate] (12.039,5.975) circle (0.28pt);
\fill[slate] (12.039,6.188) circle (0.28pt);
\fill[slate] (12.039,6.644) circle (0.28pt);
\fill[slate] (12.039,6.224) circle (0.28pt);
\fill[slate] (12.039,6.472) circle (0.28pt);
\fill[slate] (12.039,6.675) circle (0.28pt);
\fill[slate] (12.039,6.721) circle (0.28pt);
\fill[slate] (12.039,5.312) circle (0.28pt);
\fill[slate] (12.039,5.561) circle (0.28pt);
\fill[slate] (12.039,5.763) circle (0.28pt);
\fill[slate] (12.039,6.219) circle (0.28pt);
\fill[slate] (12.039,6.463) circle (0.28pt);
\fill[slate] (12.039,6.675) circle (0.28pt);
\fill[slate] (12.039,6.711) circle (0.28pt);
\fill[slate] (12.039,5.81) circle (0.28pt);
\fill[slate] (12.039,6.059) circle (0.28pt);
\fill[slate] (12.039,6.259) circle (0.28pt);
\fill[slate] (12.039,6.715) circle (0.28pt);
\fill[slate] (12.039,6.308) circle (0.28pt);
\fill[slate] (12.039,6.557) circle (0.28pt);
\fill[slate] (12.039,6.757) circle (0.28pt);
\fill[slate] (12.039,6.806) circle (0.28pt);
\fill[slate] (12.039,4.412) circle (0.28pt);
\fill[slate] (12.039,4.661) circle (0.28pt);
\fill[slate] (12.039,4.862) circle (0.28pt);
\fill[slate] (12.039,5.318) circle (0.28pt);
\fill[slate] (12.039,5.561) circle (0.28pt);
\fill[slate] (12.039,5.774) circle (0.28pt);
\fill[slate] (12.039,6.23) circle (0.28pt);
\fill[slate] (12.039,6.473) circle (0.28pt);
\fill[slate] (12.039,6.686) circle (0.28pt);
\fill[slate] (12.039,6.721) circle (0.28pt);
\fill[slate] (12.039,5.81) circle (0.28pt);
\fill[slate] (12.039,6.058) circle (0.28pt);
\fill[slate] (12.039,6.261) circle (0.28pt);
\fill[slate] (12.039,6.717) circle (0.28pt);
\fill[slate] (12.039,6.307) circle (0.28pt);
\fill[slate] (12.039,6.556) circle (0.28pt);
\fill[slate] (12.039,6.757) circle (0.28pt);
\fill[slate] (12.039,6.805) circle (0.28pt);
\fill[slate] (12.039,4.91) circle (0.28pt);
\fill[slate] (12.039,5.159) circle (0.28pt);
\fill[slate] (12.039,5.36) circle (0.28pt);
\fill[slate] (12.039,5.816) circle (0.28pt);
\fill[slate] (12.039,6.059) circle (0.28pt);
\fill[slate] (12.039,6.272) circle (0.28pt);
\fill[slate] (12.039,6.728) circle (0.28pt);
\fill[slate] (12.039,6.308) circle (0.28pt);
\fill[slate] (12.039,6.556) circle (0.28pt);
\fill[slate] (12.039,6.759) circle (0.28pt);
\fill[slate] (12.039,6.805) circle (0.28pt);
\fill[slate] (12.039,5.408) circle (0.28pt);
\fill[slate] (12.039,5.657) circle (0.28pt);
\fill[slate] (12.039,5.858) circle (0.28pt);
\fill[slate] (12.039,6.314) circle (0.28pt);
\fill[slate] (12.039,6.557) circle (0.28pt);
\fill[slate] (12.039,6.77) circle (0.28pt);
\fill[slate] (12.039,6.806) circle (0.28pt);
\fill[slate] (12.039,5.906) circle (0.28pt);
\fill[slate] (12.039,6.155) circle (0.28pt);
\fill[slate] (12.039,6.356) circle (0.28pt);
\fill[slate] (12.039,6.812) circle (0.28pt);
\fill[slate] (12.039,6.404) circle (0.28pt);
\fill[slate] (12.039,6.653) circle (0.28pt);
\fill[slate] (12.031,2.517) circle (0.28pt);
\fill[slate] (12.031,2.766) circle (0.28pt);
\fill[slate] (12.032,2.966) circle (0.28pt);
\fill[slate] (12.032,3.422) circle (0.28pt);
\fill[slate] (12.032,3.666) circle (0.28pt);
\fill[slate] (12.032,3.878) circle (0.28pt);
\fill[slate] (12.032,4.334) circle (0.28pt);
\fill[slate] (12.032,4.578) circle (0.28pt);
\fill[slate] (12.032,4.791) circle (0.28pt);
\fill[slate] (12.032,5.247) circle (0.28pt);
\fill[slate] (12.032,5.49) circle (0.28pt);
\fill[slate] (12.032,5.703) circle (0.28pt);
\fill[slate] (12.032,6.159) circle (0.28pt);
\fill[slate] (12.032,6.402) circle (0.28pt);
\fill[slate] (12.032,6.615) circle (0.28pt);
\fill[slate] (12.032,6.65) circle (0.28pt);
\fill[slate] (12.032,5.738) circle (0.28pt);
\fill[slate] (12.032,5.987) circle (0.28pt);
\fill[slate] (12.032,6.189) circle (0.28pt);
\fill[slate] (12.032,6.645) circle (0.28pt);
\fill[slate] (12.032,6.236) circle (0.28pt);
\fill[slate] (12.032,6.485) circle (0.28pt);
\fill[slate] (12.032,6.686) circle (0.28pt);
\fill[slate] (12.032,6.734) circle (0.28pt);
\fill[slate] (12.032,4.826) circle (0.28pt);
\fill[slate] (12.032,5.075) circle (0.28pt);
\fill[slate] (12.032,5.277) circle (0.28pt);
\fill[slate] (12.032,5.733) circle (0.28pt);
\fill[slate] (12.032,5.977) circle (0.28pt);
\fill[slate] (12.032,6.189) circle (0.28pt);
\fill[slate] (12.032,6.645) circle (0.28pt);
\fill[slate] (12.032,6.225) circle (0.28pt);
\fill[slate] (12.032,6.676) circle (0.28pt);
\fill[slate] (12.032,6.723) circle (0.28pt);
\fill[slate] (12.032,5.324) circle (0.28pt);
\fill[slate] (12.032,5.573) circle (0.28pt);
\fill[slate] (12.032,5.774) circle (0.28pt);
\fill[slate] (12.032,6.23) circle (0.28pt);
\fill[slate] (12.032,6.473) circle (0.28pt);
\fill[slate] (12.032,6.686) circle (0.28pt);
\fill[slate] (12.032,6.721) circle (0.28pt);
\fill[slate] (12.032,5.822) circle (0.28pt);
\fill[slate] (12.032,6.071) circle (0.28pt);
\fill[slate] (12.032,6.271) circle (0.28pt);
\fill[slate] (12.032,6.727) circle (0.28pt);
\fill[slate] (12.032,6.32) circle (0.28pt);
\fill[slate] (12.032,6.569) circle (0.28pt);
\fill[slate] (12.032,6.769) circle (0.28pt);
\fill[slate] (12.032,6.818) circle (0.28pt);
\fill[slate] (12.032,3.914) circle (0.28pt);
\fill[slate] (12.032,4.163) circle (0.28pt);
\fill[slate] (12.032,4.365) circle (0.28pt);
\fill[slate] (12.032,4.821) circle (0.28pt);
\fill[slate] (12.032,5.065) circle (0.28pt);
\fill[slate] (12.032,5.277) circle (0.28pt);
\fill[slate] (12.032,5.734) circle (0.28pt);
\fill[slate] (12.032,5.977) circle (0.28pt);
\fill[slate] (12.032,6.19) circle (0.28pt);
\fill[slate] (12.032,6.646) circle (0.28pt);
\fill[slate] (12.032,6.225) circle (0.28pt);
\fill[slate] (12.032,6.676) circle (0.28pt);
\fill[slate] (12.032,6.723) circle (0.28pt);
\fill[slate] (12.032,5.313) circle (0.28pt);
\fill[slate] (12.032,5.562) circle (0.28pt);
\fill[slate] (12.032,5.764) circle (0.28pt);
\fill[slate] (12.032,6.22) circle (0.28pt);
\fill[slate] (12.032,6.464) circle (0.28pt);
\fill[slate] (12.032,6.676) circle (0.28pt);
\fill[slate] (12.032,6.712) circle (0.28pt);
\fill[slate] (12.032,5.811) circle (0.28pt);
\fill[slate] (12.032,6.06) circle (0.28pt);
\fill[slate] (12.032,6.261) circle (0.28pt);
\fill[slate] (12.032,6.717) circle (0.28pt);
\fill[slate] (12.032,6.309) circle (0.28pt);
\fill[slate] (12.032,6.558) circle (0.28pt);
\fill[slate] (12.032,6.759) circle (0.28pt);
\fill[slate] (12.032,6.807) circle (0.28pt);
\fill[slate] (12.032,4.412) circle (0.28pt);
\fill[slate] (12.032,4.661) circle (0.28pt);
\fill[slate] (12.032,4.862) circle (0.28pt);
\fill[slate] (12.032,5.318) circle (0.28pt);
\fill[slate] (12.032,5.561) circle (0.28pt);
\fill[slate] (12.032,5.774) circle (0.28pt);
\fill[slate] (12.032,6.23) circle (0.28pt);
\fill[slate] (12.032,6.473) circle (0.28pt);
\fill[slate] (12.032,6.686) circle (0.28pt);
\fill[slate] (12.032,6.721) circle (0.28pt);
\fill[slate] (12.032,5.81) circle (0.28pt);
\fill[slate] (12.032,6.058) circle (0.28pt);
\fill[slate] (12.032,6.261) circle (0.28pt);
\fill[slate] (12.032,6.717) circle (0.28pt);
\fill[slate] (12.032,6.307) circle (0.28pt);
\fill[slate] (12.032,6.556) circle (0.28pt);
\fill[slate] (12.032,6.757) circle (0.28pt);
\fill[slate] (12.032,6.805) circle (0.28pt);
\fill[slate] (12.032,4.91) circle (0.28pt);
\fill[slate] (12.032,5.159) circle (0.28pt);
\fill[slate] (12.032,5.36) circle (0.28pt);
\fill[slate] (12.032,5.816) circle (0.28pt);
\fill[slate] (12.032,6.059) circle (0.28pt);
\fill[slate] (12.032,6.272) circle (0.28pt);
\fill[slate] (12.032,6.728) circle (0.28pt);
\fill[slate] (12.032,6.307) circle (0.28pt);
\fill[slate] (12.032,6.556) circle (0.28pt);
\fill[slate] (12.032,6.759) circle (0.28pt);
\fill[slate] (12.032,6.805) circle (0.28pt);
\fill[slate] (12.032,5.408) circle (0.28pt);
\fill[slate] (12.032,5.657) circle (0.28pt);
\fill[slate] (12.032,5.858) circle (0.28pt);
\fill[slate] (12.032,6.313) circle (0.28pt);
\fill[slate] (12.032,6.557) circle (0.28pt);
\fill[slate] (12.032,6.77) circle (0.28pt);
\fill[slate] (12.032,6.805) circle (0.28pt);
\fill[slate] (12.032,5.906) circle (0.28pt);
\fill[slate] (12.032,6.155) circle (0.28pt);
\fill[slate] (12.032,6.356) circle (0.28pt);
\fill[slate] (12.032,6.811) circle (0.28pt);
\fill[slate] (12.032,6.404) circle (0.28pt);
\fill[slate] (12.032,6.653) circle (0.28pt);
\fill[slate] (12.031,3.015) circle (0.28pt);
\fill[slate] (12.031,3.264) circle (0.28pt);
\fill[slate] (12.031,3.464) circle (0.28pt);
\fill[slate] (12.031,3.92) circle (0.28pt);
\fill[slate] (12.031,4.164) circle (0.28pt);
\fill[slate] (12.031,4.376) circle (0.28pt);
\fill[slate] (12.031,4.832) circle (0.28pt);
\fill[slate] (12.031,5.076) circle (0.28pt);
\fill[slate] (12.031,5.289) circle (0.28pt);
\fill[slate] (12.031,5.745) circle (0.28pt);
\fill[slate] (12.031,5.988) circle (0.28pt);
\fill[slate] (12.031,6.201) circle (0.28pt);
\fill[slate] (12.031,6.657) circle (0.28pt);
\fill[slate] (12.031,6.236) circle (0.28pt);
\fill[slate] (12.031,6.485) circle (0.28pt);
\fill[slate] (12.031,6.687) circle (0.28pt);
\fill[slate] (12.031,6.734) circle (0.28pt);
\fill[slate] (12.031,5.324) circle (0.28pt);
\fill[slate] (12.031,5.573) circle (0.28pt);
\fill[slate] (12.031,5.775) circle (0.28pt);
\fill[slate] (12.031,6.231) circle (0.28pt);
\fill[slate] (12.031,6.475) circle (0.28pt);
\fill[slate] (12.031,6.687) circle (0.28pt);
\fill[slate] (12.031,6.723) circle (0.28pt);
\fill[slate] (12.031,5.822) circle (0.28pt);
\fill[slate] (12.031,6.071) circle (0.28pt);
\fill[slate] (12.031,6.272) circle (0.28pt);
\fill[slate] (12.031,6.727) circle (0.28pt);
\fill[slate] (12.031,6.32) circle (0.28pt);
\fill[slate] (12.031,6.569) circle (0.28pt);
\fill[slate] (12.031,6.769) circle (0.28pt);
\fill[slate] (12.031,6.818) circle (0.28pt);
\fill[slate] (12.031,4.412) circle (0.28pt);
\fill[slate] (12.031,4.661) circle (0.28pt);
\fill[slate] (12.031,4.863) circle (0.28pt);
\fill[slate] (12.031,5.319) circle (0.28pt);
\fill[slate] (12.031,5.563) circle (0.28pt);
\fill[slate] (12.031,5.775) circle (0.28pt);
\fill[slate] (12.031,6.231) circle (0.28pt);
\fill[slate] (12.031,6.475) circle (0.28pt);
\fill[slate] (12.031,6.688) circle (0.28pt);
\fill[slate] (12.031,6.723) circle (0.28pt);
\fill[slate] (12.031,5.811) circle (0.28pt);
\fill[slate] (12.031,6.06) circle (0.28pt);
\fill[slate] (12.031,6.262) circle (0.28pt);
\fill[slate] (12.031,6.718) circle (0.28pt);
\fill[slate] (12.031,6.309) circle (0.28pt);
\fill[slate] (12.031,6.558) circle (0.28pt);
\fill[slate] (12.031,6.759) circle (0.28pt);
\fill[slate] (12.031,6.807) circle (0.28pt);
\fill[slate] (12.031,4.91) circle (0.28pt);
\fill[slate] (12.031,5.159) circle (0.28pt);
\fill[slate] (12.031,5.36) circle (0.28pt);
\fill[slate] (12.031,5.816) circle (0.28pt);
\fill[slate] (12.031,6.059) circle (0.28pt);
\fill[slate] (12.031,6.272) circle (0.28pt);
\fill[slate] (12.031,6.728) circle (0.28pt);
\fill[slate] (12.031,6.308) circle (0.28pt);
\fill[slate] (12.031,6.556) circle (0.28pt);
\fill[slate] (12.031,6.759) circle (0.28pt);
\fill[slate] (12.031,6.805) circle (0.28pt);
\fill[slate] (12.031,5.408) circle (0.28pt);
\fill[slate] (12.031,5.657) circle (0.28pt);
\fill[slate] (12.031,5.858) circle (0.28pt);
\fill[slate] (12.031,6.313) circle (0.28pt);
\fill[slate] (12.031,6.557) circle (0.28pt);
\fill[slate] (12.031,6.77) circle (0.28pt);
\fill[slate] (12.031,6.805) circle (0.28pt);
\fill[slate] (12.031,5.906) circle (0.28pt);
\fill[slate] (12.031,6.155) circle (0.28pt);
\fill[slate] (12.031,6.356) circle (0.28pt);
\fill[slate] (12.031,6.811) circle (0.28pt);
\fill[slate] (12.031,6.404) circle (0.28pt);
\fill[slate] (12.031,6.653) circle (0.28pt);
\fill[slate] (12.031,3.513) circle (0.28pt);
\fill[slate] (12.031,3.762) circle (0.28pt);
\fill[slate] (12.031,3.962) circle (0.28pt);
\fill[slate] (12.031,4.418) circle (0.28pt);
\fill[slate] (12.031,4.662) circle (0.28pt);
\fill[slate] (12.031,4.874) circle (0.28pt);
\fill[slate] (12.031,5.33) circle (0.28pt);
\fill[slate] (12.031,5.574) circle (0.28pt);
\fill[slate] (12.031,5.787) circle (0.28pt);
\fill[slate] (12.031,6.243) circle (0.28pt);
\fill[slate] (12.031,6.486) circle (0.28pt);
\fill[slate] (12.031,6.699) circle (0.28pt);
\fill[slate] (12.031,6.734) circle (0.28pt);
\fill[slate] (12.031,5.822) circle (0.28pt);
\fill[slate] (12.031,6.071) circle (0.28pt);
\fill[slate] (12.031,6.273) circle (0.28pt);
\fill[slate] (12.031,6.729) circle (0.28pt);
\fill[slate] (12.031,6.32) circle (0.28pt);
\fill[slate] (12.031,6.569) circle (0.28pt);
\fill[slate] (12.031,6.77) circle (0.28pt);
\fill[slate] (12.031,6.818) circle (0.28pt);
\fill[slate] (12.031,4.91) circle (0.28pt);
\fill[slate] (12.031,5.159) circle (0.28pt);
\fill[slate] (12.031,5.361) circle (0.28pt);
\fill[slate] (12.031,5.817) circle (0.28pt);
\fill[slate] (12.031,6.061) circle (0.28pt);
\fill[slate] (12.031,6.273) circle (0.28pt);
\fill[slate] (12.031,6.729) circle (0.28pt);
\fill[slate] (12.031,6.309) circle (0.28pt);
\fill[slate] (12.031,6.558) circle (0.28pt);
\fill[slate] (12.031,6.76) circle (0.28pt);
\fill[slate] (12.031,6.807) circle (0.28pt);
\fill[slate] (12.031,5.408) circle (0.28pt);
\fill[slate] (12.031,5.657) circle (0.28pt);
\fill[slate] (12.031,5.858) circle (0.28pt);
\fill[slate] (12.031,6.314) circle (0.28pt);
\fill[slate] (12.031,6.557) circle (0.28pt);
\fill[slate] (12.031,6.77) circle (0.28pt);
\fill[slate] (12.031,6.806) circle (0.28pt);
\fill[slate] (12.031,5.906) circle (0.28pt);
\fill[slate] (12.031,6.155) circle (0.28pt);
\fill[slate] (12.031,6.356) circle (0.28pt);
\fill[slate] (12.031,6.811) circle (0.28pt);
\fill[slate] (12.031,6.404) circle (0.28pt);
\fill[slate] (12.031,6.653) circle (0.28pt);
\fill[slate] (12.031,4.011) circle (0.28pt);
\fill[slate] (12.031,4.26) circle (0.28pt);
\fill[slate] (12.031,4.46) circle (0.28pt);
\fill[slate] (12.031,4.916) circle (0.28pt);
\fill[slate] (12.031,5.16) circle (0.28pt);
\fill[slate] (12.031,5.372) circle (0.28pt);
\fill[slate] (12.031,5.828) circle (0.28pt);
\fill[slate] (12.031,6.072) circle (0.28pt);
\fill[slate] (12.031,6.285) circle (0.28pt);
\fill[slate] (12.031,6.741) circle (0.28pt);
\fill[slate] (12.031,6.32) circle (0.28pt);
\fill[slate] (12.031,6.569) circle (0.28pt);
\fill[slate] (12.031,6.771) circle (0.28pt);
\fill[slate] (12.031,6.818) circle (0.28pt);
\fill[slate] (12.031,5.408) circle (0.28pt);
\fill[slate] (12.031,5.657) circle (0.28pt);
\fill[slate] (12.031,5.859) circle (0.28pt);
\fill[slate] (12.031,6.315) circle (0.28pt);
\fill[slate] (12.031,6.559) circle (0.28pt);
\fill[slate] (12.031,6.771) circle (0.28pt);
\fill[slate] (12.031,6.807) circle (0.28pt);
\fill[slate] (12.031,5.906) circle (0.28pt);
\fill[slate] (12.031,6.155) circle (0.28pt);
\fill[slate] (12.031,6.356) circle (0.28pt);
\fill[slate] (12.031,6.812) circle (0.28pt);
\fill[slate] (12.031,6.404) circle (0.28pt);
\fill[slate] (12.031,6.653) circle (0.28pt);
\fill[slate] (12.031,4.509) circle (0.28pt);
\fill[slate] (12.031,4.758) circle (0.28pt);
\fill[slate] (12.031,4.958) circle (0.28pt);
\fill[slate] (12.031,5.414) circle (0.28pt);
\fill[slate] (12.031,5.658) circle (0.28pt);
\fill[slate] (12.031,5.87) circle (0.28pt);
\fill[slate] (12.031,6.326) circle (0.28pt);
\fill[slate] (12.031,6.57) circle (0.28pt);
\fill[slate] (12.031,6.783) circle (0.28pt);
\fill[slate] (12.031,6.818) circle (0.28pt);
\fill[slate] (12.031,5.906) circle (0.28pt);
\fill[slate] (12.031,6.155) circle (0.28pt);
\fill[slate] (12.031,6.357) circle (0.28pt);
\fill[slate] (12.031,6.813) circle (0.28pt);
\fill[slate] (12.031,6.404) circle (0.28pt);
\fill[slate] (12.031,6.653) circle (0.28pt);
\fill[slate] (12.031,5.007) circle (0.28pt);
\fill[slate] (12.031,5.256) circle (0.28pt);
\fill[slate] (12.031,5.456) circle (0.28pt);
\fill[slate] (12.031,5.912) circle (0.28pt);
\fill[slate] (12.031,6.156) circle (0.28pt);
\fill[slate] (12.031,6.368) circle (0.28pt);
\fill[slate] (12.031,6.825) circle (0.28pt);
\fill[slate] (12.031,6.404) circle (0.28pt);
\fill[slate] (12.031,6.653) circle (0.28pt);
\fill[slate] (12.031,5.505) circle (0.28pt);
\fill[slate] (12.031,5.754) circle (0.28pt);
\fill[slate] (12.031,5.954) circle (0.28pt);
\fill[slate] (12.031,6.41) circle (0.28pt);
\fill[slate] (12.031,6.654) circle (0.28pt);
\fill[slate] (12.031,6.003) circle (0.28pt);
\fill[slate] (12.031,6.252) circle (0.28pt);
\fill[slate] (12.031,6.453) circle (0.28pt);
\fill[slate] (12.031,6.501) circle (0.28pt);
\fill[slate] (12.031,6.75) circle (0.28pt);
\fill[slate] (1.832,1.977) circle (0.28pt);
\fill[slate] (1.832,2.433) circle (0.28pt);
\fill[slate] (1.83,2.676) circle (0.28pt);
\fill[slate] (1.832,2.889) circle (0.28pt);
\fill[slate] (1.832,3.345) circle (0.28pt);
\fill[slate] (1.832,3.589) circle (0.28pt);
\fill[slate] (1.832,3.801) circle (0.28pt);
\fill[slate] (1.832,4.257) circle (0.28pt);
\fill[slate] (1.832,4.501) circle (0.28pt);
\fill[slate] (1.832,4.713) circle (0.28pt);
\fill[slate] (1.832,5.17) circle (0.28pt);
\fill[slate] (1.832,5.413) circle (0.28pt);
\fill[slate] (1.832,5.626) circle (0.28pt);
\fill[slate] (1.832,6.082) circle (0.28pt);
\fill[slate] (1.832,6.325) circle (0.28pt);
\fill[slate] (1.832,6.538) circle (0.28pt);
\fill[slate] (1.832,6.573) circle (0.28pt);
\fill[slate] (1.832,6.822) circle (0.28pt);
\fill[slate] (1.832,5.661) circle (0.28pt);
\fill[slate] (1.832,5.91) circle (0.28pt);
\fill[slate] (1.832,6.112) circle (0.28pt);
\fill[slate] (1.832,6.568) circle (0.28pt);
\fill[slate] (1.832,6.812) circle (0.28pt);
\fill[slate] (1.832,6.159) circle (0.28pt);
\fill[slate] (1.832,6.408) circle (0.28pt);
\fill[slate] (1.832,6.609) circle (0.28pt);
\fill[slate] (1.832,6.657) circle (0.28pt);
\fill[slate] (1.832,4.749) circle (0.28pt);
\fill[slate] (1.832,4.998) circle (0.28pt);
\fill[slate] (1.832,5.2) circle (0.28pt);
\fill[slate] (1.832,5.656) circle (0.28pt);
\fill[slate] (1.832,5.9) circle (0.28pt);
\fill[slate] (1.832,6.112) circle (0.28pt);
\fill[slate] (1.832,6.568) circle (0.28pt);
\fill[slate] (1.832,6.812) circle (0.28pt);
\fill[slate] (1.832,6.148) circle (0.28pt);
\fill[slate] (1.832,6.397) circle (0.28pt);
\fill[slate] (1.832,6.599) circle (0.28pt);
\fill[slate] (1.832,6.646) circle (0.28pt);
\fill[slate] (1.832,5.247) circle (0.28pt);
\fill[slate] (1.832,5.496) circle (0.28pt);
\fill[slate] (1.832,5.697) circle (0.28pt);
\fill[slate] (1.832,6.152) circle (0.28pt);
\fill[slate] (1.832,6.396) circle (0.28pt);
\fill[slate] (1.832,6.609) circle (0.28pt);
\fill[slate] (1.832,6.644) circle (0.28pt);
\fill[slate] (1.832,5.745) circle (0.28pt);
\fill[slate] (1.832,5.994) circle (0.28pt);
\fill[slate] (1.832,6.194) circle (0.28pt);
\fill[slate] (1.832,6.65) circle (0.28pt);
\fill[slate] (1.832,6.243) circle (0.28pt);
\fill[slate] (1.832,6.492) circle (0.28pt);
\fill[slate] (1.832,6.692) circle (0.28pt);
\fill[slate] (1.832,6.741) circle (0.28pt);
\fill[slate] (1.832,3.837) circle (0.28pt);
\fill[slate] (1.832,4.086) circle (0.28pt);
\fill[slate] (1.832,4.288) circle (0.28pt);
\fill[slate] (1.832,4.744) circle (0.28pt);
\fill[slate] (1.832,4.988) circle (0.28pt);
\fill[slate] (1.832,5.2) circle (0.28pt);
\fill[slate] (1.832,5.656) circle (0.28pt);
\fill[slate] (1.832,5.899) circle (0.28pt);
\fill[slate] (1.832,6.112) circle (0.28pt);
\fill[slate] (1.832,6.568) circle (0.28pt);
\fill[slate] (1.832,6.812) circle (0.28pt);
\fill[slate] (1.832,6.148) circle (0.28pt);
\fill[slate] (1.832,6.396) circle (0.28pt);
\fill[slate] (1.832,6.599) circle (0.28pt);
\fill[slate] (1.832,6.645) circle (0.28pt);
\fill[slate] (1.832,5.236) circle (0.28pt);
\fill[slate] (1.832,5.485) circle (0.28pt);
\fill[slate] (1.832,5.687) circle (0.28pt);
\fill[slate] (1.832,6.143) circle (0.28pt);
\fill[slate] (1.832,6.387) circle (0.28pt);
\fill[slate] (1.832,6.599) circle (0.28pt);
\fill[slate] (1.832,6.635) circle (0.28pt);
\fill[slate] (1.832,5.734) circle (0.28pt);
\fill[slate] (1.832,5.983) circle (0.28pt);
\fill[slate] (1.832,6.183) circle (0.28pt);
\fill[slate] (1.832,6.639) circle (0.28pt);
\fill[slate] (1.832,6.232) circle (0.28pt);
\fill[slate] (1.832,6.481) circle (0.28pt);
\fill[slate] (1.832,6.681) circle (0.28pt);
\fill[slate] (1.832,6.73) circle (0.28pt);
\fill[slate] (1.832,4.335) circle (0.28pt);
\fill[slate] (1.832,4.584) circle (0.28pt);
\fill[slate] (1.832,4.784) circle (0.28pt);
\fill[slate] (1.832,5.24) circle (0.28pt);
\fill[slate] (1.832,5.484) circle (0.28pt);
\fill[slate] (1.832,5.696) circle (0.28pt);
\fill[slate] (1.832,6.152) circle (0.28pt);
\fill[slate] (1.832,6.396) circle (0.28pt);
\fill[slate] (1.832,6.609) circle (0.28pt);
\fill[slate] (1.832,6.644) circle (0.28pt);
\fill[slate] (1.832,5.732) circle (0.28pt);
\fill[slate] (1.832,5.981) circle (0.28pt);
\fill[slate] (1.832,6.183) circle (0.28pt);
\fill[slate] (1.832,6.639) circle (0.28pt);
\fill[slate] (1.832,6.23) circle (0.28pt);
\fill[slate] (1.832,6.479) circle (0.28pt);
\fill[slate] (1.832,6.68) circle (0.28pt);
\fill[slate] (1.832,6.728) circle (0.28pt);
\fill[slate] (1.832,4.833) circle (0.28pt);
\fill[slate] (1.832,5.082) circle (0.28pt);
\fill[slate] (1.832,5.282) circle (0.28pt);
\fill[slate] (1.832,5.738) circle (0.28pt);
\fill[slate] (1.832,5.982) circle (0.28pt);
\fill[slate] (1.832,6.194) circle (0.28pt);
\fill[slate] (1.832,6.65) circle (0.28pt);
\fill[slate] (1.832,6.23) circle (0.28pt);
\fill[slate] (1.832,6.479) circle (0.28pt);
\fill[slate] (1.832,6.681) circle (0.28pt);
\fill[slate] (1.832,6.728) circle (0.28pt);
\fill[slate] (1.832,5.331) circle (0.28pt);
\fill[slate] (1.832,5.58) circle (0.28pt);
\fill[slate] (1.832,5.78) circle (0.28pt);
\fill[slate] (1.832,6.236) circle (0.28pt);
\fill[slate] (1.832,6.48) circle (0.28pt);
\fill[slate] (1.832,6.692) circle (0.28pt);
\fill[slate] (1.832,6.728) circle (0.28pt);
\fill[slate] (1.832,5.829) circle (0.28pt);
\fill[slate] (1.832,6.078) circle (0.28pt);
\fill[slate] (1.832,6.278) circle (0.28pt);
\fill[slate] (1.832,6.734) circle (0.28pt);
\fill[slate] (1.832,6.327) circle (0.28pt);
\fill[slate] (1.832,6.576) circle (0.28pt);
\fill[slate] (1.832,6.776) circle (0.28pt);
\fill[slate] (1.832,6.825) circle (0.28pt);
\fill[slate] (1.83,3.173) circle (0.28pt);
\fill[slate] (1.83,3.376) circle (0.28pt);
\fill[slate] (1.83,3.832) circle (0.28pt);
\fill[slate] (1.83,4.075) circle (0.28pt);
\fill[slate] (1.83,4.288) circle (0.28pt);
\fill[slate] (1.83,4.744) circle (0.28pt);
\fill[slate] (1.83,4.987) circle (0.28pt);
\fill[slate] (1.83,5.2) circle (0.28pt);
\fill[slate] (1.83,5.656) circle (0.28pt);
\fill[slate] (1.83,5.899) circle (0.28pt);
\fill[slate] (1.83,6.112) circle (0.28pt);
\fill[slate] (1.83,6.568) circle (0.28pt);
\fill[slate] (1.83,6.812) circle (0.28pt);
\fill[slate] (1.83,6.148) circle (0.28pt);
\fill[slate] (1.83,6.396) circle (0.28pt);
\fill[slate] (1.83,6.599) circle (0.28pt);
\fill[slate] (1.83,6.645) circle (0.28pt);
\fill[slate] (1.83,5.235) circle (0.28pt);
\fill[slate] (1.83,5.484) circle (0.28pt);
\fill[slate] (1.83,5.687) circle (0.28pt);
\fill[slate] (1.83,6.143) circle (0.28pt);
\fill[slate] (1.83,6.386) circle (0.28pt);
\fill[slate] (1.83,6.599) circle (0.28pt);
\fill[slate] (1.83,6.634) circle (0.28pt);
\fill[slate] (1.83,5.733) circle (0.28pt);
\fill[slate] (1.83,5.982) circle (0.28pt);
\fill[slate] (1.83,6.183) circle (0.28pt);
\fill[slate] (1.83,6.639) circle (0.28pt);
\fill[slate] (1.83,6.231) circle (0.28pt);
\fill[slate] (1.83,6.48) circle (0.28pt);
\fill[slate] (1.83,6.681) circle (0.28pt);
\fill[slate] (1.83,6.729) circle (0.28pt);
\fill[slate] (1.83,4.572) circle (0.28pt);
\fill[slate] (1.83,4.775) circle (0.28pt);
\fill[slate] (1.83,5.231) circle (0.28pt);
\fill[slate] (1.83,5.474) circle (0.28pt);
\fill[slate] (1.83,5.687) circle (0.28pt);
\fill[slate] (1.83,6.143) circle (0.28pt);
\fill[slate] (1.83,6.386) circle (0.28pt);
\fill[slate] (1.83,6.599) circle (0.28pt);
\fill[slate] (1.83,6.634) circle (0.28pt);
\fill[slate] (1.83,5.971) circle (0.28pt);
\fill[slate] (1.83,6.174) circle (0.28pt);
\fill[slate] (1.83,6.63) circle (0.28pt);
\fill[slate] (1.83,6.22) circle (0.28pt);
\fill[slate] (1.83,6.469) circle (0.28pt);
\fill[slate] (1.83,6.67) circle (0.28pt);
\fill[slate] (1.83,6.718) circle (0.28pt);
\fill[slate] (1.83,4.821) circle (0.28pt);
\fill[slate] (1.83,5.07) circle (0.28pt);
\fill[slate] (1.83,5.271) circle (0.28pt);
\fill[slate] (1.83,5.971) circle (0.28pt);
\fill[slate] (1.83,6.183) circle (0.28pt);
\fill[slate] (1.83,6.639) circle (0.28pt);
\fill[slate] (1.83,6.219) circle (0.28pt);
\fill[slate] (1.83,6.468) circle (0.28pt);
\fill[slate] (1.83,6.67) circle (0.28pt);
\fill[slate] (1.83,6.717) circle (0.28pt);
\fill[slate] (1.83,5.319) circle (0.28pt);
\fill[slate] (1.83,5.568) circle (0.28pt);
\fill[slate] (1.83,5.769) circle (0.28pt);
\fill[slate] (1.83,6.225) circle (0.28pt);
\fill[slate] (1.83,6.469) circle (0.28pt);
\fill[slate] (1.83,6.681) circle (0.28pt);
\fill[slate] (1.83,6.717) circle (0.28pt);
\fill[slate] (1.83,5.817) circle (0.28pt);
\fill[slate] (1.83,6.067) circle (0.28pt);
\fill[slate] (1.83,6.267) circle (0.28pt);
\fill[slate] (1.83,6.723) circle (0.28pt);
\fill[slate] (1.83,6.316) circle (0.28pt);
\fill[slate] (1.83,6.565) circle (0.28pt);
\fill[slate] (1.83,6.765) circle (0.28pt);
\fill[slate] (1.83,6.814) circle (0.28pt);
\fill[slate] (1.83,3.422) circle (0.28pt);
\fill[slate] (1.83,3.671) circle (0.28pt);
\fill[slate] (1.83,3.872) circle (0.28pt);
\fill[slate] (1.83,4.572) circle (0.28pt);
\fill[slate] (1.83,4.784) circle (0.28pt);
\fill[slate] (1.83,5.24) circle (0.28pt);
\fill[slate] (1.83,5.484) circle (0.28pt);
\fill[slate] (1.83,5.696) circle (0.28pt);
\fill[slate] (1.83,6.152) circle (0.28pt);
\fill[slate] (1.83,6.396) circle (0.28pt);
\fill[slate] (1.83,6.609) circle (0.28pt);
\fill[slate] (1.83,6.644) circle (0.28pt);
\fill[slate] (1.83,5.981) circle (0.28pt);
\fill[slate] (1.83,6.183) circle (0.28pt);
\fill[slate] (1.83,6.639) circle (0.28pt);
\fill[slate] (1.83,6.23) circle (0.28pt);
\fill[slate] (1.83,6.479) circle (0.28pt);
\fill[slate] (1.83,6.679) circle (0.28pt);
\fill[slate] (1.83,6.728) circle (0.28pt);
\fill[slate] (1.83,4.82) circle (0.28pt);
\fill[slate] (1.83,5.069) circle (0.28pt);
\fill[slate] (1.83,5.271) circle (0.28pt);
\fill[slate] (1.83,5.971) circle (0.28pt);
\fill[slate] (1.83,6.183) circle (0.28pt);
\fill[slate] (1.83,6.639) circle (0.28pt);
\fill[slate] (1.83,6.219) circle (0.28pt);
\fill[slate] (1.83,6.468) circle (0.28pt);
\fill[slate] (1.83,6.67) circle (0.28pt);
\fill[slate] (1.83,6.717) circle (0.28pt);
\fill[slate] (1.83,5.318) circle (0.28pt);
\fill[slate] (1.83,5.567) circle (0.28pt);
\fill[slate] (1.83,5.768) circle (0.28pt);
\fill[slate] (1.83,6.224) circle (0.28pt);
\fill[slate] (1.83,6.467) circle (0.28pt);
\fill[slate] (1.83,6.68) circle (0.28pt);
\fill[slate] (1.83,6.715) circle (0.28pt);
\fill[slate] (1.83,5.816) circle (0.28pt);
\fill[slate] (1.83,6.065) circle (0.28pt);
\fill[slate] (1.83,6.265) circle (0.28pt);
\fill[slate] (1.83,6.721) circle (0.28pt);
\fill[slate] (1.83,6.314) circle (0.28pt);
\fill[slate] (1.83,6.563) circle (0.28pt);
\fill[slate] (1.83,6.763) circle (0.28pt);
\fill[slate] (1.83,6.812) circle (0.28pt);
\fill[slate] (1.83,3.92) circle (0.28pt);
\fill[slate] (1.83,4.169) circle (0.28pt);
\fill[slate] (1.83,4.37) circle (0.28pt);
\fill[slate] (1.83,4.826) circle (0.28pt);
\fill[slate] (1.83,5.07) circle (0.28pt);
\fill[slate] (1.83,5.282) circle (0.28pt);
\fill[slate] (1.83,5.738) circle (0.28pt);
\fill[slate] (1.83,5.981) circle (0.28pt);
\fill[slate] (1.83,6.194) circle (0.28pt);
\fill[slate] (1.83,6.65) circle (0.28pt);
\fill[slate] (1.83,6.23) circle (0.28pt);
\fill[slate] (1.83,6.478) circle (0.28pt);
\fill[slate] (1.83,6.681) circle (0.28pt);
\fill[slate] (1.83,6.727) circle (0.28pt);
\fill[slate] (1.83,5.318) circle (0.28pt);
\fill[slate] (1.83,5.567) circle (0.28pt);
\fill[slate] (1.83,5.769) circle (0.28pt);
\fill[slate] (1.83,6.225) circle (0.28pt);
\fill[slate] (1.83,6.469) circle (0.28pt);
\fill[slate] (1.83,6.681) circle (0.28pt);
\fill[slate] (1.83,6.717) circle (0.28pt);
\fill[slate] (1.83,5.816) circle (0.28pt);
\fill[slate] (1.83,6.065) circle (0.28pt);
\fill[slate] (1.83,6.265) circle (0.28pt);
\fill[slate] (1.83,6.721) circle (0.28pt);
\fill[slate] (1.83,6.314) circle (0.28pt);
\fill[slate] (1.83,6.563) circle (0.28pt);
\fill[slate] (1.83,6.763) circle (0.28pt);
\fill[slate] (1.83,6.812) circle (0.28pt);
\fill[slate] (1.83,4.418) circle (0.28pt);
\fill[slate] (1.83,4.667) circle (0.28pt);
\fill[slate] (1.83,4.868) circle (0.28pt);
\fill[slate] (1.83,5.324) circle (0.28pt);
\fill[slate] (1.83,5.568) circle (0.28pt);
\fill[slate] (1.83,5.78) circle (0.28pt);
\fill[slate] (1.83,6.236) circle (0.28pt);
\fill[slate] (1.83,6.479) circle (0.28pt);
\fill[slate] (1.83,6.692) circle (0.28pt);
\fill[slate] (1.83,6.727) circle (0.28pt);
\fill[slate] (1.83,5.816) circle (0.28pt);
\fill[slate] (1.83,6.065) circle (0.28pt);
\fill[slate] (1.83,6.267) circle (0.28pt);
\fill[slate] (1.83,6.723) circle (0.28pt);
\fill[slate] (1.83,6.314) circle (0.28pt);
\fill[slate] (1.83,6.563) circle (0.28pt);
\fill[slate] (1.83,6.763) circle (0.28pt);
\fill[slate] (1.83,6.812) circle (0.28pt);
\fill[slate] (1.83,4.916) circle (0.28pt);
\fill[slate] (1.83,5.165) circle (0.28pt);
\fill[slate] (1.83,5.366) circle (0.28pt);
\fill[slate] (1.83,5.822) circle (0.28pt);
\fill[slate] (1.83,6.066) circle (0.28pt);
\fill[slate] (1.83,6.278) circle (0.28pt);
\fill[slate] (1.83,6.734) circle (0.28pt);
\fill[slate] (1.83,6.314) circle (0.28pt);
\fill[slate] (1.83,6.563) circle (0.28pt);
\fill[slate] (1.83,6.765) circle (0.28pt);
\fill[slate] (1.83,6.812) circle (0.28pt);
\fill[slate] (1.83,5.414) circle (0.28pt);
\fill[slate] (1.83,5.663) circle (0.28pt);
\fill[slate] (1.83,5.864) circle (0.28pt);
\fill[slate] (1.83,6.32) circle (0.28pt);
\fill[slate] (1.83,6.564) circle (0.28pt);
\fill[slate] (1.83,6.776) circle (0.28pt);
\fill[slate] (1.83,6.812) circle (0.28pt);
\fill[slate] (1.83,5.913) circle (0.28pt);
\fill[slate] (1.83,6.162) circle (0.28pt);
\fill[slate] (1.83,6.362) circle (0.28pt);
\fill[slate] (1.83,6.818) circle (0.28pt);
\fill[slate] (1.83,6.411) circle (0.28pt);
\fill[slate] (1.83,6.66) circle (0.28pt);
\fill[slate] (1.785,3.376) circle (0.28pt);
\fill[slate] (1.785,3.832) circle (0.28pt);
\fill[slate] (1.785,4.075) circle (0.28pt);
\fill[slate] (1.785,4.288) circle (0.28pt);
\fill[slate] (1.785,4.744) circle (0.28pt);
\fill[slate] (1.785,4.988) circle (0.28pt);
\fill[slate] (1.785,5.2) circle (0.28pt);
\fill[slate] (1.785,5.656) circle (0.28pt);
\fill[slate] (1.785,5.9) circle (0.28pt);
\fill[slate] (1.785,6.112) circle (0.28pt);
\fill[slate] (1.785,6.569) circle (0.28pt);
\fill[slate] (1.785,6.812) circle (0.28pt);
\fill[slate] (1.785,6.148) circle (0.28pt);
\fill[slate] (1.785,6.397) circle (0.28pt);
\fill[slate] (1.785,6.599) circle (0.28pt);
\fill[slate] (1.785,6.646) circle (0.28pt);
\fill[slate] (1.785,5.236) circle (0.28pt);
\fill[slate] (1.785,5.485) circle (0.28pt);
\fill[slate] (1.785,5.687) circle (0.28pt);
\fill[slate] (1.785,6.143) circle (0.28pt);
\fill[slate] (1.785,6.387) circle (0.28pt);
\fill[slate] (1.785,6.599) circle (0.28pt);
\fill[slate] (1.785,6.635) circle (0.28pt);
\fill[slate] (1.785,5.734) circle (0.28pt);
\fill[slate] (1.785,5.983) circle (0.28pt);
\fill[slate] (1.785,6.183) circle (0.28pt);
\fill[slate] (1.785,6.639) circle (0.28pt);
\fill[slate] (1.785,6.232) circle (0.28pt);
\fill[slate] (1.785,6.481) circle (0.28pt);
\fill[slate] (1.785,6.681) circle (0.28pt);
\fill[slate] (1.785,6.73) circle (0.28pt);
\fill[slate] (1.785,4.572) circle (0.28pt);
\fill[slate] (1.785,4.775) circle (0.28pt);
\fill[slate] (1.785,5.231) circle (0.28pt);
\fill[slate] (1.785,5.474) circle (0.28pt);
\fill[slate] (1.785,5.687) circle (0.28pt);
\fill[slate] (1.785,6.143) circle (0.28pt);
\fill[slate] (1.785,6.386) circle (0.28pt);
\fill[slate] (1.785,6.599) circle (0.28pt);
\fill[slate] (1.785,6.634) circle (0.28pt);
\fill[slate] (1.785,5.971) circle (0.28pt);
\fill[slate] (1.785,6.174) circle (0.28pt);
\fill[slate] (1.785,6.63) circle (0.28pt);
\fill[slate] (1.785,6.22) circle (0.28pt);
\fill[slate] (1.785,6.469) circle (0.28pt);
\fill[slate] (1.785,6.67) circle (0.28pt);
\fill[slate] (1.785,6.718) circle (0.28pt);
\fill[slate] (1.785,4.821) circle (0.28pt);
\fill[slate] (1.785,5.07) circle (0.28pt);
\fill[slate] (1.785,5.271) circle (0.28pt);
\fill[slate] (1.785,5.971) circle (0.28pt);
\fill[slate] (1.785,6.183) circle (0.28pt);
\fill[slate] (1.785,6.639) circle (0.28pt);
\fill[slate] (1.785,6.219) circle (0.28pt);
\fill[slate] (1.785,6.468) circle (0.28pt);
\fill[slate] (1.785,6.67) circle (0.28pt);
\fill[slate] (1.785,6.717) circle (0.28pt);
\fill[slate] (1.785,5.319) circle (0.28pt);
\fill[slate] (1.785,5.568) circle (0.28pt);
\fill[slate] (1.785,5.769) circle (0.28pt);
\fill[slate] (1.785,6.225) circle (0.28pt);
\fill[slate] (1.785,6.469) circle (0.28pt);
\fill[slate] (1.785,6.681) circle (0.28pt);
\fill[slate] (1.785,6.717) circle (0.28pt);
\fill[slate] (1.785,5.817) circle (0.28pt);
\fill[slate] (1.785,6.066) circle (0.28pt);
\fill[slate] (1.785,6.267) circle (0.28pt);
\fill[slate] (1.785,6.723) circle (0.28pt);
\fill[slate] (1.785,6.315) circle (0.28pt);
\fill[slate] (1.785,6.564) circle (0.28pt);
\fill[slate] (1.785,6.765) circle (0.28pt);
\fill[slate] (1.785,6.813) circle (0.28pt);
\fill[slate] (1.784,3.661) circle (0.28pt);
\fill[slate] (1.785,4.775) circle (0.28pt);
\fill[slate] (1.785,5.231) circle (0.28pt);
\fill[slate] (1.785,5.474) circle (0.28pt);
\fill[slate] (1.785,5.687) circle (0.28pt);
\fill[slate] (1.785,6.143) circle (0.28pt);
\fill[slate] (1.785,6.387) circle (0.28pt);
\fill[slate] (1.785,6.599) circle (0.28pt);
\fill[slate] (1.785,6.635) circle (0.28pt);
\fill[slate] (1.785,5.971) circle (0.28pt);
\fill[slate] (1.785,6.174) circle (0.28pt);
\fill[slate] (1.785,6.63) circle (0.28pt);
\fill[slate] (1.785,6.22) circle (0.28pt);
\fill[slate] (1.785,6.469) circle (0.28pt);
\fill[slate] (1.785,6.67) circle (0.28pt);
\fill[slate] (1.785,6.718) circle (0.28pt);
\fill[slate] (1.785,5.06) circle (0.28pt);
\fill[slate] (1.785,6.174) circle (0.28pt);
\fill[slate] (1.785,6.63) circle (0.28pt);
\fill[slate] (1.785,6.459) circle (0.28pt);
\fill[slate] (1.785,6.708) circle (0.28pt);
\fill[slate] (1.785,5.309) circle (0.28pt);
\fill[slate] (1.785,5.558) circle (0.28pt);
\fill[slate] (1.785,5.758) circle (0.28pt);
\fill[slate] (1.785,6.214) circle (0.28pt);
\fill[slate] (1.785,6.67) circle (0.28pt);
\fill[slate] (1.785,6.706) circle (0.28pt);
\fill[slate] (1.785,5.807) circle (0.28pt);
\fill[slate] (1.785,6.056) circle (0.28pt);
\fill[slate] (1.785,6.256) circle (0.28pt);
\fill[slate] (1.785,6.712) circle (0.28pt);
\fill[slate] (1.785,6.305) circle (0.28pt);
\fill[slate] (1.785,6.554) circle (0.28pt);
\fill[slate] (1.785,6.754) circle (0.28pt);
\fill[slate] (1.785,6.803) circle (0.28pt);
\fill[slate] (1.784,3.91) circle (0.28pt);
\fill[slate] (1.784,4.159) circle (0.28pt);
\fill[slate] (1.784,4.815) circle (0.28pt);
\fill[slate] (1.784,5.271) circle (0.28pt);
\fill[slate] (1.784,5.727) circle (0.28pt);
\fill[slate] (1.784,5.971) circle (0.28pt);
\fill[slate] (1.784,6.184) circle (0.28pt);
\fill[slate] (1.784,6.64) circle (0.28pt);
\fill[slate] (1.784,6.219) circle (0.28pt);
\fill[slate] (1.784,6.468) circle (0.28pt);
\fill[slate] (1.784,6.67) circle (0.28pt);
\fill[slate] (1.784,6.717) circle (0.28pt);
\fill[slate] (1.784,5.307) circle (0.28pt);
\fill[slate] (1.784,5.556) circle (0.28pt);
\fill[slate] (1.784,6.214) circle (0.28pt);
\fill[slate] (1.784,6.67) circle (0.28pt);
\fill[slate] (1.784,5.805) circle (0.28pt);
\fill[slate] (1.784,6.054) circle (0.28pt);
\fill[slate] (1.784,6.255) circle (0.28pt);
\fill[slate] (1.784,6.711) circle (0.28pt);
\fill[slate] (1.784,6.303) circle (0.28pt);
\fill[slate] (1.784,6.552) circle (0.28pt);
\fill[slate] (1.784,6.753) circle (0.28pt);
\fill[slate] (1.784,6.801) circle (0.28pt);
\fill[slate] (1.784,4.408) circle (0.28pt);
\fill[slate] (1.784,4.657) circle (0.28pt);
\fill[slate] (1.784,4.857) circle (0.28pt);
\fill[slate] (1.784,5.313) circle (0.28pt);
\fill[slate] (1.784,5.557) circle (0.28pt);
\fill[slate] (1.784,5.769) circle (0.28pt);
\fill[slate] (1.784,6.225) circle (0.28pt);
\fill[slate] (1.784,6.469) circle (0.28pt);
\fill[slate] (1.784,6.681) circle (0.28pt);
\fill[slate] (1.784,6.717) circle (0.28pt);
\fill[slate] (1.784,5.805) circle (0.28pt);
\fill[slate] (1.784,6.054) circle (0.28pt);
\fill[slate] (1.784,6.256) circle (0.28pt);
\fill[slate] (1.784,6.712) circle (0.28pt);
\fill[slate] (1.784,6.303) circle (0.28pt);
\fill[slate] (1.784,6.552) circle (0.28pt);
\fill[slate] (1.784,6.753) circle (0.28pt);
\fill[slate] (1.784,6.801) circle (0.28pt);
\fill[slate] (1.784,4.906) circle (0.28pt);
\fill[slate] (1.784,5.155) circle (0.28pt);
\fill[slate] (1.784,5.355) circle (0.28pt);
\fill[slate] (1.784,5.811) circle (0.28pt);
\fill[slate] (1.784,6.055) circle (0.28pt);
\fill[slate] (1.784,6.267) circle (0.28pt);
\fill[slate] (1.784,6.723) circle (0.28pt);
\fill[slate] (1.784,6.303) circle (0.28pt);
\fill[slate] (1.784,6.552) circle (0.28pt);
\fill[slate] (1.784,6.754) circle (0.28pt);
\fill[slate] (1.784,6.801) circle (0.28pt);
\fill[slate] (1.784,5.404) circle (0.28pt);
\fill[slate] (1.784,5.653) circle (0.28pt);
\fill[slate] (1.784,5.853) circle (0.28pt);
\fill[slate] (1.784,6.309) circle (0.28pt);
\fill[slate] (1.784,6.553) circle (0.28pt);
\fill[slate] (1.784,6.765) circle (0.28pt);
\fill[slate] (1.784,6.801) circle (0.28pt);
\fill[slate] (1.784,5.902) circle (0.28pt);
\fill[slate] (1.784,6.151) circle (0.28pt);
\fill[slate] (1.784,6.351) circle (0.28pt);
\fill[slate] (1.784,6.807) circle (0.28pt);
\fill[slate] (1.784,6.4) circle (0.28pt);
\fill[slate] (1.784,6.649) circle (0.28pt);
\fill[slate] (1.784,6.849) circle (0.28pt);
\fill[slate] (1.777,2.511) circle (0.28pt);
\fill[slate] (1.777,2.76) circle (0.28pt);
\fill[slate] (1.778,3.416) circle (0.28pt);
\fill[slate] (1.778,3.872) circle (0.28pt);
\fill[slate] (1.778,4.328) circle (0.28pt);
\fill[slate] (1.778,4.572) circle (0.28pt);
\fill[slate] (1.778,4.785) circle (0.28pt);
\fill[slate] (1.778,5.241) circle (0.28pt);
\fill[slate] (1.778,5.484) circle (0.28pt);
\fill[slate] (1.778,5.697) circle (0.28pt);
\fill[slate] (1.778,6.153) circle (0.28pt);
\fill[slate] (1.778,6.396) circle (0.28pt);
\fill[slate] (1.778,6.609) circle (0.28pt);
\fill[slate] (1.778,6.644) circle (0.28pt);
\fill[slate] (1.778,5.732) circle (0.28pt);
\fill[slate] (1.778,5.981) circle (0.28pt);
\fill[slate] (1.778,6.183) circle (0.28pt);
\fill[slate] (1.778,6.639) circle (0.28pt);
\fill[slate] (1.778,6.23) circle (0.28pt);
\fill[slate] (1.778,6.479) circle (0.28pt);
\fill[slate] (1.778,6.68) circle (0.28pt);
\fill[slate] (1.778,6.728) circle (0.28pt);
\fill[slate] (1.778,4.82) circle (0.28pt);
\fill[slate] (1.778,5.069) circle (0.28pt);
\fill[slate] (1.778,5.271) circle (0.28pt);
\fill[slate] (1.778,5.727) circle (0.28pt);
\fill[slate] (1.778,5.971) circle (0.28pt);
\fill[slate] (1.778,6.183) circle (0.28pt);
\fill[slate] (1.778,6.639) circle (0.28pt);
\fill[slate] (1.778,6.219) circle (0.28pt);
\fill[slate] (1.778,6.468) circle (0.28pt);
\fill[slate] (1.778,6.67) circle (0.28pt);
\fill[slate] (1.778,6.717) circle (0.28pt);
\fill[slate] (1.778,5.318) circle (0.28pt);
\fill[slate] (1.778,5.567) circle (0.28pt);
\fill[slate] (1.778,5.768) circle (0.28pt);
\fill[slate] (1.778,6.224) circle (0.28pt);
\fill[slate] (1.778,6.467) circle (0.28pt);
\fill[slate] (1.778,6.68) circle (0.28pt);
\fill[slate] (1.778,6.715) circle (0.28pt);
\fill[slate] (1.778,5.816) circle (0.28pt);
\fill[slate] (1.778,6.065) circle (0.28pt);
\fill[slate] (1.778,6.265) circle (0.28pt);
\fill[slate] (1.778,6.721) circle (0.28pt);
\fill[slate] (1.778,6.314) circle (0.28pt);
\fill[slate] (1.778,6.563) circle (0.28pt);
\fill[slate] (1.778,6.763) circle (0.28pt);
\fill[slate] (1.778,6.812) circle (0.28pt);
\fill[slate] (1.778,4.815) circle (0.28pt);
\fill[slate] (1.778,5.059) circle (0.28pt);
\fill[slate] (1.778,5.271) circle (0.28pt);
\fill[slate] (1.778,5.728) circle (0.28pt);
\fill[slate] (1.778,5.971) circle (0.28pt);
\fill[slate] (1.778,6.184) circle (0.28pt);
\fill[slate] (1.778,6.64) circle (0.28pt);
\fill[slate] (1.778,6.219) circle (0.28pt);
\fill[slate] (1.778,6.468) circle (0.28pt);
\fill[slate] (1.778,6.67) circle (0.28pt);
\fill[slate] (1.778,6.717) circle (0.28pt);
\fill[slate] (1.778,5.556) circle (0.28pt);
\fill[slate] (1.778,5.758) circle (0.28pt);
\fill[slate] (1.778,6.214) circle (0.28pt);
\fill[slate] (1.778,6.458) circle (0.28pt);
\fill[slate] (1.778,6.67) circle (0.28pt);
\fill[slate] (1.778,6.706) circle (0.28pt);
\fill[slate] (1.778,5.805) circle (0.28pt);
\fill[slate] (1.778,6.054) circle (0.28pt);
\fill[slate] (1.778,6.711) circle (0.28pt);
\fill[slate] (1.778,6.303) circle (0.28pt);
\fill[slate] (1.778,6.552) circle (0.28pt);
\fill[slate] (1.778,6.753) circle (0.28pt);
\fill[slate] (1.778,6.801) circle (0.28pt);
\fill[slate] (1.778,4.406) circle (0.28pt);
\fill[slate] (1.778,4.655) circle (0.28pt);
\fill[slate] (1.778,5.312) circle (0.28pt);
\fill[slate] (1.778,5.768) circle (0.28pt);
\fill[slate] (1.778,6.224) circle (0.28pt);
\fill[slate] (1.778,6.467) circle (0.28pt);
\fill[slate] (1.778,6.68) circle (0.28pt);
\fill[slate] (1.778,6.715) circle (0.28pt);
\fill[slate] (1.778,6.255) circle (0.28pt);
\fill[slate] (1.778,6.711) circle (0.28pt);
\fill[slate] (1.778,6.301) circle (0.28pt);
\fill[slate] (1.778,6.55) circle (0.28pt);
\fill[slate] (1.778,6.799) circle (0.28pt);
\fill[slate] (1.778,4.904) circle (0.28pt);
\fill[slate] (1.778,5.153) circle (0.28pt);
\fill[slate] (1.778,5.354) circle (0.28pt);
\fill[slate] (1.778,5.81) circle (0.28pt);
\fill[slate] (1.778,6.266) circle (0.28pt);
\fill[slate] (1.778,6.722) circle (0.28pt);
\fill[slate] (1.778,6.301) circle (0.28pt);
\fill[slate] (1.778,6.55) circle (0.28pt);
\fill[slate] (1.778,6.753) circle (0.28pt);
\fill[slate] (1.778,6.799) circle (0.28pt);
\fill[slate] (1.778,5.402) circle (0.28pt);
\fill[slate] (1.778,5.651) circle (0.28pt);
\fill[slate] (1.778,5.852) circle (0.28pt);
\fill[slate] (1.778,6.308) circle (0.28pt);
\fill[slate] (1.778,6.551) circle (0.28pt);
\fill[slate] (1.778,6.764) circle (0.28pt);
\fill[slate] (1.778,6.799) circle (0.28pt);
\fill[slate] (1.778,5.9) circle (0.28pt);
\fill[slate] (1.778,6.149) circle (0.28pt);
\fill[slate] (1.778,6.35) circle (0.28pt);
\fill[slate] (1.778,6.806) circle (0.28pt);
\fill[slate] (1.778,6.398) circle (0.28pt);
\fill[slate] (1.778,6.647) circle (0.28pt);
\fill[slate] (1.778,6.848) circle (0.28pt);
\fill[slate] (1.777,3.009) circle (0.28pt);
\fill[slate] (1.777,3.258) circle (0.28pt);
\fill[slate] (1.777,3.458) circle (0.28pt);
\fill[slate] (1.777,3.914) circle (0.28pt);
\fill[slate] (1.777,4.37) circle (0.28pt);
\fill[slate] (1.777,4.826) circle (0.28pt);
\fill[slate] (1.777,5.07) circle (0.28pt);
\fill[slate] (1.777,5.282) circle (0.28pt);
\fill[slate] (1.777,5.738) circle (0.28pt);
\fill[slate] (1.777,5.982) circle (0.28pt);
\fill[slate] (1.777,6.194) circle (0.28pt);
\fill[slate] (1.777,6.651) circle (0.28pt);
\fill[slate] (1.777,6.23) circle (0.28pt);
\fill[slate] (1.777,6.479) circle (0.28pt);
\fill[slate] (1.777,6.681) circle (0.28pt);
\fill[slate] (1.777,6.728) circle (0.28pt);
\fill[slate] (1.777,5.318) circle (0.28pt);
\fill[slate] (1.777,5.567) circle (0.28pt);
\fill[slate] (1.777,5.769) circle (0.28pt);
\fill[slate] (1.777,6.225) circle (0.28pt);
\fill[slate] (1.777,6.469) circle (0.28pt);
\fill[slate] (1.777,6.681) circle (0.28pt);
\fill[slate] (1.777,6.717) circle (0.28pt);
\fill[slate] (1.777,5.816) circle (0.28pt);
\fill[slate] (1.777,6.065) circle (0.28pt);
\fill[slate] (1.777,6.265) circle (0.28pt);
\fill[slate] (1.777,6.721) circle (0.28pt);
\fill[slate] (1.777,6.314) circle (0.28pt);
\fill[slate] (1.777,6.563) circle (0.28pt);
\fill[slate] (1.777,6.763) circle (0.28pt);
\fill[slate] (1.777,6.812) circle (0.28pt);
\fill[slate] (1.777,4.406) circle (0.28pt);
\fill[slate] (1.777,4.655) circle (0.28pt);
\fill[slate] (1.777,4.857) circle (0.28pt);
\fill[slate] (1.777,5.313) circle (0.28pt);
\fill[slate] (1.777,5.557) circle (0.28pt);
\fill[slate] (1.777,5.769) circle (0.28pt);
\fill[slate] (1.777,6.225) circle (0.28pt);
\fill[slate] (1.777,6.469) circle (0.28pt);
\fill[slate] (1.777,6.681) circle (0.28pt);
\fill[slate] (1.777,6.717) circle (0.28pt);
\fill[slate] (1.777,5.805) circle (0.28pt);
\fill[slate] (1.777,6.256) circle (0.28pt);
\fill[slate] (1.777,6.712) circle (0.28pt);
\fill[slate] (1.777,6.303) circle (0.28pt);
\fill[slate] (1.777,6.552) circle (0.28pt);
\fill[slate] (1.777,6.753) circle (0.28pt);
\fill[slate] (1.777,6.801) circle (0.28pt);
\fill[slate] (1.777,4.904) circle (0.28pt);
\fill[slate] (1.777,5.153) circle (0.28pt);
\fill[slate] (1.777,5.354) circle (0.28pt);
\fill[slate] (1.777,5.81) circle (0.28pt);
\fill[slate] (1.777,6.266) circle (0.28pt);
\fill[slate] (1.777,6.722) circle (0.28pt);
\fill[slate] (1.777,6.301) circle (0.28pt);
\fill[slate] (1.777,6.55) circle (0.28pt);
\fill[slate] (1.777,6.753) circle (0.28pt);
\fill[slate] (1.777,6.799) circle (0.28pt);
\fill[slate] (1.777,5.402) circle (0.28pt);
\fill[slate] (1.777,5.651) circle (0.28pt);
\fill[slate] (1.777,5.851) circle (0.28pt);
\fill[slate] (1.777,6.307) circle (0.28pt);
\fill[slate] (1.777,6.551) circle (0.28pt);
\fill[slate] (1.777,6.763) circle (0.28pt);
\fill[slate] (1.777,6.799) circle (0.28pt);
\fill[slate] (1.777,5.9) circle (0.28pt);
\fill[slate] (1.777,6.149) circle (0.28pt);
\fill[slate] (1.777,6.349) circle (0.28pt);
\fill[slate] (1.777,6.805) circle (0.28pt);
\fill[slate] (1.777,6.398) circle (0.28pt);
\fill[slate] (1.777,6.647) circle (0.28pt);
\fill[slate] (1.777,6.847) circle (0.28pt);
\fill[slate] (1.777,3.507) circle (0.28pt);
\fill[slate] (1.777,3.756) circle (0.28pt);
\fill[slate] (1.777,3.956) circle (0.28pt);
\fill[slate] (1.777,4.412) circle (0.28pt);
\fill[slate] (1.777,4.656) circle (0.28pt);
\fill[slate] (1.777,4.868) circle (0.28pt);
\fill[slate] (1.777,5.324) circle (0.28pt);
\fill[slate] (1.777,5.568) circle (0.28pt);
\fill[slate] (1.777,5.78) circle (0.28pt);
\fill[slate] (1.777,6.236) circle (0.28pt);
\fill[slate] (1.777,6.48) circle (0.28pt);
\fill[slate] (1.777,6.692) circle (0.28pt);
\fill[slate] (1.777,6.728) circle (0.28pt);
\fill[slate] (1.777,5.816) circle (0.28pt);
\fill[slate] (1.777,6.065) circle (0.28pt);
\fill[slate] (1.777,6.267) circle (0.28pt);
\fill[slate] (1.777,6.723) circle (0.28pt);
\fill[slate] (1.777,6.314) circle (0.28pt);
\fill[slate] (1.777,6.563) circle (0.28pt);
\fill[slate] (1.777,6.763) circle (0.28pt);
\fill[slate] (1.777,6.812) circle (0.28pt);
\fill[slate] (1.777,4.904) circle (0.28pt);
\fill[slate] (1.777,5.153) circle (0.28pt);
\fill[slate] (1.777,5.355) circle (0.28pt);
\fill[slate] (1.777,5.811) circle (0.28pt);
\fill[slate] (1.777,6.055) circle (0.28pt);
\fill[slate] (1.777,6.267) circle (0.28pt);
\fill[slate] (1.777,6.723) circle (0.28pt);
\fill[slate] (1.777,6.303) circle (0.28pt);
\fill[slate] (1.777,6.552) circle (0.28pt);
\fill[slate] (1.777,6.754) circle (0.28pt);
\fill[slate] (1.777,6.801) circle (0.28pt);
\fill[slate] (1.777,5.402) circle (0.28pt);
\fill[slate] (1.777,5.651) circle (0.28pt);
\fill[slate] (1.777,5.852) circle (0.28pt);
\fill[slate] (1.777,6.308) circle (0.28pt);
\fill[slate] (1.777,6.551) circle (0.28pt);
\fill[slate] (1.777,6.764) circle (0.28pt);
\fill[slate] (1.777,6.799) circle (0.28pt);
\fill[slate] (1.777,5.9) circle (0.28pt);
\fill[slate] (1.777,6.149) circle (0.28pt);
\fill[slate] (1.777,6.349) circle (0.28pt);
\fill[slate] (1.777,6.805) circle (0.28pt);
\fill[slate] (1.777,6.398) circle (0.28pt);
\fill[slate] (1.777,6.647) circle (0.28pt);
\fill[slate] (1.777,6.847) circle (0.28pt);
\fill[slate] (1.777,4.005) circle (0.28pt);
\fill[slate] (1.777,4.254) circle (0.28pt);
\fill[slate] (1.777,4.454) circle (0.28pt);
\fill[slate] (1.777,4.91) circle (0.28pt);
\fill[slate] (1.777,5.154) circle (0.28pt);
\fill[slate] (1.777,5.366) circle (0.28pt);
\fill[slate] (1.777,5.822) circle (0.28pt);
\fill[slate] (1.777,6.066) circle (0.28pt);
\fill[slate] (1.777,6.278) circle (0.28pt);
\fill[slate] (1.777,6.734) circle (0.28pt);
\fill[slate] (1.777,6.314) circle (0.28pt);
\fill[slate] (1.777,6.563) circle (0.28pt);
\fill[slate] (1.777,6.765) circle (0.28pt);
\fill[slate] (1.777,6.812) circle (0.28pt);
\fill[slate] (1.777,5.402) circle (0.28pt);
\fill[slate] (1.777,5.651) circle (0.28pt);
\fill[slate] (1.777,5.853) circle (0.28pt);
\fill[slate] (1.777,6.309) circle (0.28pt);
\fill[slate] (1.777,6.553) circle (0.28pt);
\fill[slate] (1.777,6.765) circle (0.28pt);
\fill[slate] (1.777,6.801) circle (0.28pt);
\fill[slate] (1.777,5.9) circle (0.28pt);
\fill[slate] (1.777,6.149) circle (0.28pt);
\fill[slate] (1.777,6.35) circle (0.28pt);
\fill[slate] (1.777,6.806) circle (0.28pt);
\fill[slate] (1.777,6.398) circle (0.28pt);
\fill[slate] (1.777,6.647) circle (0.28pt);
\fill[slate] (1.777,6.847) circle (0.28pt);
\fill[slate] (1.777,4.503) circle (0.28pt);
\fill[slate] (1.777,4.752) circle (0.28pt);
\fill[slate] (1.777,4.952) circle (0.28pt);
\fill[slate] (1.777,5.408) circle (0.28pt);
\fill[slate] (1.777,5.652) circle (0.28pt);
\fill[slate] (1.777,5.864) circle (0.28pt);
\fill[slate] (1.777,6.32) circle (0.28pt);
\fill[slate] (1.777,6.564) circle (0.28pt);
\fill[slate] (1.777,6.776) circle (0.28pt);
\fill[slate] (1.777,6.812) circle (0.28pt);
\fill[slate] (1.777,5.9) circle (0.28pt);
\fill[slate] (1.777,6.149) circle (0.28pt);
\fill[slate] (1.777,6.351) circle (0.28pt);
\fill[slate] (1.777,6.807) circle (0.28pt);
\fill[slate] (1.777,6.398) circle (0.28pt);
\fill[slate] (1.777,6.647) circle (0.28pt);
\fill[slate] (1.777,6.848) circle (0.28pt);
\fill[slate] (1.777,5.001) circle (0.28pt);
\fill[slate] (1.777,5.25) circle (0.28pt);
\fill[slate] (1.777,5.45) circle (0.28pt);
\fill[slate] (1.777,5.906) circle (0.28pt);
\fill[slate] (1.777,6.15) circle (0.28pt);
\fill[slate] (1.777,6.362) circle (0.28pt);
\fill[slate] (1.777,6.818) circle (0.28pt);
\fill[slate] (1.777,6.398) circle (0.28pt);
\fill[slate] (1.777,6.647) circle (0.28pt);
\fill[slate] (1.777,6.849) circle (0.28pt);
\fill[slate] (1.777,5.499) circle (0.28pt);
\fill[slate] (1.777,5.748) circle (0.28pt);
\fill[slate] (1.777,5.948) circle (0.28pt);
\fill[slate] (1.777,6.404) circle (0.28pt);
\fill[slate] (1.777,6.648) circle (0.28pt);
\fill[slate] (1.777,5.997) circle (0.28pt);
\fill[slate] (1.777,6.246) circle (0.28pt);
\fill[slate] (1.777,6.446) circle (0.28pt);
\fill[slate] (1.777,6.495) circle (0.28pt);
\fill[slate] (1.777,6.744) circle (0.28pt);
\fill[slate] (0.267,2.018) circle (0.28pt);
\fill[slate] (0.268,2.474) circle (0.28pt);
\fill[slate] (0.268,2.93) circle (0.28pt);
\fill[slate] (0.268,3.173) circle (0.28pt);
\fill[slate] (0.268,3.386) circle (0.28pt);
\fill[slate] (0.268,3.842) circle (0.28pt);
\fill[slate] (0.268,4.086) circle (0.28pt);
\fill[slate] (0.268,4.298) circle (0.28pt);
\fill[slate] (0.268,4.754) circle (0.28pt);
\fill[slate] (0.268,4.998) circle (0.28pt);
\fill[slate] (0.268,5.211) circle (0.28pt);
\fill[slate] (0.268,5.667) circle (0.28pt);
\fill[slate] (0.268,5.91) circle (0.28pt);
\fill[slate] (0.268,6.123) circle (0.28pt);
\fill[slate] (0.268,6.579) circle (0.28pt);
\fill[slate] (0.268,6.822) circle (0.28pt);
\fill[slate] (0.268,6.158) circle (0.28pt);
\fill[slate] (0.268,6.407) circle (0.28pt);
\fill[slate] (0.268,6.609) circle (0.28pt);
\fill[slate] (0.268,6.656) circle (0.28pt);
\fill[slate] (0.268,5.246) circle (0.28pt);
\fill[slate] (0.268,5.495) circle (0.28pt);
\fill[slate] (0.268,5.697) circle (0.28pt);
\fill[slate] (0.268,6.153) circle (0.28pt);
\fill[slate] (0.268,6.397) circle (0.28pt);
\fill[slate] (0.268,6.609) circle (0.28pt);
\fill[slate] (0.268,6.645) circle (0.28pt);
\fill[slate] (0.268,5.744) circle (0.28pt);
\fill[slate] (0.268,5.993) circle (0.28pt);
\fill[slate] (0.268,6.194) circle (0.28pt);
\fill[slate] (0.268,6.65) circle (0.28pt);
\fill[slate] (0.268,6.242) circle (0.28pt);
\fill[slate] (0.268,6.491) circle (0.28pt);
\fill[slate] (0.268,6.691) circle (0.28pt);
\fill[slate] (0.268,6.74) circle (0.28pt);
\fill[slate] (0.268,4.334) circle (0.28pt);
\fill[slate] (0.268,4.583) circle (0.28pt);
\fill[slate] (0.268,4.785) circle (0.28pt);
\fill[slate] (0.268,5.241) circle (0.28pt);
\fill[slate] (0.268,5.485) circle (0.28pt);
\fill[slate] (0.268,5.697) circle (0.28pt);
\fill[slate] (0.268,6.153) circle (0.28pt);
\fill[slate] (0.268,6.396) circle (0.28pt);
\fill[slate] (0.268,6.609) circle (0.28pt);
\fill[slate] (0.268,6.645) circle (0.28pt);
\fill[slate] (0.268,5.733) circle (0.28pt);
\fill[slate] (0.268,5.982) circle (0.28pt);
\fill[slate] (0.268,6.184) circle (0.28pt);
\fill[slate] (0.268,6.64) circle (0.28pt);
\fill[slate] (0.268,6.231) circle (0.28pt);
\fill[slate] (0.268,6.48) circle (0.28pt);
\fill[slate] (0.268,6.68) circle (0.28pt);
\fill[slate] (0.268,6.729) circle (0.28pt);
\fill[slate] (0.268,4.832) circle (0.28pt);
\fill[slate] (0.268,5.081) circle (0.28pt);
\fill[slate] (0.268,5.281) circle (0.28pt);
\fill[slate] (0.268,5.737) circle (0.28pt);
\fill[slate] (0.268,5.981) circle (0.28pt);
\fill[slate] (0.268,6.193) circle (0.28pt);
\fill[slate] (0.268,6.65) circle (0.28pt);
\fill[slate] (0.268,6.229) circle (0.28pt);
\fill[slate] (0.268,6.478) circle (0.28pt);
\fill[slate] (0.268,6.68) circle (0.28pt);
\fill[slate] (0.268,6.727) circle (0.28pt);
\fill[slate] (0.268,5.33) circle (0.28pt);
\fill[slate] (0.268,5.579) circle (0.28pt);
\fill[slate] (0.268,5.779) circle (0.28pt);
\fill[slate] (0.268,6.235) circle (0.28pt);
\fill[slate] (0.268,6.479) circle (0.28pt);
\fill[slate] (0.268,6.691) circle (0.28pt);
\fill[slate] (0.268,6.727) circle (0.28pt);
\fill[slate] (0.268,5.828) circle (0.28pt);
\fill[slate] (0.268,6.077) circle (0.28pt);
\fill[slate] (0.268,6.277) circle (0.28pt);
\fill[slate] (0.268,6.733) circle (0.28pt);
\fill[slate] (0.268,6.326) circle (0.28pt);
\fill[slate] (0.268,6.575) circle (0.28pt);
\fill[slate] (0.268,6.775) circle (0.28pt);
\fill[slate] (0.268,6.824) circle (0.28pt);
\fill[slate] (0.267,3.422) circle (0.28pt);
\fill[slate] (0.267,3.671) circle (0.28pt);
\fill[slate] (0.268,3.873) circle (0.28pt);
\fill[slate] (0.268,4.329) circle (0.28pt);
\fill[slate] (0.268,4.572) circle (0.28pt);
\fill[slate] (0.268,4.785) circle (0.28pt);
\fill[slate] (0.268,5.241) circle (0.28pt);
\fill[slate] (0.268,5.484) circle (0.28pt);
\fill[slate] (0.268,5.697) circle (0.28pt);
\fill[slate] (0.268,6.153) circle (0.28pt);
\fill[slate] (0.268,6.396) circle (0.28pt);
\fill[slate] (0.268,6.609) circle (0.28pt);
\fill[slate] (0.268,6.645) circle (0.28pt);
\fill[slate] (0.268,5.732) circle (0.28pt);
\fill[slate] (0.268,5.981) circle (0.28pt);
\fill[slate] (0.268,6.184) circle (0.28pt);
\fill[slate] (0.268,6.64) circle (0.28pt);
\fill[slate] (0.268,6.23) circle (0.28pt);
\fill[slate] (0.268,6.479) circle (0.28pt);
\fill[slate] (0.268,6.68) circle (0.28pt);
\fill[slate] (0.268,6.728) circle (0.28pt);
\fill[slate] (0.268,4.821) circle (0.28pt);
\fill[slate] (0.268,5.07) circle (0.28pt);
\fill[slate] (0.268,5.272) circle (0.28pt);
\fill[slate] (0.268,5.728) circle (0.28pt);
\fill[slate] (0.268,5.971) circle (0.28pt);
\fill[slate] (0.268,6.184) circle (0.28pt);
\fill[slate] (0.268,6.64) circle (0.28pt);
\fill[slate] (0.268,6.22) circle (0.28pt);
\fill[slate] (0.268,6.469) circle (0.28pt);
\fill[slate] (0.268,6.671) circle (0.28pt);
\fill[slate] (0.268,6.718) circle (0.28pt);
\fill[slate] (0.268,5.319) circle (0.28pt);
\fill[slate] (0.268,5.568) circle (0.28pt);
\fill[slate] (0.268,5.768) circle (0.28pt);
\fill[slate] (0.268,6.224) circle (0.28pt);
\fill[slate] (0.268,6.468) circle (0.28pt);
\fill[slate] (0.268,6.68) circle (0.28pt);
\fill[slate] (0.268,6.716) circle (0.28pt);
\fill[slate] (0.268,5.817) circle (0.28pt);
\fill[slate] (0.268,6.066) circle (0.28pt);
\fill[slate] (0.268,6.266) circle (0.28pt);
\fill[slate] (0.268,6.722) circle (0.28pt);
\fill[slate] (0.268,6.315) circle (0.28pt);
\fill[slate] (0.268,6.564) circle (0.28pt);
\fill[slate] (0.268,6.764) circle (0.28pt);
\fill[slate] (0.268,6.813) circle (0.28pt);
\fill[slate] (0.267,3.92) circle (0.28pt);
\fill[slate] (0.267,4.169) circle (0.28pt);
\fill[slate] (0.267,4.369) circle (0.28pt);
\fill[slate] (0.267,4.825) circle (0.28pt);
\fill[slate] (0.267,5.069) circle (0.28pt);
\fill[slate] (0.267,5.281) circle (0.28pt);
\fill[slate] (0.267,5.737) circle (0.28pt);
\fill[slate] (0.267,5.981) circle (0.28pt);
\fill[slate] (0.267,6.193) circle (0.28pt);
\fill[slate] (0.267,6.65) circle (0.28pt);
\fill[slate] (0.267,6.229) circle (0.28pt);
\fill[slate] (0.267,6.478) circle (0.28pt);
\fill[slate] (0.267,6.68) circle (0.28pt);
\fill[slate] (0.267,6.727) circle (0.28pt);
\fill[slate] (0.267,5.317) circle (0.28pt);
\fill[slate] (0.267,5.566) circle (0.28pt);
\fill[slate] (0.267,5.768) circle (0.28pt);
\fill[slate] (0.267,6.224) circle (0.28pt);
\fill[slate] (0.267,6.468) circle (0.28pt);
\fill[slate] (0.267,6.68) circle (0.28pt);
\fill[slate] (0.267,6.716) circle (0.28pt);
\fill[slate] (0.267,5.815) circle (0.28pt);
\fill[slate] (0.267,6.064) circle (0.28pt);
\fill[slate] (0.267,6.265) circle (0.28pt);
\fill[slate] (0.267,6.721) circle (0.28pt);
\fill[slate] (0.267,6.313) circle (0.28pt);
\fill[slate] (0.267,6.562) circle (0.28pt);
\fill[slate] (0.267,6.763) circle (0.28pt);
\fill[slate] (0.267,6.811) circle (0.28pt);
\fill[slate] (0.267,4.418) circle (0.28pt);
\fill[slate] (0.267,4.667) circle (0.28pt);
\fill[slate] (0.267,4.867) circle (0.28pt);
\fill[slate] (0.267,5.323) circle (0.28pt);
\fill[slate] (0.267,5.567) circle (0.28pt);
\fill[slate] (0.267,5.779) circle (0.28pt);
\fill[slate] (0.267,6.235) circle (0.28pt);
\fill[slate] (0.267,6.478) circle (0.28pt);
\fill[slate] (0.267,6.691) circle (0.28pt);
\fill[slate] (0.267,6.727) circle (0.28pt);
\fill[slate] (0.267,5.815) circle (0.28pt);
\fill[slate] (0.267,6.064) circle (0.28pt);
\fill[slate] (0.267,6.266) circle (0.28pt);
\fill[slate] (0.267,6.722) circle (0.28pt);
\fill[slate] (0.267,6.313) circle (0.28pt);
\fill[slate] (0.267,6.562) circle (0.28pt);
\fill[slate] (0.267,6.763) circle (0.28pt);
\fill[slate] (0.267,6.811) circle (0.28pt);
\fill[slate] (0.267,4.916) circle (0.28pt);
\fill[slate] (0.267,5.165) circle (0.28pt);
\fill[slate] (0.267,5.365) circle (0.28pt);
\fill[slate] (0.267,5.821) circle (0.28pt);
\fill[slate] (0.267,6.065) circle (0.28pt);
\fill[slate] (0.267,6.277) circle (0.28pt);
\fill[slate] (0.267,6.733) circle (0.28pt);
\fill[slate] (0.267,6.313) circle (0.28pt);
\fill[slate] (0.267,6.562) circle (0.28pt);
\fill[slate] (0.267,6.764) circle (0.28pt);
\fill[slate] (0.267,6.811) circle (0.28pt);
\fill[slate] (0.267,5.414) circle (0.28pt);
\fill[slate] (0.267,5.663) circle (0.28pt);
\fill[slate] (0.267,5.863) circle (0.28pt);
\fill[slate] (0.267,6.319) circle (0.28pt);
\fill[slate] (0.267,6.563) circle (0.28pt);
\fill[slate] (0.267,6.775) circle (0.28pt);
\fill[slate] (0.267,6.811) circle (0.28pt);
\fill[slate] (0.267,5.912) circle (0.28pt);
\fill[slate] (0.267,6.161) circle (0.28pt);
\fill[slate] (0.267,6.361) circle (0.28pt);
\fill[slate] (0.267,6.817) circle (0.28pt);
\fill[slate] (0.267,6.41) circle (0.28pt);
\fill[slate] (0.267,6.659) circle (0.28pt);
\fill[slate] (0.261,2.961) circle (0.28pt);
\fill[slate] (0.261,3.417) circle (0.28pt);
\fill[slate] (0.261,3.661) circle (0.28pt);
\fill[slate] (0.261,3.873) circle (0.28pt);
\fill[slate] (0.261,4.329) circle (0.28pt);
\fill[slate] (0.261,4.573) circle (0.28pt);
\fill[slate] (0.261,4.785) circle (0.28pt);
\fill[slate] (0.261,5.241) circle (0.28pt);
\fill[slate] (0.261,5.485) circle (0.28pt);
\fill[slate] (0.261,5.697) circle (0.28pt);
\fill[slate] (0.261,6.154) circle (0.28pt);
\fill[slate] (0.261,6.397) circle (0.28pt);
\fill[slate] (0.261,6.61) circle (0.28pt);
\fill[slate] (0.261,6.645) circle (0.28pt);
\fill[slate] (0.261,5.733) circle (0.28pt);
\fill[slate] (0.261,5.982) circle (0.28pt);
\fill[slate] (0.261,6.184) circle (0.28pt);
\fill[slate] (0.261,6.64) circle (0.28pt);
\fill[slate] (0.261,6.231) circle (0.28pt);
\fill[slate] (0.261,6.48) circle (0.28pt);
\fill[slate] (0.261,6.681) circle (0.28pt);
\fill[slate] (0.261,6.729) circle (0.28pt);
\fill[slate] (0.261,4.821) circle (0.28pt);
\fill[slate] (0.261,5.07) circle (0.28pt);
\fill[slate] (0.261,5.272) circle (0.28pt);
\fill[slate] (0.261,5.728) circle (0.28pt);
\fill[slate] (0.261,5.971) circle (0.28pt);
\fill[slate] (0.261,6.184) circle (0.28pt);
\fill[slate] (0.261,6.64) circle (0.28pt);
\fill[slate] (0.261,6.22) circle (0.28pt);
\fill[slate] (0.261,6.469) circle (0.28pt);
\fill[slate] (0.261,6.671) circle (0.28pt);
\fill[slate] (0.261,6.718) circle (0.28pt);
\fill[slate] (0.261,5.319) circle (0.28pt);
\fill[slate] (0.261,5.568) circle (0.28pt);
\fill[slate] (0.261,5.768) circle (0.28pt);
\fill[slate] (0.261,6.224) circle (0.28pt);
\fill[slate] (0.261,6.468) circle (0.28pt);
\fill[slate] (0.261,6.68) circle (0.28pt);
\fill[slate] (0.261,6.716) circle (0.28pt);
\fill[slate] (0.261,5.817) circle (0.28pt);
\fill[slate] (0.261,6.066) circle (0.28pt);
\fill[slate] (0.261,6.266) circle (0.28pt);
\fill[slate] (0.261,6.722) circle (0.28pt);
\fill[slate] (0.261,6.315) circle (0.28pt);
\fill[slate] (0.261,6.564) circle (0.28pt);
\fill[slate] (0.261,6.764) circle (0.28pt);
\fill[slate] (0.261,6.813) circle (0.28pt);
\fill[slate] (0.261,3.909) circle (0.28pt);
\fill[slate] (0.261,4.36) circle (0.28pt);
\fill[slate] (0.261,4.816) circle (0.28pt);
\fill[slate] (0.261,5.06) circle (0.28pt);
\fill[slate] (0.261,5.272) circle (0.28pt);
\fill[slate] (0.261,5.728) circle (0.28pt);
\fill[slate] (0.261,5.971) circle (0.28pt);
\fill[slate] (0.261,6.184) circle (0.28pt);
\fill[slate] (0.261,6.64) circle (0.28pt);
\fill[slate] (0.261,6.22) circle (0.28pt);
\fill[slate] (0.261,6.469) circle (0.28pt);
\fill[slate] (0.261,6.671) circle (0.28pt);
\fill[slate] (0.261,6.718) circle (0.28pt);
\fill[slate] (0.261,5.308) circle (0.28pt);
\fill[slate] (0.261,5.557) circle (0.28pt);
\fill[slate] (0.261,5.759) circle (0.28pt);
\fill[slate] (0.261,6.215) circle (0.28pt);
\fill[slate] (0.261,6.459) circle (0.28pt);
\fill[slate] (0.261,6.671) circle (0.28pt);
\fill[slate] (0.261,6.707) circle (0.28pt);
\fill[slate] (0.261,5.806) circle (0.28pt);
\fill[slate] (0.261,6.055) circle (0.28pt);
\fill[slate] (0.261,6.256) circle (0.28pt);
\fill[slate] (0.261,6.711) circle (0.28pt);
\fill[slate] (0.261,6.304) circle (0.28pt);
\fill[slate] (0.261,6.553) circle (0.28pt);
\fill[slate] (0.261,6.753) circle (0.28pt);
\fill[slate] (0.261,6.802) circle (0.28pt);
\fill[slate] (0.261,4.407) circle (0.28pt);
\fill[slate] (0.261,4.656) circle (0.28pt);
\fill[slate] (0.261,4.857) circle (0.28pt);
\fill[slate] (0.261,5.312) circle (0.28pt);
\fill[slate] (0.261,5.556) circle (0.28pt);
\fill[slate] (0.261,5.768) circle (0.28pt);
\fill[slate] (0.261,6.225) circle (0.28pt);
\fill[slate] (0.261,6.468) circle (0.28pt);
\fill[slate] (0.261,6.681) circle (0.28pt);
\fill[slate] (0.261,6.716) circle (0.28pt);
\fill[slate] (0.261,5.804) circle (0.28pt);
\fill[slate] (0.261,6.256) circle (0.28pt);
\fill[slate] (0.261,6.711) circle (0.28pt);
\fill[slate] (0.261,6.302) circle (0.28pt);
\fill[slate] (0.261,6.551) circle (0.28pt);
\fill[slate] (0.261,6.752) circle (0.28pt);
\fill[slate] (0.261,6.8) circle (0.28pt);
\fill[slate] (0.261,4.905) circle (0.28pt);
\fill[slate] (0.261,5.154) circle (0.28pt);
\fill[slate] (0.261,5.354) circle (0.28pt);
\fill[slate] (0.261,5.81) circle (0.28pt);
\fill[slate] (0.261,6.266) circle (0.28pt);
\fill[slate] (0.261,6.722) circle (0.28pt);
\fill[slate] (0.261,6.302) circle (0.28pt);
\fill[slate] (0.261,6.551) circle (0.28pt);
\fill[slate] (0.261,6.753) circle (0.28pt);
\fill[slate] (0.261,6.8) circle (0.28pt);
\fill[slate] (0.261,5.403) circle (0.28pt);
\fill[slate] (0.261,5.652) circle (0.28pt);
\fill[slate] (0.261,5.852) circle (0.28pt);
\fill[slate] (0.261,6.308) circle (0.28pt);
\fill[slate] (0.261,6.552) circle (0.28pt);
\fill[slate] (0.261,6.764) circle (0.28pt);
\fill[slate] (0.261,6.8) circle (0.28pt);
\fill[slate] (0.261,5.901) circle (0.28pt);
\fill[slate] (0.261,6.15) circle (0.28pt);
\fill[slate] (0.261,6.35) circle (0.28pt);
\fill[slate] (0.261,6.806) circle (0.28pt);
\fill[slate] (0.261,6.399) circle (0.28pt);
\fill[slate] (0.261,6.648) circle (0.28pt);
\fill[slate] (0.261,6.848) circle (0.28pt);
\fill[slate] (0.26,3.257) circle (0.28pt);
\fill[slate] (0.26,3.913) circle (0.28pt);
\fill[slate] (0.26,4.37) circle (0.28pt);
\fill[slate] (0.26,4.826) circle (0.28pt);
\fill[slate] (0.26,5.069) circle (0.28pt);
\fill[slate] (0.26,5.282) circle (0.28pt);
\fill[slate] (0.26,5.738) circle (0.28pt);
\fill[slate] (0.26,5.981) circle (0.28pt);
\fill[slate] (0.26,6.194) circle (0.28pt);
\fill[slate] (0.26,6.65) circle (0.28pt);
\fill[slate] (0.26,6.229) circle (0.28pt);
\fill[slate] (0.26,6.478) circle (0.28pt);
\fill[slate] (0.26,6.681) circle (0.28pt);
\fill[slate] (0.26,6.727) circle (0.28pt);
\fill[slate] (0.26,5.317) circle (0.28pt);
\fill[slate] (0.26,5.566) circle (0.28pt);
\fill[slate] (0.26,5.768) circle (0.28pt);
\fill[slate] (0.26,6.224) circle (0.28pt);
\fill[slate] (0.26,6.468) circle (0.28pt);
\fill[slate] (0.26,6.68) circle (0.28pt);
\fill[slate] (0.26,6.716) circle (0.28pt);
\fill[slate] (0.26,5.815) circle (0.28pt);
\fill[slate] (0.26,6.064) circle (0.28pt);
\fill[slate] (0.26,6.265) circle (0.28pt);
\fill[slate] (0.26,6.721) circle (0.28pt);
\fill[slate] (0.26,6.313) circle (0.28pt);
\fill[slate] (0.26,6.562) circle (0.28pt);
\fill[slate] (0.26,6.763) circle (0.28pt);
\fill[slate] (0.26,6.811) circle (0.28pt);
\fill[slate] (0.26,4.857) circle (0.28pt);
\fill[slate] (0.26,5.312) circle (0.28pt);
\fill[slate] (0.26,5.556) circle (0.28pt);
\fill[slate] (0.26,5.769) circle (0.28pt);
\fill[slate] (0.26,6.225) circle (0.28pt);
\fill[slate] (0.26,6.468) circle (0.28pt);
\fill[slate] (0.26,6.681) circle (0.28pt);
\fill[slate] (0.26,6.716) circle (0.28pt);
\fill[slate] (0.26,5.804) circle (0.28pt);
\fill[slate] (0.26,6.256) circle (0.28pt);
\fill[slate] (0.26,6.711) circle (0.28pt);
\fill[slate] (0.26,6.302) circle (0.28pt);
\fill[slate] (0.26,6.551) circle (0.28pt);
\fill[slate] (0.26,6.752) circle (0.28pt);
\fill[slate] (0.26,6.8) circle (0.28pt);
\fill[slate] (0.26,4.903) circle (0.28pt);
\fill[slate] (0.26,5.152) circle (0.28pt);
\fill[slate] (0.26,5.809) circle (0.28pt);
\fill[slate] (0.26,6.265) circle (0.28pt);
\fill[slate] (0.26,6.721) circle (0.28pt);
\fill[slate] (0.26,6.752) circle (0.28pt);
\fill[slate] (0.26,6.799) circle (0.28pt);
\fill[slate] (0.26,5.401) circle (0.28pt);
\fill[slate] (0.26,5.65) circle (0.28pt);
\fill[slate] (0.26,6.307) circle (0.28pt);
\fill[slate] (0.26,6.55) circle (0.28pt);
\fill[slate] (0.26,6.763) circle (0.28pt);
\fill[slate] (0.26,5.899) circle (0.28pt);
\fill[slate] (0.26,6.148) circle (0.28pt);
\fill[slate] (0.26,6.349) circle (0.28pt);
\fill[slate] (0.26,6.805) circle (0.28pt);
\fill[slate] (0.26,6.397) circle (0.28pt);
\fill[slate] (0.26,6.646) circle (0.28pt);
\fill[slate] (0.26,6.847) circle (0.28pt);
\fill[slate] (0.26,3.506) circle (0.28pt);
\fill[slate] (0.26,3.755) circle (0.28pt);
\fill[slate] (0.26,4.411) circle (0.28pt);
\fill[slate] (0.26,4.655) circle (0.28pt);
\fill[slate] (0.26,4.867) circle (0.28pt);
\fill[slate] (0.26,5.323) circle (0.28pt);
\fill[slate] (0.26,5.567) circle (0.28pt);
\fill[slate] (0.26,5.779) circle (0.28pt);
\fill[slate] (0.26,6.236) circle (0.28pt);
\fill[slate] (0.26,6.479) circle (0.28pt);
\fill[slate] (0.26,6.692) circle (0.28pt);
\fill[slate] (0.26,6.727) circle (0.28pt);
\fill[slate] (0.26,5.815) circle (0.28pt);
\fill[slate] (0.26,6.064) circle (0.28pt);
\fill[slate] (0.26,6.266) circle (0.28pt);
\fill[slate] (0.26,6.722) circle (0.28pt);
\fill[slate] (0.26,6.313) circle (0.28pt);
\fill[slate] (0.26,6.562) circle (0.28pt);
\fill[slate] (0.26,6.763) circle (0.28pt);
\fill[slate] (0.26,6.811) circle (0.28pt);
\fill[slate] (0.26,5.354) circle (0.28pt);
\fill[slate] (0.26,5.81) circle (0.28pt);
\fill[slate] (0.26,6.054) circle (0.28pt);
\fill[slate] (0.26,6.266) circle (0.28pt);
\fill[slate] (0.26,6.722) circle (0.28pt);
\fill[slate] (0.26,6.302) circle (0.28pt);
\fill[slate] (0.26,6.551) circle (0.28pt);
\fill[slate] (0.26,6.753) circle (0.28pt);
\fill[slate] (0.26,6.8) circle (0.28pt);
\fill[slate] (0.26,5.65) circle (0.28pt);
\fill[slate] (0.26,6.307) circle (0.28pt);
\fill[slate] (0.26,6.55) circle (0.28pt);
\fill[slate] (0.26,6.763) circle (0.28pt);
\fill[slate] (0.26,6.798) circle (0.28pt);
\fill[slate] (0.26,5.899) circle (0.28pt);
\fill[slate] (0.26,6.148) circle (0.28pt);
\fill[slate] (0.26,6.804) circle (0.28pt);
\fill[slate] (0.26,6.397) circle (0.28pt);
\fill[slate] (0.26,6.646) circle (0.28pt);
\fill[slate] (0.26,6.847) circle (0.28pt);
\fill[slate] (0.26,4.004) circle (0.28pt);
\fill[slate] (0.26,4.253) circle (0.28pt);
\fill[slate] (0.26,4.453) circle (0.28pt);
\fill[slate] (0.26,4.909) circle (0.28pt);
\fill[slate] (0.26,5.153) circle (0.28pt);
\fill[slate] (0.26,5.365) circle (0.28pt);
\fill[slate] (0.26,5.821) circle (0.28pt);
\fill[slate] (0.26,6.065) circle (0.28pt);
\fill[slate] (0.26,6.277) circle (0.28pt);
\fill[slate] (0.26,6.733) circle (0.28pt);
\fill[slate] (0.26,6.313) circle (0.28pt);
\fill[slate] (0.26,6.562) circle (0.28pt);
\fill[slate] (0.26,6.764) circle (0.28pt);
\fill[slate] (0.26,6.811) circle (0.28pt);
\fill[slate] (0.26,5.65) circle (0.28pt);
\fill[slate] (0.26,5.852) circle (0.28pt);
\fill[slate] (0.26,6.308) circle (0.28pt);
\fill[slate] (0.26,6.552) circle (0.28pt);
\fill[slate] (0.26,6.764) circle (0.28pt);
\fill[slate] (0.26,6.8) circle (0.28pt);
\fill[slate] (0.26,5.899) circle (0.28pt);
\fill[slate] (0.26,6.148) circle (0.28pt);
\fill[slate] (0.26,6.349) circle (0.28pt);
\fill[slate] (0.26,6.805) circle (0.28pt);
\fill[slate] (0.26,6.397) circle (0.28pt);
\fill[slate] (0.26,6.646) circle (0.28pt);
\fill[slate] (0.26,6.847) circle (0.28pt);
\fill[slate] (0.26,4.502) circle (0.28pt);
\fill[slate] (0.26,4.751) circle (0.28pt);
\fill[slate] (0.26,4.951) circle (0.28pt);
\fill[slate] (0.26,5.407) circle (0.28pt);
\fill[slate] (0.26,5.651) circle (0.28pt);
\fill[slate] (0.26,5.863) circle (0.28pt);
\fill[slate] (0.26,6.319) circle (0.28pt);
\fill[slate] (0.26,6.563) circle (0.28pt);
\fill[slate] (0.26,6.775) circle (0.28pt);
\fill[slate] (0.26,6.811) circle (0.28pt);
\fill[slate] (0.26,5.899) circle (0.28pt);
\fill[slate] (0.26,6.148) circle (0.28pt);
\fill[slate] (0.26,6.35) circle (0.28pt);
\fill[slate] (0.26,6.806) circle (0.28pt);
\fill[slate] (0.26,6.397) circle (0.28pt);
\fill[slate] (0.26,6.646) circle (0.28pt);
\fill[slate] (0.26,6.847) circle (0.28pt);
\fill[slate] (0.26,5) circle (0.28pt);
\fill[slate] (0.26,5.249) circle (0.28pt);
\fill[slate] (0.26,5.449) circle (0.28pt);
\fill[slate] (0.26,5.905) circle (0.28pt);
\fill[slate] (0.26,6.149) circle (0.28pt);
\fill[slate] (0.26,6.361) circle (0.28pt);
\fill[slate] (0.26,6.817) circle (0.28pt);
\fill[slate] (0.26,6.397) circle (0.28pt);
\fill[slate] (0.26,6.646) circle (0.28pt);
\fill[slate] (0.26,6.848) circle (0.28pt);
\fill[slate] (0.26,5.498) circle (0.28pt);
\fill[slate] (0.26,5.747) circle (0.28pt);
\fill[slate] (0.26,5.947) circle (0.28pt);
\fill[slate] (0.26,6.403) circle (0.28pt);
\fill[slate] (0.26,6.647) circle (0.28pt);
\fill[slate] (0.26,5.996) circle (0.28pt);
\fill[slate] (0.26,6.245) circle (0.28pt);
\fill[slate] (0.26,6.445) circle (0.28pt);
\fill[slate] (0.26,6.494) circle (0.28pt);
\fill[slate] (0.26,6.743) circle (0.28pt);
\fill[slate] (0.039,2.516) circle (0.28pt);
\fill[slate] (0.038,2.76) circle (0.28pt);
\fill[slate] (0.039,2.972) circle (0.28pt);
\fill[slate] (0.039,3.428) circle (0.28pt);
\fill[slate] (0.039,3.671) circle (0.28pt);
\fill[slate] (0.039,3.884) circle (0.28pt);
\fill[slate] (0.039,4.34) circle (0.28pt);
\fill[slate] (0.039,4.583) circle (0.28pt);
\fill[slate] (0.039,4.796) circle (0.28pt);
\fill[slate] (0.039,5.252) circle (0.28pt);
\fill[slate] (0.039,5.496) circle (0.28pt);
\fill[slate] (0.039,5.708) circle (0.28pt);
\fill[slate] (0.039,6.165) circle (0.28pt);
\fill[slate] (0.039,6.408) circle (0.28pt);
\fill[slate] (0.039,6.621) circle (0.28pt);
\fill[slate] (0.039,6.656) circle (0.28pt);
\fill[slate] (0.039,5.744) circle (0.28pt);
\fill[slate] (0.039,5.993) circle (0.28pt);
\fill[slate] (0.039,6.195) circle (0.28pt);
\fill[slate] (0.039,6.651) circle (0.28pt);
\fill[slate] (0.039,6.242) circle (0.28pt);
\fill[slate] (0.039,6.491) circle (0.28pt);
\fill[slate] (0.039,6.691) circle (0.28pt);
\fill[slate] (0.039,6.74) circle (0.28pt);
\fill[slate] (0.039,4.832) circle (0.28pt);
\fill[slate] (0.039,5.081) circle (0.28pt);
\fill[slate] (0.039,5.283) circle (0.28pt);
\fill[slate] (0.039,5.739) circle (0.28pt);
\fill[slate] (0.039,5.982) circle (0.28pt);
\fill[slate] (0.039,6.195) circle (0.28pt);
\fill[slate] (0.039,6.651) circle (0.28pt);
\fill[slate] (0.039,6.231) circle (0.28pt);
\fill[slate] (0.039,6.48) circle (0.28pt);
\fill[slate] (0.039,6.682) circle (0.28pt);
\fill[slate] (0.039,6.729) circle (0.28pt);
\fill[slate] (0.039,5.33) circle (0.28pt);
\fill[slate] (0.039,5.579) circle (0.28pt);
\fill[slate] (0.039,5.779) circle (0.28pt);
\fill[slate] (0.039,6.235) circle (0.28pt);
\fill[slate] (0.039,6.479) circle (0.28pt);
\fill[slate] (0.039,6.691) circle (0.28pt);
\fill[slate] (0.039,6.727) circle (0.28pt);
\fill[slate] (0.039,5.828) circle (0.28pt);
\fill[slate] (0.039,6.077) circle (0.28pt);
\fill[slate] (0.039,6.277) circle (0.28pt);
\fill[slate] (0.039,6.733) circle (0.28pt);
\fill[slate] (0.039,6.326) circle (0.28pt);
\fill[slate] (0.039,6.575) circle (0.28pt);
\fill[slate] (0.039,6.775) circle (0.28pt);
\fill[slate] (0.039,6.824) circle (0.28pt);
\fill[slate] (0.039,3.92) circle (0.28pt);
\fill[slate] (0.039,4.168) circle (0.28pt);
\fill[slate] (0.039,4.371) circle (0.28pt);
\fill[slate] (0.039,4.827) circle (0.28pt);
\fill[slate] (0.039,5.07) circle (0.28pt);
\fill[slate] (0.039,5.283) circle (0.28pt);
\fill[slate] (0.039,5.739) circle (0.28pt);
\fill[slate] (0.039,5.982) circle (0.28pt);
\fill[slate] (0.039,6.195) circle (0.28pt);
\fill[slate] (0.039,6.651) circle (0.28pt);
\fill[slate] (0.039,6.23) circle (0.28pt);
\fill[slate] (0.039,6.479) circle (0.28pt);
\fill[slate] (0.039,6.682) circle (0.28pt);
\fill[slate] (0.039,6.728) circle (0.28pt);
\fill[slate] (0.039,5.319) circle (0.28pt);
\fill[slate] (0.039,5.567) circle (0.28pt);
\fill[slate] (0.039,5.77) circle (0.28pt);
\fill[slate] (0.039,6.226) circle (0.28pt);
\fill[slate] (0.039,6.469) circle (0.28pt);
\fill[slate] (0.039,6.682) circle (0.28pt);
\fill[slate] (0.039,6.718) circle (0.28pt);
\fill[slate] (0.039,5.816) circle (0.28pt);
\fill[slate] (0.039,6.065) circle (0.28pt);
\fill[slate] (0.039,6.266) circle (0.28pt);
\fill[slate] (0.039,6.722) circle (0.28pt);
\fill[slate] (0.039,6.314) circle (0.28pt);
\fill[slate] (0.039,6.563) circle (0.28pt);
\fill[slate] (0.039,6.764) circle (0.28pt);
\fill[slate] (0.039,6.812) circle (0.28pt);
\fill[slate] (0.039,4.417) circle (0.28pt);
\fill[slate] (0.039,4.666) circle (0.28pt);
\fill[slate] (0.039,4.867) circle (0.28pt);
\fill[slate] (0.039,5.323) circle (0.28pt);
\fill[slate] (0.039,5.567) circle (0.28pt);
\fill[slate] (0.039,5.779) circle (0.28pt);
\fill[slate] (0.039,6.235) circle (0.28pt);
\fill[slate] (0.039,6.479) circle (0.28pt);
\fill[slate] (0.039,6.691) circle (0.28pt);
\fill[slate] (0.039,6.727) circle (0.28pt);
\fill[slate] (0.039,5.815) circle (0.28pt);
\fill[slate] (0.039,6.064) circle (0.28pt);
\fill[slate] (0.039,6.266) circle (0.28pt);
\fill[slate] (0.039,6.722) circle (0.28pt);
\fill[slate] (0.039,6.313) circle (0.28pt);
\fill[slate] (0.039,6.562) circle (0.28pt);
\fill[slate] (0.039,6.763) circle (0.28pt);
\fill[slate] (0.039,6.811) circle (0.28pt);
\fill[slate] (0.039,4.915) circle (0.28pt);
\fill[slate] (0.039,5.164) circle (0.28pt);
\fill[slate] (0.039,5.365) circle (0.28pt);
\fill[slate] (0.039,5.821) circle (0.28pt);
\fill[slate] (0.039,6.064) circle (0.28pt);
\fill[slate] (0.039,6.277) circle (0.28pt);
\fill[slate] (0.039,6.733) circle (0.28pt);
\fill[slate] (0.039,6.313) circle (0.28pt);
\fill[slate] (0.039,6.562) circle (0.28pt);
\fill[slate] (0.039,6.764) circle (0.28pt);
\fill[slate] (0.039,6.811) circle (0.28pt);
\fill[slate] (0.039,5.413) circle (0.28pt);
\fill[slate] (0.039,5.662) circle (0.28pt);
\fill[slate] (0.039,5.863) circle (0.28pt);
\fill[slate] (0.039,6.319) circle (0.28pt);
\fill[slate] (0.039,6.562) circle (0.28pt);
\fill[slate] (0.039,6.775) circle (0.28pt);
\fill[slate] (0.039,6.811) circle (0.28pt);
\fill[slate] (0.039,5.911) circle (0.28pt);
\fill[slate] (0.039,6.16) circle (0.28pt);
\fill[slate] (0.039,6.361) circle (0.28pt);
\fill[slate] (0.039,6.817) circle (0.28pt);
\fill[slate] (0.039,6.409) circle (0.28pt);
\fill[slate] (0.039,6.658) circle (0.28pt);
\fill[slate] (0.038,3.459) circle (0.28pt);
\fill[slate] (0.038,3.915) circle (0.28pt);
\fill[slate] (0.038,4.159) circle (0.28pt);
\fill[slate] (0.038,4.371) circle (0.28pt);
\fill[slate] (0.038,4.827) circle (0.28pt);
\fill[slate] (0.038,5.07) circle (0.28pt);
\fill[slate] (0.038,5.283) circle (0.28pt);
\fill[slate] (0.038,5.739) circle (0.28pt);
\fill[slate] (0.038,5.983) circle (0.28pt);
\fill[slate] (0.038,6.195) circle (0.28pt);
\fill[slate] (0.038,6.651) circle (0.28pt);
\fill[slate] (0.038,6.231) circle (0.28pt);
\fill[slate] (0.038,6.48) circle (0.28pt);
\fill[slate] (0.038,6.682) circle (0.28pt);
\fill[slate] (0.038,6.729) circle (0.28pt);
\fill[slate] (0.038,5.319) circle (0.28pt);
\fill[slate] (0.038,5.567) circle (0.28pt);
\fill[slate] (0.038,5.77) circle (0.28pt);
\fill[slate] (0.038,6.226) circle (0.28pt);
\fill[slate] (0.038,6.469) circle (0.28pt);
\fill[slate] (0.038,6.682) circle (0.28pt);
\fill[slate] (0.038,6.718) circle (0.28pt);
\fill[slate] (0.038,5.816) circle (0.28pt);
\fill[slate] (0.038,6.065) circle (0.28pt);
\fill[slate] (0.038,6.266) circle (0.28pt);
\fill[slate] (0.038,6.722) circle (0.28pt);
\fill[slate] (0.038,6.314) circle (0.28pt);
\fill[slate] (0.038,6.563) circle (0.28pt);
\fill[slate] (0.038,6.764) circle (0.28pt);
\fill[slate] (0.038,6.812) circle (0.28pt);
\fill[slate] (0.038,4.407) circle (0.28pt);
\fill[slate] (0.038,4.656) circle (0.28pt);
\fill[slate] (0.038,4.858) circle (0.28pt);
\fill[slate] (0.038,5.314) circle (0.28pt);
\fill[slate] (0.038,5.558) circle (0.28pt);
\fill[slate] (0.038,5.77) circle (0.28pt);
\fill[slate] (0.038,6.226) circle (0.28pt);
\fill[slate] (0.038,6.469) circle (0.28pt);
\fill[slate] (0.038,6.682) circle (0.28pt);
\fill[slate] (0.038,6.718) circle (0.28pt);
\fill[slate] (0.038,5.806) circle (0.28pt);
\fill[slate] (0.038,6.055) circle (0.28pt);
\fill[slate] (0.038,6.257) circle (0.28pt);
\fill[slate] (0.038,6.713) circle (0.28pt);
\fill[slate] (0.038,6.304) circle (0.28pt);
\fill[slate] (0.038,6.553) circle (0.28pt);
\fill[slate] (0.038,6.753) circle (0.28pt);
\fill[slate] (0.038,6.802) circle (0.28pt);
\fill[slate] (0.038,4.905) circle (0.28pt);
\fill[slate] (0.038,5.154) circle (0.28pt);
\fill[slate] (0.038,5.354) circle (0.28pt);
\fill[slate] (0.038,5.81) circle (0.28pt);
\fill[slate] (0.038,6.054) circle (0.28pt);
\fill[slate] (0.038,6.266) circle (0.28pt);
\fill[slate] (0.038,6.722) circle (0.28pt);
\fill[slate] (0.038,6.302) circle (0.28pt);
\fill[slate] (0.038,6.551) circle (0.28pt);
\fill[slate] (0.038,6.753) circle (0.28pt);
\fill[slate] (0.038,6.8) circle (0.28pt);
\fill[slate] (0.038,5.403) circle (0.28pt);
\fill[slate] (0.038,5.652) circle (0.28pt);
\fill[slate] (0.038,5.852) circle (0.28pt);
\fill[slate] (0.038,6.308) circle (0.28pt);
\fill[slate] (0.038,6.552) circle (0.28pt);
\fill[slate] (0.038,6.764) circle (0.28pt);
\fill[slate] (0.038,6.8) circle (0.28pt);
\fill[slate] (0.038,5.901) circle (0.28pt);
\fill[slate] (0.038,6.15) circle (0.28pt);
\fill[slate] (0.038,6.35) circle (0.28pt);
\fill[slate] (0.038,6.806) circle (0.28pt);
\fill[slate] (0.038,6.399) circle (0.28pt);
\fill[slate] (0.038,6.648) circle (0.28pt);
\fill[slate] (0.038,6.848) circle (0.28pt);
\fill[slate] (0.038,3.955) circle (0.28pt);
\fill[slate] (0.038,4.411) circle (0.28pt);
\fill[slate] (0.038,4.655) circle (0.28pt);
\fill[slate] (0.038,4.867) circle (0.28pt);
\fill[slate] (0.038,5.323) circle (0.28pt);
\fill[slate] (0.038,5.567) circle (0.28pt);
\fill[slate] (0.038,5.78) circle (0.28pt);
\fill[slate] (0.038,6.236) circle (0.28pt);
\fill[slate] (0.038,6.479) circle (0.28pt);
\fill[slate] (0.038,6.692) circle (0.28pt);
\fill[slate] (0.038,6.727) circle (0.28pt);
\fill[slate] (0.038,5.815) circle (0.28pt);
\fill[slate] (0.038,6.064) circle (0.28pt);
\fill[slate] (0.038,6.266) circle (0.28pt);
\fill[slate] (0.038,6.722) circle (0.28pt);
\fill[slate] (0.038,6.313) circle (0.28pt);
\fill[slate] (0.038,6.562) circle (0.28pt);
\fill[slate] (0.038,6.763) circle (0.28pt);
\fill[slate] (0.038,6.811) circle (0.28pt);
\fill[slate] (0.038,4.903) circle (0.28pt);
\fill[slate] (0.038,5.152) circle (0.28pt);
\fill[slate] (0.038,5.354) circle (0.28pt);
\fill[slate] (0.038,5.81) circle (0.28pt);
\fill[slate] (0.038,6.054) circle (0.28pt);
\fill[slate] (0.038,6.266) circle (0.28pt);
\fill[slate] (0.038,6.723) circle (0.28pt);
\fill[slate] (0.038,6.302) circle (0.28pt);
\fill[slate] (0.038,6.551) circle (0.28pt);
\fill[slate] (0.038,6.753) circle (0.28pt);
\fill[slate] (0.038,6.8) circle (0.28pt);
\fill[slate] (0.038,5.65) circle (0.28pt);
\fill[slate] (0.038,5.851) circle (0.28pt);
\fill[slate] (0.038,6.307) circle (0.28pt);
\fill[slate] (0.038,6.55) circle (0.28pt);
\fill[slate] (0.038,6.763) circle (0.28pt);
\fill[slate] (0.038,6.799) circle (0.28pt);
\fill[slate] (0.038,5.899) circle (0.28pt);
\fill[slate] (0.038,6.148) circle (0.28pt);
\fill[slate] (0.038,6.349) circle (0.28pt);
\fill[slate] (0.038,6.805) circle (0.28pt);
\fill[slate] (0.038,6.397) circle (0.28pt);
\fill[slate] (0.038,6.646) circle (0.28pt);
\fill[slate] (0.038,6.847) circle (0.28pt);
\fill[slate] (0.038,4.004) circle (0.28pt);
\fill[slate] (0.038,4.253) circle (0.28pt);
\fill[slate] (0.038,4.909) circle (0.28pt);
\fill[slate] (0.038,5.153) circle (0.28pt);
\fill[slate] (0.038,5.365) circle (0.28pt);
\fill[slate] (0.038,5.821) circle (0.28pt);
\fill[slate] (0.038,6.065) circle (0.28pt);
\fill[slate] (0.038,6.277) circle (0.28pt);
\fill[slate] (0.038,6.733) circle (0.28pt);
\fill[slate] (0.038,6.313) circle (0.28pt);
\fill[slate] (0.038,6.562) circle (0.28pt);
\fill[slate] (0.038,6.764) circle (0.28pt);
\fill[slate] (0.038,6.811) circle (0.28pt);
\fill[slate] (0.038,5.852) circle (0.28pt);
\fill[slate] (0.038,6.308) circle (0.28pt);
\fill[slate] (0.038,6.552) circle (0.28pt);
\fill[slate] (0.038,6.764) circle (0.28pt);
\fill[slate] (0.038,6.8) circle (0.28pt);
\fill[slate] (0.038,6.148) circle (0.28pt);
\fill[slate] (0.038,6.349) circle (0.28pt);
\fill[slate] (0.038,6.805) circle (0.28pt);
\fill[slate] (0.038,6.397) circle (0.28pt);
\fill[slate] (0.038,6.646) circle (0.28pt);
\fill[slate] (0.038,4.502) circle (0.28pt);
\fill[slate] (0.038,4.751) circle (0.28pt);
\fill[slate] (0.038,5.407) circle (0.28pt);
\fill[slate] (0.038,5.651) circle (0.28pt);
\fill[slate] (0.038,5.863) circle (0.28pt);
\fill[slate] (0.038,6.319) circle (0.28pt);
\fill[slate] (0.038,6.563) circle (0.28pt);
\fill[slate] (0.038,6.775) circle (0.28pt);
\fill[slate] (0.038,6.811) circle (0.28pt);
\fill[slate] (0.038,5.899) circle (0.28pt);
\fill[slate] (0.038,6.35) circle (0.28pt);
\fill[slate] (0.038,6.806) circle (0.28pt);
\fill[slate] (0.038,6.847) circle (0.28pt);
\fill[slate] (0.038,5) circle (0.28pt);
\fill[slate] (0.038,5.249) circle (0.28pt);
\fill[slate] (0.038,5.449) circle (0.28pt);
\fill[slate] (0.038,5.905) circle (0.28pt);
\fill[slate] (0.038,6.149) circle (0.28pt);
\fill[slate] (0.038,6.361) circle (0.28pt);
\fill[slate] (0.038,6.817) circle (0.28pt);
\fill[slate] (0.038,6.397) circle (0.28pt);
\fill[slate] (0.038,6.848) circle (0.28pt);
\fill[slate] (0.038,5.498) circle (0.28pt);
\fill[slate] (0.038,5.747) circle (0.28pt);
\fill[slate] (0.038,5.947) circle (0.28pt);
\fill[slate] (0.038,6.403) circle (0.28pt);
\fill[slate] (0.038,6.647) circle (0.28pt);
\fill[slate] (0.038,5.996) circle (0.28pt);
\fill[slate] (0.038,6.245) circle (0.28pt);
\fill[slate] (0.038,6.445) circle (0.28pt);
\fill[slate] (0.038,6.494) circle (0.28pt);
\fill[slate] (0.038,6.743) circle (0.28pt);
\fill[slate] (0.006,3.014) circle (0.28pt);
\fill[slate] (0.006,3.258) circle (0.28pt);
\fill[slate] (0.006,3.47) circle (0.28pt);
\fill[slate] (0.006,3.926) circle (0.28pt);
\fill[slate] (0.006,4.169) circle (0.28pt);
\fill[slate] (0.006,4.382) circle (0.28pt);
\fill[slate] (0.006,4.838) circle (0.28pt);
\fill[slate] (0.006,5.081) circle (0.28pt);
\fill[slate] (0.006,5.294) circle (0.28pt);
\fill[slate] (0.006,5.75) circle (0.28pt);
\fill[slate] (0.006,5.994) circle (0.28pt);
\fill[slate] (0.006,6.206) circle (0.28pt);
\fill[slate] (0.006,6.663) circle (0.28pt);
\fill[slate] (0.006,6.242) circle (0.28pt);
\fill[slate] (0.006,6.491) circle (0.28pt);
\fill[slate] (0.006,6.693) circle (0.28pt);
\fill[slate] (0.006,6.74) circle (0.28pt);
\fill[slate] (0.006,5.33) circle (0.28pt);
\fill[slate] (0.006,5.579) circle (0.28pt);
\fill[slate] (0.006,5.781) circle (0.28pt);
\fill[slate] (0.006,6.237) circle (0.28pt);
\fill[slate] (0.006,6.48) circle (0.28pt);
\fill[slate] (0.006,6.693) circle (0.28pt);
\fill[slate] (0.006,6.729) circle (0.28pt);
\fill[slate] (0.006,5.828) circle (0.28pt);
\fill[slate] (0.006,6.077) circle (0.28pt);
\fill[slate] (0.006,6.277) circle (0.28pt);
\fill[slate] (0.006,6.733) circle (0.28pt);
\fill[slate] (0.006,6.326) circle (0.28pt);
\fill[slate] (0.006,6.575) circle (0.28pt);
\fill[slate] (0.006,6.775) circle (0.28pt);
\fill[slate] (0.006,6.824) circle (0.28pt);
\fill[slate] (0.006,4.418) circle (0.28pt);
\fill[slate] (0.006,4.666) circle (0.28pt);
\fill[slate] (0.006,4.869) circle (0.28pt);
\fill[slate] (0.006,5.325) circle (0.28pt);
\fill[slate] (0.006,5.568) circle (0.28pt);
\fill[slate] (0.006,5.781) circle (0.28pt);
\fill[slate] (0.006,6.237) circle (0.28pt);
\fill[slate] (0.006,6.48) circle (0.28pt);
\fill[slate] (0.006,6.693) circle (0.28pt);
\fill[slate] (0.006,6.728) circle (0.28pt);
\fill[slate] (0.006,5.817) circle (0.28pt);
\fill[slate] (0.006,6.065) circle (0.28pt);
\fill[slate] (0.006,6.268) circle (0.28pt);
\fill[slate] (0.006,6.724) circle (0.28pt);
\fill[slate] (0.006,6.314) circle (0.28pt);
\fill[slate] (0.006,6.563) circle (0.28pt);
\fill[slate] (0.006,6.764) circle (0.28pt);
\fill[slate] (0.006,6.812) circle (0.28pt);
\fill[slate] (0.006,4.915) circle (0.28pt);
\fill[slate] (0.006,5.164) circle (0.28pt);
\fill[slate] (0.006,5.365) circle (0.28pt);
\fill[slate] (0.006,5.821) circle (0.28pt);
\fill[slate] (0.006,6.065) circle (0.28pt);
\fill[slate] (0.006,6.277) circle (0.28pt);
\fill[slate] (0.006,6.733) circle (0.28pt);
\fill[slate] (0.006,6.313) circle (0.28pt);
\fill[slate] (0.006,6.562) circle (0.28pt);
\fill[slate] (0.006,6.764) circle (0.28pt);
\fill[slate] (0.006,6.811) circle (0.28pt);
\fill[slate] (0.006,5.413) circle (0.28pt);
\fill[slate] (0.006,5.662) circle (0.28pt);
\fill[slate] (0.006,5.863) circle (0.28pt);
\fill[slate] (0.006,6.319) circle (0.28pt);
\fill[slate] (0.006,6.562) circle (0.28pt);
\fill[slate] (0.006,6.775) circle (0.28pt);
\fill[slate] (0.006,6.811) circle (0.28pt);
\fill[slate] (0.006,5.911) circle (0.28pt);
\fill[slate] (0.006,6.16) circle (0.28pt);
\fill[slate] (0.006,6.361) circle (0.28pt);
\fill[slate] (0.006,6.817) circle (0.28pt);
\fill[slate] (0.006,6.409) circle (0.28pt);
\fill[slate] (0.006,6.658) circle (0.28pt);
\fill[slate] (0.006,3.506) circle (0.28pt);
\fill[slate] (0.006,3.957) circle (0.28pt);
\fill[slate] (0.006,4.413) circle (0.28pt);
\fill[slate] (0.006,4.657) circle (0.28pt);
\fill[slate] (0.006,4.869) circle (0.28pt);
\fill[slate] (0.006,5.325) circle (0.28pt);
\fill[slate] (0.006,5.568) circle (0.28pt);
\fill[slate] (0.006,5.781) circle (0.28pt);
\fill[slate] (0.006,6.237) circle (0.28pt);
\fill[slate] (0.006,6.481) circle (0.28pt);
\fill[slate] (0.006,6.693) circle (0.28pt);
\fill[slate] (0.006,6.729) circle (0.28pt);
\fill[slate] (0.006,5.817) circle (0.28pt);
\fill[slate] (0.006,6.065) circle (0.28pt);
\fill[slate] (0.006,6.268) circle (0.28pt);
\fill[slate] (0.006,6.724) circle (0.28pt);
\fill[slate] (0.006,6.314) circle (0.28pt);
\fill[slate] (0.006,6.563) circle (0.28pt);
\fill[slate] (0.006,6.764) circle (0.28pt);
\fill[slate] (0.006,6.812) circle (0.28pt);
\fill[slate] (0.006,4.905) circle (0.28pt);
\fill[slate] (0.006,5.154) circle (0.28pt);
\fill[slate] (0.006,5.356) circle (0.28pt);
\fill[slate] (0.006,5.812) circle (0.28pt);
\fill[slate] (0.006,6.056) circle (0.28pt);
\fill[slate] (0.006,6.268) circle (0.28pt);
\fill[slate] (0.006,6.724) circle (0.28pt);
\fill[slate] (0.006,6.304) circle (0.28pt);
\fill[slate] (0.006,6.553) circle (0.28pt);
\fill[slate] (0.006,6.755) circle (0.28pt);
\fill[slate] (0.006,6.802) circle (0.28pt);
\fill[slate] (0.006,5.403) circle (0.28pt);
\fill[slate] (0.006,5.652) circle (0.28pt);
\fill[slate] (0.006,5.852) circle (0.28pt);
\fill[slate] (0.006,6.308) circle (0.28pt);
\fill[slate] (0.006,6.552) circle (0.28pt);
\fill[slate] (0.006,6.764) circle (0.28pt);
\fill[slate] (0.006,6.8) circle (0.28pt);
\fill[slate] (0.006,5.901) circle (0.28pt);
\fill[slate] (0.006,6.15) circle (0.28pt);
\fill[slate] (0.006,6.35) circle (0.28pt);
\fill[slate] (0.006,6.806) circle (0.28pt);
\fill[slate] (0.006,6.399) circle (0.28pt);
\fill[slate] (0.006,6.648) circle (0.28pt);
\fill[slate] (0.006,6.848) circle (0.28pt);
\fill[slate] (0.006,4.453) circle (0.28pt);
\fill[slate] (0.006,4.909) circle (0.28pt);
\fill[slate] (0.006,5.153) circle (0.28pt);
\fill[slate] (0.006,5.365) circle (0.28pt);
\fill[slate] (0.006,5.821) circle (0.28pt);
\fill[slate] (0.006,6.065) circle (0.28pt);
\fill[slate] (0.006,6.278) circle (0.28pt);
\fill[slate] (0.006,6.734) circle (0.28pt);
\fill[slate] (0.006,6.313) circle (0.28pt);
\fill[slate] (0.006,6.562) circle (0.28pt);
\fill[slate] (0.006,6.764) circle (0.28pt);
\fill[slate] (0.006,6.811) circle (0.28pt);
\fill[slate] (0.006,5.401) circle (0.28pt);
\fill[slate] (0.006,5.65) circle (0.28pt);
\fill[slate] (0.006,5.852) circle (0.28pt);
\fill[slate] (0.006,6.308) circle (0.28pt);
\fill[slate] (0.006,6.552) circle (0.28pt);
\fill[slate] (0.006,6.764) circle (0.28pt);
\fill[slate] (0.006,6.8) circle (0.28pt);
\fill[slate] (0.006,5.899) circle (0.28pt);
\fill[slate] (0.006,6.148) circle (0.28pt);
\fill[slate] (0.006,6.349) circle (0.28pt);
\fill[slate] (0.006,6.805) circle (0.28pt);
\fill[slate] (0.006,6.397) circle (0.28pt);
\fill[slate] (0.006,6.646) circle (0.28pt);
\fill[slate] (0.006,6.847) circle (0.28pt);
\fill[slate] (0.006,4.751) circle (0.28pt);
\fill[slate] (0.006,4.951) circle (0.28pt);
\fill[slate] (0.006,5.407) circle (0.28pt);
\fill[slate] (0.006,5.651) circle (0.28pt);
\fill[slate] (0.006,5.863) circle (0.28pt);
\fill[slate] (0.006,6.319) circle (0.28pt);
\fill[slate] (0.006,6.563) circle (0.28pt);
\fill[slate] (0.006,6.775) circle (0.28pt);
\fill[slate] (0.006,6.811) circle (0.28pt);
\fill[slate] (0.006,5.899) circle (0.28pt);
\fill[slate] (0.006,6.148) circle (0.28pt);
\fill[slate] (0.006,6.35) circle (0.28pt);
\fill[slate] (0.006,6.806) circle (0.28pt);
\fill[slate] (0.006,6.397) circle (0.28pt);
\fill[slate] (0.006,6.847) circle (0.28pt);
\fill[slate] (0.006,5) circle (0.28pt);
\fill[slate] (0.006,5.249) circle (0.28pt);
\fill[slate] (0.006,5.905) circle (0.28pt);
\fill[slate] (0.006,6.149) circle (0.28pt);
\fill[slate] (0.006,6.361) circle (0.28pt);
\fill[slate] (0.006,6.817) circle (0.28pt);
\fill[slate] (0.006,6.397) circle (0.28pt);
\fill[slate] (0.006,6.848) circle (0.28pt);
\fill[slate] (0.006,5.498) circle (0.28pt);
\fill[slate] (0.006,5.747) circle (0.28pt);
\fill[slate] (0.006,6.403) circle (0.28pt);
\fill[slate] (0.006,6.647) circle (0.28pt);
\fill[slate] (0.006,5.996) circle (0.28pt);
\fill[slate] (0.006,6.245) circle (0.28pt);
\fill[slate] (0.006,6.445) circle (0.28pt);
\fill[slate] (0.006,6.494) circle (0.28pt);
\fill[slate] (0.006,6.743) circle (0.28pt);
\fill[slate] (0.001,3.512) circle (0.28pt);
\fill[slate] (0.001,3.756) circle (0.28pt);
\fill[slate] (0.001,3.968) circle (0.28pt);
\fill[slate] (0.001,4.424) circle (0.28pt);
\fill[slate] (0.001,4.667) circle (0.28pt);
\fill[slate] (0.001,4.88) circle (0.28pt);
\fill[slate] (0.001,5.336) circle (0.28pt);
\fill[slate] (0.001,5.579) circle (0.28pt);
\fill[slate] (0.001,5.792) circle (0.28pt);
\fill[slate] (0.001,6.248) circle (0.28pt);
\fill[slate] (0.001,6.492) circle (0.28pt);
\fill[slate] (0.001,6.704) circle (0.28pt);
\fill[slate] (0.001,6.74) circle (0.28pt);
\fill[slate] (0.001,5.828) circle (0.28pt);
\fill[slate] (0.001,6.077) circle (0.28pt);
\fill[slate] (0.001,6.279) circle (0.28pt);
\fill[slate] (0.001,6.735) circle (0.28pt);
\fill[slate] (0.001,6.326) circle (0.28pt);
\fill[slate] (0.001,6.575) circle (0.28pt);
\fill[slate] (0.001,6.775) circle (0.28pt);
\fill[slate] (0.001,6.824) circle (0.28pt);
\fill[slate] (0.001,4.916) circle (0.28pt);
\fill[slate] (0.001,5.164) circle (0.28pt);
\fill[slate] (0.001,5.367) circle (0.28pt);
\fill[slate] (0.001,5.823) circle (0.28pt);
\fill[slate] (0.001,6.066) circle (0.28pt);
\fill[slate] (0.001,6.279) circle (0.28pt);
\fill[slate] (0.001,6.735) circle (0.28pt);
\fill[slate] (0.001,6.315) circle (0.28pt);
\fill[slate] (0.001,6.563) circle (0.28pt);
\fill[slate] (0.001,6.766) circle (0.28pt);
\fill[slate] (0.001,6.812) circle (0.28pt);
\fill[slate] (0.001,5.413) circle (0.28pt);
\fill[slate] (0.001,5.662) circle (0.28pt);
\fill[slate] (0.001,5.863) circle (0.28pt);
\fill[slate] (0.001,6.319) circle (0.28pt);
\fill[slate] (0.001,6.563) circle (0.28pt);
\fill[slate] (0.001,6.775) circle (0.28pt);
\fill[slate] (0.001,6.811) circle (0.28pt);
\fill[slate] (0.001,5.911) circle (0.28pt);
\fill[slate] (0.001,6.16) circle (0.28pt);
\fill[slate] (0.001,6.361) circle (0.28pt);
\fill[slate] (0.001,6.817) circle (0.28pt);
\fill[slate] (0.001,6.409) circle (0.28pt);
\fill[slate] (0.001,6.658) circle (0.28pt);
\fill[slate] (0.001,4.004) circle (0.28pt);
\fill[slate] (0.001,4.253) circle (0.28pt);
\fill[slate] (0.001,4.455) circle (0.28pt);
\fill[slate] (0.001,4.911) circle (0.28pt);
\fill[slate] (0.001,5.155) circle (0.28pt);
\fill[slate] (0.001,5.367) circle (0.28pt);
\fill[slate] (0.001,5.823) circle (0.28pt);
\fill[slate] (0.001,6.066) circle (0.28pt);
\fill[slate] (0.001,6.279) circle (0.28pt);
\fill[slate] (0.001,6.735) circle (0.28pt);
\fill[slate] (0.001,6.315) circle (0.28pt);
\fill[slate] (0.001,6.563) circle (0.28pt);
\fill[slate] (0.001,6.766) circle (0.28pt);
\fill[slate] (0.001,6.812) circle (0.28pt);
\fill[slate] (0.001,5.403) circle (0.28pt);
\fill[slate] (0.001,5.652) circle (0.28pt);
\fill[slate] (0.001,5.854) circle (0.28pt);
\fill[slate] (0.001,6.31) circle (0.28pt);
\fill[slate] (0.001,6.554) circle (0.28pt);
\fill[slate] (0.001,6.766) circle (0.28pt);
\fill[slate] (0.001,6.802) circle (0.28pt);
\fill[slate] (0.001,5.901) circle (0.28pt);
\fill[slate] (0.001,6.15) circle (0.28pt);
\fill[slate] (0.001,6.35) circle (0.28pt);
\fill[slate] (0.001,6.806) circle (0.28pt);
\fill[slate] (0.001,6.399) circle (0.28pt);
\fill[slate] (0.001,6.648) circle (0.28pt);
\fill[slate] (0.001,6.848) circle (0.28pt);
\fill[slate] (0.001,4.951) circle (0.28pt);
\fill[slate] (0.001,5.407) circle (0.28pt);
\fill[slate] (0.001,5.651) circle (0.28pt);
\fill[slate] (0.001,5.863) circle (0.28pt);
\fill[slate] (0.001,6.319) circle (0.28pt);
\fill[slate] (0.001,6.563) circle (0.28pt);
\fill[slate] (0.001,6.776) circle (0.28pt);
\fill[slate] (0.001,6.811) circle (0.28pt);
\fill[slate] (0.001,5.899) circle (0.28pt);
\fill[slate] (0.001,6.148) circle (0.28pt);
\fill[slate] (0.001,6.35) circle (0.28pt);
\fill[slate] (0.001,6.806) circle (0.28pt);
\fill[slate] (0.001,6.397) circle (0.28pt);
\fill[slate] (0.001,6.646) circle (0.28pt);
\fill[slate] (0.001,6.847) circle (0.28pt);
\fill[slate] (0.001,5.449) circle (0.28pt);
\fill[slate] (0.001,5.905) circle (0.28pt);
\fill[slate] (0.001,6.149) circle (0.28pt);
\fill[slate] (0.001,6.361) circle (0.28pt);
\fill[slate] (0.001,6.817) circle (0.28pt);
\fill[slate] (0.001,6.397) circle (0.28pt);
\fill[slate] (0.001,6.646) circle (0.28pt);
\fill[slate] (0.001,6.848) circle (0.28pt);
\fill[slate] (0.001,5.498) circle (0.28pt);
\fill[slate] (0.001,5.747) circle (0.28pt);
\fill[slate] (0.001,5.947) circle (0.28pt);
\fill[slate] (0.001,6.403) circle (0.28pt);
\fill[slate] (0.001,6.647) circle (0.28pt);
\fill[slate] (0.001,5.996) circle (0.28pt);
\fill[slate] (0.001,6.245) circle (0.28pt);
\fill[slate] (0.001,6.494) circle (0.28pt);
\fill[slate] (0.001,6.743) circle (0.28pt);
\fill[slate] (0,4.01) circle (0.28pt);
\fill[slate] (0,4.254) circle (0.28pt);
\fill[slate] (0,4.466) circle (0.28pt);
\fill[slate] (0,4.922) circle (0.28pt);
\fill[slate] (0,5.165) circle (0.28pt);
\fill[slate] (0,5.378) circle (0.28pt);
\fill[slate] (0,5.834) circle (0.28pt);
\fill[slate] (0,6.078) circle (0.28pt);
\fill[slate] (0,6.29) circle (0.28pt);
\fill[slate] (0,6.746) circle (0.28pt);
\fill[slate] (0,6.326) circle (0.28pt);
\fill[slate] (0,6.575) circle (0.28pt);
\fill[slate] (0,6.777) circle (0.28pt);
\fill[slate] (0,6.824) circle (0.28pt);
\fill[slate] (0,5.414) circle (0.28pt);
\fill[slate] (0,5.662) circle (0.28pt);
\fill[slate] (0,5.865) circle (0.28pt);
\fill[slate] (0,6.321) circle (0.28pt);
\fill[slate] (0,6.564) circle (0.28pt);
\fill[slate] (0,6.777) circle (0.28pt);
\fill[slate] (0,6.813) circle (0.28pt);
\fill[slate] (0,5.911) circle (0.28pt);
\fill[slate] (0,6.16) circle (0.28pt);
\fill[slate] (0,6.361) circle (0.28pt);
\fill[slate] (0,6.817) circle (0.28pt);
\fill[slate] (0,6.409) circle (0.28pt);
\fill[slate] (0,6.658) circle (0.28pt);
\fill[slate] (0,4.502) circle (0.28pt);
\fill[slate] (0,4.751) circle (0.28pt);
\fill[slate] (0,4.953) circle (0.28pt);
\fill[slate] (0,5.409) circle (0.28pt);
\fill[slate] (0,5.653) circle (0.28pt);
\fill[slate] (0,5.865) circle (0.28pt);
\fill[slate] (0,6.321) circle (0.28pt);
\fill[slate] (0,6.564) circle (0.28pt);
\fill[slate] (0,6.777) circle (0.28pt);
\fill[slate] (0,6.813) circle (0.28pt);
\fill[slate] (0,5.901) circle (0.28pt);
\fill[slate] (0,6.15) circle (0.28pt);
\fill[slate] (0,6.352) circle (0.28pt);
\fill[slate] (0,6.808) circle (0.28pt);
\fill[slate] (0,6.399) circle (0.28pt);
\fill[slate] (0,6.648) circle (0.28pt);
\fill[slate] (0,6.848) circle (0.28pt);
\fill[slate] (0,5) circle (0.28pt);
\fill[slate] (0,5.449) circle (0.28pt);
\fill[slate] (0,5.905) circle (0.28pt);
\fill[slate] (0,6.149) circle (0.28pt);
\fill[slate] (0,6.361) circle (0.28pt);
\fill[slate] (0,6.817) circle (0.28pt);
\fill[slate] (0,6.397) circle (0.28pt);
\fill[slate] (0,6.646) circle (0.28pt);
\fill[slate] (0,6.848) circle (0.28pt);
\fill[slate] (0,5.947) circle (0.28pt);
\fill[slate] (0,6.403) circle (0.28pt);
\fill[slate] (0,6.647) circle (0.28pt);
\fill[slate] (0,6.245) circle (0.28pt);
\fill[slate] (0,6.445) circle (0.28pt);
\fill[slate] (0,6.494) circle (0.28pt);
\fill[slate] (0,6.743) circle (0.28pt);
\fill[slate] (0,4.508) circle (0.28pt);
\fill[slate] (0,4.752) circle (0.28pt);
\fill[slate] (0,4.964) circle (0.28pt);
\fill[slate] (0,5.42) circle (0.28pt);
\fill[slate] (0,5.663) circle (0.28pt);
\fill[slate] (0,5.876) circle (0.28pt);
\fill[slate] (0,6.332) circle (0.28pt);
\fill[slate] (0,6.576) circle (0.28pt);
\fill[slate] (0,6.788) circle (0.28pt);
\fill[slate] (0,6.824) circle (0.28pt);
\fill[slate] (0,5.912) circle (0.28pt);
\fill[slate] (0,6.16) circle (0.28pt);
\fill[slate] (0,6.363) circle (0.28pt);
\fill[slate] (0,6.819) circle (0.28pt);
\fill[slate] (0,6.409) circle (0.28pt);
\fill[slate] (0,6.658) circle (0.28pt);
\fill[slate] (0,5) circle (0.28pt);
\fill[slate] (0,5.249) circle (0.28pt);
\fill[slate] (0,5.451) circle (0.28pt);
\fill[slate] (0,5.907) circle (0.28pt);
\fill[slate] (0,6.151) circle (0.28pt);
\fill[slate] (0,6.363) circle (0.28pt);
\fill[slate] (0,6.819) circle (0.28pt);
\fill[slate] (0,6.399) circle (0.28pt);
\fill[slate] (0,6.648) circle (0.28pt);
\fill[slate] (0,6.85) circle (0.28pt);
\fill[slate] (0,5.498) circle (0.28pt);
\fill[slate] (0,5.747) circle (0.28pt);
\fill[slate] (0,5.947) circle (0.28pt);
\fill[slate] (0,6.403) circle (0.28pt);
\fill[slate] (0,6.647) circle (0.28pt);
\fill[slate] (0,6.445) circle (0.28pt);
\fill[slate] (0,5.006) circle (0.28pt);
\fill[slate] (0,5.25) circle (0.28pt);
\fill[slate] (0,5.462) circle (0.28pt);
\fill[slate] (0,5.918) circle (0.28pt);
\fill[slate] (0,6.161) circle (0.28pt);
\fill[slate] (0,6.374) circle (0.28pt);
\fill[slate] (0,6.83) circle (0.28pt);
\fill[slate] (0,6.41) circle (0.28pt);
\fill[slate] (0,6.658) circle (0.28pt);
\fill[slate] (0,5.498) circle (0.28pt);
\fill[slate] (0,5.747) circle (0.28pt);
\fill[slate] (0,5.949) circle (0.28pt);
\fill[slate] (0,6.405) circle (0.28pt);
\fill[slate] (0,6.649) circle (0.28pt);
\fill[slate] (0,5.996) circle (0.28pt);
\fill[slate] (0,6.245) circle (0.28pt);
\fill[slate] (0,6.445) circle (0.28pt);
\fill[slate] (0,6.494) circle (0.28pt);
\fill[slate] (0,5.504) circle (0.28pt);
\fill[slate] (0,5.748) circle (0.28pt);
\fill[slate] (0,5.96) circle (0.28pt);
\fill[slate] (0,6.416) circle (0.28pt);
\fill[slate] (0,6.659) circle (0.28pt);
\fill[slate] (0,5.996) circle (0.28pt);
\fill[slate] (0,6.245) circle (0.28pt);
\fill[slate] (0,6.447) circle (0.28pt);
\fill[slate] (0,6.494) circle (0.28pt);
\fill[slate] (0,6.743) circle (0.28pt);
\fill[slate] (0,6.002) circle (0.28pt);
\fill[slate] (0,6.246) circle (0.28pt);
\fill[slate] (0,6.458) circle (0.28pt);
\fill[slate] (0,6.494) circle (0.28pt);
\fill[slate] (0,6.743) circle (0.28pt);
\fill[slate] (0,6.5) circle (0.28pt);
\fill[slate] (0,6.744) circle (0.28pt);
\fill[navy] (6.935,0.119) circle (0.75pt);
\fill[navy] (1.019,0.366) circle (0.75pt);
\fill[navy] (0.149,0.614) circle (0.75pt);
\fill[navy] (0.022,0.863) circle (0.75pt);
\fill[navy] (0.003,1.112) circle (0.75pt);
\fill[navy] (0,1.361) circle (0.75pt);
\fill[navy] (0,1.61) circle (0.75pt);
\fill[navy] (0,1.859) circle (0.75pt);
\fill[navy] (0,2.108) circle (0.75pt);
\fill[navy] (0,2.357) circle (0.75pt);
\fill[navy] (0,2.607) circle (0.75pt);
\fill[navy] (0,2.856) circle (0.75pt);
\fill[navy] (0,3.105) circle (0.75pt);
\fill[navy] (0,3.354) circle (0.75pt);
\fill[navy] (0,3.603) circle (0.75pt);
\fill[navy] (0,3.852) circle (0.75pt);
\fill[navy] (0,4.101) circle (0.75pt);
\fill[navy] (0,4.35) circle (0.75pt);
\fill[navy] (0,4.599) circle (0.75pt);
\fill[navy] (0,4.848) circle (0.75pt);
\fill[navy] (0,5.097) circle (0.75pt);
\fill[navy] (0,5.346) circle (0.75pt);
\fill[navy] (0,5.595) circle (0.75pt);
\fill[navy] (0,5.844) circle (0.75pt);
\fill[navy] (0,6.093) circle (0.75pt);
\fill[navy] (0,6.342) circle (0.75pt);
\fill[navy] (0,6.591) circle (0.75pt);
\fill[navy] (0,6.84) circle (0.75pt);
\fill[navy] (0,6.542) circle (0.75pt);
\fill[navy] (0,6.044) circle (0.75pt);
\fill[navy] (0,5.546) circle (0.75pt);
\fill[navy] (0,5.048) circle (0.75pt);
\fill[navy] (0,4.55) circle (0.75pt);
\fill[navy] (0,4.052) circle (0.75pt);
\fill[navy] (0,6.494) circle (0.75pt);
\fill[navy] (0,6.245) circle (0.75pt);
\fill[navy] (0,3.554) circle (0.75pt);
\fill[navy] (0,5.747) circle (0.75pt);
\fill[navy] (0,5.498) circle (0.75pt);
\fill[navy] (0.001,3.056) circle (0.75pt);
\fill[navy] (0.001,5) circle (0.75pt);
\fill[navy] (0.001,4.751) circle (0.75pt);
\fill[navy] (0.001,6.445) circle (0.75pt);
\fill[navy] (0.006,2.558) circle (0.75pt);
\fill[navy] (0.006,4.253) circle (0.75pt);
\fill[navy] (0.006,5.947) circle (0.75pt);
\fill[navy] (0.006,4.004) circle (0.75pt);
\fill[navy] (0.006,5.449) circle (0.75pt);
\fill[navy] (0.038,2.06) circle (0.75pt);
\fill[navy] (0.038,3.506) circle (0.75pt);
\fill[navy] (0.038,4.951) circle (0.75pt);
\fill[navy] (0.038,6.397) circle (0.75pt);
\fill[navy] (0.038,6.148) circle (0.75pt);
\fill[navy] (0.038,3.257) circle (0.75pt);
\fill[navy] (0.038,5.899) circle (0.75pt);
\fill[navy] (0.038,4.453) circle (0.75pt);
\fill[navy] (0.038,5.65) circle (0.75pt);
\fill[navy] (0.038,6.846) circle (0.75pt);
\fill[navy] (0.26,1.562) circle (0.75pt);
\fill[navy] (0.26,2.759) circle (0.75pt);
\fill[navy] (0.26,3.955) circle (0.75pt);
\fill[navy] (0.26,5.152) circle (0.75pt);
\fill[navy] (0.26,6.349) circle (0.75pt);
\fill[navy] (0.26,4.903) circle (0.75pt);
\fill[navy] (0.26,2.51) circle (0.75pt);
\fill[navy] (0.26,4.654) circle (0.75pt);
\fill[navy] (0.26,6.798) circle (0.75pt);
\fill[navy] (0.26,3.458) circle (0.75pt);
\fill[navy] (0.26,6.55) circle (0.75pt);
\fill[navy] (0.26,4.405) circle (0.75pt);
\fill[navy] (0.26,5.353) circle (0.75pt);
\fill[navy] (0.26,6.301) circle (0.75pt);
\fill[navy] (1.781,1.065) circle (0.75pt);
\fill[navy] (1.778,2.013) circle (0.75pt);
\fill[navy] (1.778,2.96) circle (0.75pt);
\fill[navy] (1.778,3.908) circle (0.75pt);
\fill[navy] (1.778,4.856) circle (0.75pt);
\fill[navy] (1.778,5.804) circle (0.75pt);
\fill[navy] (1.778,6.751) circle (0.75pt);
\fill[navy] (1.778,5.555) circle (0.75pt);
\fill[navy] (1.778,3.66) circle (0.75pt);
\fill[navy] (1.778,6.255) circle (0.75pt);
\fill[navy] (1.778,5.307) circle (0.75pt);
\fill[navy] (1.785,1.765) circle (0.75pt);
\fill[navy] (1.784,3.412) circle (0.75pt);
\fill[navy] (1.784,5.059) circle (0.75pt);
\fill[navy] (1.784,6.706) circle (0.75pt);
\fill[navy] (1.784,5.758) circle (0.75pt);
\fill[navy] (1.785,2.464) circle (0.75pt);
\fill[navy] (1.785,4.811) circle (0.75pt);
\fill[navy] (1.785,3.164) circle (0.75pt);
\fill[navy] (1.785,6.21) circle (0.75pt);
\fill[navy] (1.785,3.863) circle (0.75pt);
\fill[navy] (1.785,4.563) circle (0.75pt);
\fill[navy] (1.785,5.262) circle (0.75pt);
\fill[navy] (1.785,5.962) circle (0.75pt);
\fill[navy] (1.785,6.661) circle (0.75pt);
\fill[navy] (12.238,0.573) circle (0.75pt);
\fill[navy] (12.085,1.273) circle (0.75pt);
\fill[navy] (12.085,1.972) circle (0.75pt);
\fill[navy] (12.085,2.672) circle (0.75pt);
\fill[navy] (12.085,3.371) circle (0.75pt);
\fill[navy] (12.085,4.071) circle (0.75pt);
\fill[navy] (12.085,4.77) circle (0.75pt);
\fill[navy] (12.085,5.47) circle (0.75pt);
\fill[navy] (12.085,6.169) circle (0.75pt);
\fill[navy] (12.085,6.625) circle (0.75pt);
\fill[navy] (12.085,5.226) circle (0.75pt);
\fill[navy] (12.085,3.827) circle (0.75pt);
\fill[navy] (12.085,6.382) circle (0.75pt);
\fill[navy] (12.085,5.682) circle (0.75pt);
\fill[navy] (12.086,2.428) circle (0.75pt);
\fill[navy] (12.086,4.283) circle (0.75pt);
\fill[navy] (12.086,6.138) circle (0.75pt);
\fill[navy] (12.086,3.584) circle (0.75pt);
\fill[navy] (12.086,6.594) circle (0.75pt);
\fill[navy] (12.086,4.739) circle (0.75pt);
\fill[navy] (12.086,5.894) circle (0.75pt);
\fill[navy] (12.395,1.029) circle (0.75pt);
\fill[navy] (12.391,2.185) circle (0.75pt);
\fill[navy] (12.391,3.34) circle (0.75pt);
\fill[navy] (12.391,4.496) circle (0.75pt);
\fill[navy] (12.391,5.651) circle (0.75pt);
\fill[navy] (12.391,6.806) circle (0.75pt);
\fill[navy] (12.391,6.107) circle (0.75pt);
\fill[navy] (12.391,3.796) circle (0.75pt);
\fill[navy] (12.391,6.563) circle (0.75pt);
\fill[navy] (12.391,5.408) circle (0.75pt);
\fill[navy] (12.4,1.485) circle (0.75pt);
\fill[navy] (12.4,3.097) circle (0.75pt);
\fill[navy] (12.4,4.708) circle (0.75pt);
\fill[navy] (12.4,6.32) circle (0.75pt);
\fill[navy] (12.4,5.165) circle (0.75pt);
\fill[navy] (12.4,1.942) circle (0.75pt);
\fill[navy] (12.4,4.009) circle (0.75pt);
\fill[navy] (12.4,6.077) circle (0.75pt);
\fill[navy] (12.4,6.533) circle (0.75pt);
\fill[navy] (12.4,2.398) circle (0.75pt);
\fill[navy] (12.4,4.921) circle (0.75pt);
\fill[navy] (12.4,2.854) circle (0.75pt);
\fill[navy] (12.4,5.833) circle (0.75pt);
\fill[navy] (12.4,3.31) circle (0.75pt);
\fill[navy] (12.4,6.746) circle (0.75pt);
\fill[navy] (12.4,3.766) circle (0.75pt);
\fill[navy] (12.4,4.222) circle (0.75pt);
\fill[navy] (12.4,4.678) circle (0.75pt);
\fill[navy] (12.4,5.134) circle (0.75pt);
\fill[navy] (12.4,5.59) circle (0.75pt);
\fill[navy] (12.4,6.046) circle (0.75pt);
\fill[navy] (12.4,6.502) circle (0.75pt);
\draw[slate, densely dotted, line width=0.35pt] (6.935,0.119) -- (6.935,1.521);
\draw[slate, densely dotted, line width=0.35pt] (1.019,0.366) -- (1.019,2.262);
\draw[slate, densely dotted, line width=0.35pt] (12.238,0.573) -- (12.238,2.884);
\draw[slate, densely dotted, line width=0.35pt] (6.935,0.119) -- (6.935,2.92);
\draw[slate, densely dotted, line width=0.35pt] (6.935,0.119) -- (6.935,2.925);
\draw[slate, densely dotted, line width=0.35pt] (0.149,0.614) -- (0.149,3.008);
\draw[slate, densely dotted, line width=0.35pt] (0.022,0.863) -- (0.022,3.755);
\draw[slate, densely dotted, line width=0.35pt] (1.019,0.366) -- (1.019,4.157);
\draw[slate, densely dotted, line width=0.35pt] (1.019,0.366) -- (1.019,4.158);
\draw[slate, densely dotted, line width=0.35pt] (12.395,1.029) -- (12.395,4.252);
\draw[slate, densely dotted, line width=0.35pt] (6.935,0.119) -- (6.935,4.319);
\draw[slate, densely dotted, line width=0.35pt] (6.935,0.119) -- (6.935,4.324);
\draw[slate, densely dotted, line width=0.35pt] (6.935,0.119) -- (6.935,4.324);
\draw[slate, densely dotted, line width=0.35pt] (6.935,0.119) -- (6.935,4.328);
\draw[slate, densely dotted, line width=0.35pt] (1.781,1.065) -- (1.781,4.359);
\draw[slate, densely dotted, line width=0.35pt] (0.003,1.112) -- (0.003,4.502);
\draw[slate, densely dotted, line width=0.35pt] (12.085,1.273) -- (12.085,4.983);
\draw[slate, densely dotted, line width=0.35pt] (12.238,0.573) -- (12.238,5.195);
\draw[slate, densely dotted, line width=0.35pt] (12.238,0.573) -- (12.238,5.195);
\draw[slate, densely dotted, line width=0.35pt] (0,1.361) -- (0,5.249);
\draw[slate, densely dotted, line width=0.35pt] (0.149,0.614) -- (0.149,5.401);
\draw[slate, densely dotted, line width=0.35pt] (0.149,0.614) -- (0.149,5.401);
\draw[slate, densely dotted, line width=0.35pt] (12.4,1.485) -- (12.4,5.621);
\draw[slate, densely dotted, line width=0.35pt] (6.935,0.119) -- (6.935,5.718);
\draw[slate, densely dotted, line width=0.35pt] (6.935,0.119) -- (6.935,5.723);
\draw[slate, densely dotted, line width=0.35pt] (6.935,0.119) -- (6.935,5.723);
\draw[slate, densely dotted, line width=0.35pt] (6.935,0.119) -- (6.935,5.723);
\draw[slate, densely dotted, line width=0.35pt] (6.935,0.119) -- (6.935,5.727);
\draw[slate, densely dotted, line width=0.35pt] (1.83,1.521) -- (1.83,5.727);
\draw[slate, densely dotted, line width=0.35pt] (6.935,0.119) -- (6.935,5.727);
\draw[slate, densely dotted, line width=0.35pt] (6.935,0.119) -- (6.935,5.727);
\draw[slate, densely dotted, line width=0.35pt] (12.039,1.521) -- (12.039,5.727);
\draw[slate, densely dotted, line width=0.35pt] (6.935,0.119) -- (6.935,5.732);
\draw[slate, densely dotted, line width=0.35pt] (0.26,1.562) -- (0.26,5.851);
\draw[slate, densely dotted, line width=0.35pt] (0,1.61) -- (0,5.996);
\draw[slate, densely dotted, line width=0.35pt] (1.019,0.366) -- (1.019,6.052);
\draw[slate, densely dotted, line width=0.35pt] (1.019,0.366) -- (1.019,6.053);
\draw[slate, densely dotted, line width=0.35pt] (1.019,0.366) -- (1.019,6.053);
\draw[slate, densely dotted, line width=0.35pt] (1.019,0.366) -- (1.019,6.054);
\draw[slate, densely dotted, line width=0.35pt] (1.785,1.765) -- (1.785,6.458);
\draw[slate, densely dotted, line width=0.35pt] (12.032,1.77) -- (12.032,6.474);
\draw[slate, densely dotted, line width=0.35pt] (0.022,0.863) -- (0.022,6.646);
\draw[slate, densely dotted, line width=0.35pt] (0.022,0.863) -- (0.022,6.646);
\draw[slate, densely dotted, line width=0.35pt] (0,1.859) -- (0,6.743);
\draw[wine, line width=0.7pt] (1.83,1.521) -- (12.039,1.521);
\draw[wine, line width=0.7pt] (0.261,2.262) -- (1.777,2.262);
\draw[wine, line width=0.7pt] (12.086,2.884) -- (12.391,2.884);
\draw[wine, line width=0.7pt] (1.785,2.92) -- (12.084,2.92);
\draw[wine, line width=0.7pt] (1.83,2.925) -- (12.039,2.925);
\draw[wine, line width=0.7pt] (0.038,3.008) -- (0.26,3.008);
\draw[wine, line width=0.7pt] (0.006,3.755) -- (0.038,3.755);
\draw[wine, line width=0.7pt] (0.26,4.157) -- (1.778,4.157);
\draw[wine, line width=0.7pt] (0.261,4.158) -- (1.777,4.158);
\draw[wine, line width=0.7pt] (12.391,4.252) -- (12.4,4.252);
\draw[wine, line width=0.7pt] (1.785,4.319) -- (12.085,4.319);
\draw[wine, line width=0.7pt] (1.785,4.324) -- (12.084,4.324);
\draw[wine, line width=0.7pt] (1.83,4.324) -- (12.039,4.324);
\draw[wine, line width=0.7pt] (1.83,4.328) -- (12.039,4.328);
\draw[wine, line width=0.7pt] (1.778,4.359) -- (1.784,4.359);
\draw[wine, line width=0.7pt] (0.001,4.502) -- (0.006,4.502);
\draw[wine, line width=0.7pt] (12.085,4.983) -- (12.086,4.983);
\draw[wine, line width=0.7pt] (12.086,5.195) -- (12.391,5.195);
\draw[wine, line width=0.7pt] (12.086,5.195) -- (12.391,5.195);
\draw[wine, line width=0.7pt] (0,5.249) -- (0.001,5.249);
\draw[wine, line width=0.7pt] (0.038,5.401) -- (0.26,5.401);
\draw[wine, line width=0.7pt] (0.038,5.401) -- (0.26,5.401);
\draw[wine, line width=0.7pt] (12.4,5.621) -- (12.4,5.621);
\draw[wine, line width=0.7pt] (1.785,5.718) -- (12.085,5.718);
\draw[wine, line width=0.7pt] (1.785,5.723) -- (12.085,5.723);
\draw[wine, line width=0.7pt] (1.83,5.723) -- (12.039,5.723);
\draw[wine, line width=0.7pt] (1.785,5.723) -- (12.084,5.723);
\draw[wine, line width=0.7pt] (1.785,5.727) -- (12.084,5.727);
\draw[wine, line width=0.7pt] (1.83,5.727) -- (1.83,5.727);
\draw[wine, line width=0.7pt] (1.83,5.727) -- (12.039,5.727);
\draw[wine, line width=0.7pt] (1.83,5.727) -- (12.039,5.727);
\draw[wine, line width=0.7pt] (12.039,5.727) -- (12.039,5.727);
\draw[wine, line width=0.7pt] (1.83,5.732) -- (12.039,5.732);
\draw[wine, line width=0.7pt] (0.26,5.851) -- (0.26,5.851);
\draw[wine, line width=0.7pt] (0,5.996) -- (0,5.996);
\draw[wine, line width=0.7pt] (0.26,6.052) -- (1.778,6.052);
\draw[wine, line width=0.7pt] (0.26,6.053) -- (1.778,6.053);
\draw[wine, line width=0.7pt] (0.261,6.053) -- (1.777,6.053);
\draw[wine, line width=0.7pt] (0.261,6.054) -- (1.777,6.054);
\draw[wine, line width=0.7pt] (1.784,6.458) -- (1.785,6.458);
\draw[wine, line width=0.7pt] (12.032,6.474) -- (12.032,6.474);
\draw[wine, line width=0.7pt] (0.006,6.646) -- (0.038,6.646);
\draw[wine, line width=0.7pt] (0.006,6.646) -- (0.038,6.646);
\draw[wine, line width=0.7pt] (0,6.743) -- (0,6.743);
\draw[wine, fill=white, line width=0.45pt] (0,6.743) circle (1.1pt);
\draw[wine, fill=white, line width=0.45pt] (0,6.743) circle (1.1pt);
\draw[wine, fill=white, line width=0.45pt] (0,5.996) circle (1.1pt);
\draw[wine, fill=white, line width=0.45pt] (0,5.996) circle (1.1pt);
\draw[wine, fill=white, line width=0.45pt] (0,5.249) circle (1.1pt);
\draw[wine, fill=white, line width=0.45pt] (0.001,5.249) circle (1.1pt);
\draw[wine, fill=white, line width=0.45pt] (0.001,4.502) circle (1.1pt);
\draw[wine, fill=white, line width=0.45pt] (0.006,4.502) circle (1.1pt);
\draw[wine, fill=white, line width=0.45pt] (0.006,6.646) circle (1.1pt);
\draw[wine, fill=white, line width=0.45pt] (0.006,3.755) circle (1.1pt);
\draw[wine, fill=white, line width=0.45pt] (0.006,6.646) circle (1.1pt);
\draw[wine, fill=white, line width=0.45pt] (0.038,3.755) circle (1.1pt);
\draw[wine, fill=white, line width=0.45pt] (0.038,6.646) circle (1.1pt);
\draw[wine, fill=white, line width=0.45pt] (0.038,6.646) circle (1.1pt);
\draw[wine, fill=white, line width=0.45pt] (0.038,5.401) circle (1.1pt);
\draw[wine, fill=white, line width=0.45pt] (0.038,3.008) circle (1.1pt);
\draw[wine, fill=white, line width=0.45pt] (0.038,5.401) circle (1.1pt);
\draw[wine, fill=white, line width=0.45pt] (0.26,3.008) circle (1.1pt);
\draw[wine, fill=white, line width=0.45pt] (0.26,5.401) circle (1.1pt);
\draw[wine, fill=white, line width=0.45pt] (0.26,5.401) circle (1.1pt);
\draw[wine, fill=white, line width=0.45pt] (0.26,5.851) circle (1.1pt);
\draw[wine, fill=white, line width=0.45pt] (0.26,5.851) circle (1.1pt);
\draw[wine, fill=white, line width=0.45pt] (0.26,6.052) circle (1.1pt);
\draw[wine, fill=white, line width=0.45pt] (0.26,6.053) circle (1.1pt);
\draw[wine, fill=white, line width=0.45pt] (0.26,4.157) circle (1.1pt);
\draw[wine, fill=white, line width=0.45pt] (0.261,2.262) circle (1.1pt);
\draw[wine, fill=white, line width=0.45pt] (0.261,4.158) circle (1.1pt);
\draw[wine, fill=white, line width=0.45pt] (0.261,6.054) circle (1.1pt);
\draw[wine, fill=white, line width=0.45pt] (0.261,6.053) circle (1.1pt);
\draw[wine, fill=white, line width=0.45pt] (1.777,2.262) circle (1.1pt);
\draw[wine, fill=white, line width=0.45pt] (1.777,6.053) circle (1.1pt);
\draw[wine, fill=white, line width=0.45pt] (1.777,4.158) circle (1.1pt);
\draw[wine, fill=white, line width=0.45pt] (1.777,6.054) circle (1.1pt);
\draw[wine, fill=white, line width=0.45pt] (1.778,4.157) circle (1.1pt);
\draw[wine, fill=white, line width=0.45pt] (1.778,6.053) circle (1.1pt);
\draw[wine, fill=white, line width=0.45pt] (1.778,6.052) circle (1.1pt);
\draw[wine, fill=white, line width=0.45pt] (1.778,4.359) circle (1.1pt);
\draw[wine, fill=white, line width=0.45pt] (1.784,4.359) circle (1.1pt);
\draw[wine, fill=white, line width=0.45pt] (1.784,6.458) circle (1.1pt);
\draw[wine, fill=white, line width=0.45pt] (1.785,6.458) circle (1.1pt);
\draw[wine, fill=white, line width=0.45pt] (1.785,5.718) circle (1.1pt);
\draw[wine, fill=white, line width=0.45pt] (1.785,5.723) circle (1.1pt);
\draw[wine, fill=white, line width=0.45pt] (1.785,4.319) circle (1.1pt);
\draw[wine, fill=white, line width=0.45pt] (1.785,5.727) circle (1.1pt);
\draw[wine, fill=white, line width=0.45pt] (1.785,4.324) circle (1.1pt);
\draw[wine, fill=white, line width=0.45pt] (1.785,2.92) circle (1.1pt);
\draw[wine, fill=white, line width=0.45pt] (1.785,5.723) circle (1.1pt);
\draw[wine, fill=white, line width=0.45pt] (1.83,1.521) circle (1.1pt);
\draw[wine, fill=white, line width=0.45pt] (1.83,5.727) circle (1.1pt);
\draw[wine, fill=white, line width=0.45pt] (1.83,2.925) circle (1.1pt);
\draw[wine, fill=white, line width=0.45pt] (1.83,4.328) circle (1.1pt);
\draw[wine, fill=white, line width=0.45pt] (1.83,5.732) circle (1.1pt);
\draw[wine, fill=white, line width=0.45pt] (1.83,4.324) circle (1.1pt);
\draw[wine, fill=white, line width=0.45pt] (1.83,5.727) circle (1.1pt);
\draw[wine, fill=white, line width=0.45pt] (1.83,5.723) circle (1.1pt);
\draw[wine, fill=white, line width=0.45pt] (12.032,6.474) circle (1.1pt);
\draw[wine, fill=white, line width=0.45pt] (12.032,6.474) circle (1.1pt);
\draw[wine, fill=white, line width=0.45pt] (12.039,1.521) circle (1.1pt);
\draw[wine, fill=white, line width=0.45pt] (12.039,5.723) circle (1.1pt);
\draw[wine, fill=white, line width=0.45pt] (12.039,5.727) circle (1.1pt);
\draw[wine, fill=white, line width=0.45pt] (12.039,4.324) circle (1.1pt);
\draw[wine, fill=white, line width=0.45pt] (12.039,2.925) circle (1.1pt);
\draw[wine, fill=white, line width=0.45pt] (12.039,4.328) circle (1.1pt);
\draw[wine, fill=white, line width=0.45pt] (12.039,5.732) circle (1.1pt);
\draw[wine, fill=white, line width=0.45pt] (12.039,5.727) circle (1.1pt);
\draw[wine, fill=white, line width=0.45pt] (12.084,2.92) circle (1.1pt);
\draw[wine, fill=white, line width=0.45pt] (12.084,5.723) circle (1.1pt);
\draw[wine, fill=white, line width=0.45pt] (12.084,4.324) circle (1.1pt);
\draw[wine, fill=white, line width=0.45pt] (12.084,5.727) circle (1.1pt);
\draw[wine, fill=white, line width=0.45pt] (12.085,4.319) circle (1.1pt);
\draw[wine, fill=white, line width=0.45pt] (12.085,5.723) circle (1.1pt);
\draw[wine, fill=white, line width=0.45pt] (12.085,5.718) circle (1.1pt);
\draw[wine, fill=white, line width=0.45pt] (12.085,4.983) circle (1.1pt);
\draw[wine, fill=white, line width=0.45pt] (12.086,4.983) circle (1.1pt);
\draw[wine, fill=white, line width=0.45pt] (12.086,5.195) circle (1.1pt);
\draw[wine, fill=white, line width=0.45pt] (12.086,2.884) circle (1.1pt);
\draw[wine, fill=white, line width=0.45pt] (12.086,5.195) circle (1.1pt);
\draw[wine, fill=white, line width=0.45pt] (12.391,2.884) circle (1.1pt);
\draw[wine, fill=white, line width=0.45pt] (12.391,5.195) circle (1.1pt);
\draw[wine, fill=white, line width=0.45pt] (12.391,5.195) circle (1.1pt);
\draw[wine, fill=white, line width=0.45pt] (12.391,4.252) circle (1.1pt);
\draw[wine, fill=white, line width=0.45pt] (12.4,4.252) circle (1.1pt);
\draw[wine, fill=white, line width=0.45pt] (12.4,5.621) circle (1.1pt);
\draw[wine, fill=white, line width=0.45pt] (12.4,5.621) circle (1.1pt);
\filldraw[ochre, draw=ink, line width=0.25pt] (6.935,0.119) +(0,1.9pt) -- +(1.9pt,0) -- +(0,-1.9pt) -- +(-1.9pt,0) -- cycle;
\filldraw[ochre, draw=ink, line width=0.25pt] (1.019,0.366) +(0,1.9pt) -- +(1.9pt,0) -- +(0,-1.9pt) -- +(-1.9pt,0) -- cycle;
\filldraw[ochre, draw=ink, line width=0.25pt] (0.149,0.614) +(0,1.9pt) -- +(1.9pt,0) -- +(0,-1.9pt) -- +(-1.9pt,0) -- cycle;
\filldraw[ochre, draw=ink, line width=0.25pt] (0.022,0.863) +(0,1.9pt) -- +(1.9pt,0) -- +(0,-1.9pt) -- +(-1.9pt,0) -- cycle;
\filldraw[ochre, draw=ink, line width=0.25pt] (0.003,1.112) +(0,1.9pt) -- +(1.9pt,0) -- +(0,-1.9pt) -- +(-1.9pt,0) -- cycle;
\filldraw[ochre, draw=ink, line width=0.25pt] (0,1.361) +(0,1.9pt) -- +(1.9pt,0) -- +(0,-1.9pt) -- +(-1.9pt,0) -- cycle;
\filldraw[ochre, draw=ink, line width=0.25pt] (0,1.61) +(0,1.9pt) -- +(1.9pt,0) -- +(0,-1.9pt) -- +(-1.9pt,0) -- cycle;
\filldraw[ochre, draw=ink, line width=0.25pt] (0,1.859) +(0,1.9pt) -- +(1.9pt,0) -- +(0,-1.9pt) -- +(-1.9pt,0) -- cycle;
\filldraw[ochre, draw=ink, line width=0.25pt] (0.26,1.562) +(0,1.9pt) -- +(1.9pt,0) -- +(0,-1.9pt) -- +(-1.9pt,0) -- cycle;
\filldraw[ochre, draw=ink, line width=0.25pt] (1.781,1.065) +(0,1.9pt) -- +(1.9pt,0) -- +(0,-1.9pt) -- +(-1.9pt,0) -- cycle;
\filldraw[ochre, draw=ink, line width=0.25pt] (1.785,1.765) +(0,1.9pt) -- +(1.9pt,0) -- +(0,-1.9pt) -- +(-1.9pt,0) -- cycle;
\filldraw[ochre, draw=ink, line width=0.25pt] (1.83,1.521) +(0,1.9pt) -- +(1.9pt,0) -- +(0,-1.9pt) -- +(-1.9pt,0) -- cycle;
\draw[wine, line width=0.45pt] (1.83,1.521) circle (2.7pt);
\filldraw[ochre, draw=ink, line width=0.25pt] (12.238,0.573) +(0,1.9pt) -- +(1.9pt,0) -- +(0,-1.9pt) -- +(-1.9pt,0) -- cycle;
\filldraw[ochre, draw=ink, line width=0.25pt] (12.032,1.77) +(0,1.9pt) -- +(1.9pt,0) -- +(0,-1.9pt) -- +(-1.9pt,0) -- cycle;
\filldraw[ochre, draw=ink, line width=0.25pt] (12.039,1.521) +(0,1.9pt) -- +(1.9pt,0) -- +(0,-1.9pt) -- +(-1.9pt,0) -- cycle;
\draw[wine, line width=0.45pt] (12.039,1.521) circle (2.7pt);
\filldraw[ochre, draw=ink, line width=0.25pt] (12.085,1.273) +(0,1.9pt) -- +(1.9pt,0) -- +(0,-1.9pt) -- +(-1.9pt,0) -- cycle;
\filldraw[ochre, draw=ink, line width=0.25pt] (12.395,1.029) +(0,1.9pt) -- +(1.9pt,0) -- +(0,-1.9pt) -- +(-1.9pt,0) -- cycle;
\filldraw[ochre, draw=ink, line width=0.25pt] (12.4,1.485) +(0,1.9pt) -- +(1.9pt,0) -- +(0,-1.9pt) -- +(-1.9pt,0) -- cycle;
\node[ink, right=2.2pt] at (6.935,0.119) {$\varepsilon$};
\node[ink, right=2.2pt] at (1.019,0.366) {$a$};
\node[ink, left=2.2pt] at (12.238,0.573) {$b$};
\fill[navy] (0.2,7.42) circle (0.75pt); \node[ink, right] at (0.28,7.42) {central word};
\filldraw[ochre, draw=ink, line width=0.25pt] (2.6,7.42) +(0,1.9pt) -- +(1.9pt,0) -- +(0,-1.9pt) -- +(-1.9pt,0) -- cycle; \node[ink, right] at (2.7,7.42) {centre $p$};
\draw[wine, line width=0.7pt] (4.5,7.42) -- (5.1,7.42); \draw[wine, fill=white, line width=0.45pt] (4.5,7.42) circle (1.1pt) (5.1,7.42) circle (1.1pt); \node[ink, right] at (5.2,7.42) {twins and their rung};
\fill[slate] (8.85,7.42) circle (0.5pt); \node[ink, right] at (8.93,7.42) {other palindrome};
\end{tikzpicture}

\caption{All palindromes with Gram label at most $10^{12}$, placed by the ratio $\rho/M$ and
their Gram label $M$. The ratio $\rho/M=t/(2t+1)$ increases with the slope $t$ of \cref{lem:ladder}, and it is
the frequency of the letter $b$ in $\Omega_z$. Each horizontal segment joins two twins exchanged
around a palindrome $p$ (diamond), at the height of their common label; the dotted post shows
that the segment is centred at the ratio $\rho_p/M_p$ of $p$, which by \cref{lem:ladder} is the
lowest point between its ends. Ladders rise above the empty word, above $a$ and above $b$. The
two ringed diamonds, $abba$ and $baab$, are twins themselves and centres of the exchanges at
$12\,986\,074\,130$, the square of \cref{fig:nested-square}. No central word lies on a rung.}
\label{fig:ladders}
\end{figure}

\Cref{fig:ladders} shows the twins in the plane of the ratio $\rho/M$ and the label. They appear
as the rungs of ladders, which are already visible in the alphabetical order of the directive
words, and the following lemma explains why each rung stands on a palindrome of smaller
label.

\begin{lemma}[Ladders]
\label{lem:ladder}
For a word $w$ let $t(w)$ be the slope of the column $\mu(w)e$. For every word $q$, the slopes
of the words beginning with $q$ lie in the interval $\mu(q)(\mathcal J)$, which contains the
slope of no other word. Consequently, if $Z$ and $Z'$ are twins and $q$ is their longest
common prefix, then every word whose slope lies strictly between $t(Z)$ and $t(Z')$ begins
with $q$, the slope $t(q)$ is one of them, and $q$ has strictly the smallest label among
these words.
\end{lemma}

\begin{proof}
By \cref{sec:decoder}, $A(\mathcal J)$ and $B(\mathcal J)$ lie in $\mathcal J$, on either side
of $2$, and both maps are increasing. A word $qw'$ has slope
$\mu(q)(t(w'))\in\mu(q)(\mathcal J)$. A word that is not of this form either first differs
from $q$ at some letter, after a common prefix $u$, and then its slope lies in
$\mu(u)\mu(c)(\mathcal J)$ while $\mu(q)(\mathcal J)\subset\mu(u)\mu(c')(\mathcal J)$ with
$c\neq c'$, two disjoint intervals; or it is a proper prefix $u$ of $q=uy$, and then its
slope $\mu(u)(2)$ avoids $\mu(u)\mu(y)(\mathcal J)$ because $2\notin A(\mathcal J)\cup B(\mathcal J)$.
This proves the first assertion. Next, for a nonempty word $y=cy'$ the column $\mu(y)e-e$
has positive entries: by induction $\mu(y')e-e$ has nonnegative entries, and then
$\mu(y)e-e=\mu(c)(\mu(y')e-e)+(\mu(c)e-e)$, where $\mu(c)$ has nonnegative entries and
$\mu(c)e-e$ is $(3,2)^{\mathsf T}$ or $(10,4)^{\mathsf T}$. Since the row $e^{\mathsf T}\mu(q)$ has
positive entries, $e^{\mathsf T}\mu(qy)e-e^{\mathsf T}\mu(q)e=e^{\mathsf T}\mu(q)(\mu(y)e-e)>0$.
In particular a twin cannot be a proper prefix of another twin. Hence $Z=q\,c\,Z_1$ and $Z'=q\,c'\,Z_1'$ with $c\neq c'$, say $c=a$. Both slopes
lie in the interval $\mu(q)(\mathcal J)$, and so does everything between them; and
$t(Z)<\mu(q)(2)=t(q)<t(Z')$ because $t(aZ_1)<2<t(bZ_1')$. The words in between begin with $q$,
and the same inequality of labels shows that $q$ is the lowest among them.
\end{proof}

For an exchange pair around $p$ the foot of the rung is $p$ itself, and by
\cref{prop:any-centre} the rung is centred on it in the ratio $\rho/M$, which is an increasing
function of the slope $t$, so that the lemma applies to it unchanged. Since the lemma holds for
every pair of twins, the ladders alone cannot decide the conjecture. In this picture, the
conjecture implies that every palindrome with a twin has at least one rung centred on a
palindrome, and \cref{subsec:centre-decoding} shows that this consequence is equivalent to
part~(i).

\subsection{The midpoint statement}
\label{subsec:centre-decoding}

\Cref{thm:centre-decoder} proves the decoder at every palindromic centre, and
\cref{cor:centre-classification} turns it into a classification: two palindromes with equal
labels whose root midpoint is the ratio $\rho_p/M_p$ of a palindrome $p$ are exchanged around $p$.
This makes part~(i) of \cref{conj:structure} equivalent to one statement about midpoints:
\emph{every palindrome with a twin has a twin whose root midpoint is the ratio of a palindrome.}
Indeed, if a palindrome $Z$ has a twin $Z'$ whose root midpoint is $\rho_p/M_p$, then $Z$ is
structured around $p$ by \cref{cor:centre-classification}; conversely a structured palindrome
has its exchange as a twin with that midpoint, by \cref{prop:any-centre}. No condition on a
common prefix is needed: the corollary forces both palindromes to begin with $p$. Part~(ii) is
equivalent in the same way to the statement that any two twins are joined by a chain of twins
in which consecutive words have their root midpoint at the ratio of a palindrome. This is the
displaced problem of \cref{sec:displaced}, now posed for all palindromes. \Cref{rem:quadruple}
shows that the twin has to be chosen, since some twins are displaced, and the enumeration finds
that every displaced pair is a diagonal of a square of \cref{cor:nested-square}, that is, a
composition of two exchanges. For each pair the centre, when it exists, is unique, because the
ratio of a palindrome determines it.

\subsection{Lengths of directing words}
\label{subsec:lengths}

In the Thue--Morse coding a coincidence of labels is a coincidence of the lengths
$\abs{\operatorname{Pal}(\theta(av))}$ and $\abs{\operatorname{Pal}(\theta(av'))}$ of two central
words, and it is natural to ask whether it forces the two directing words $av$ and $av'$ to have
the same length. The question matters because the directing words of one length form a finite
family, so that a positive answer would allow the coincidences to be studied one length at a
time. This subsection gives the length its meaning in the Cohn coding, shows that
\cref{conj:structure} implies a positive answer, and shows that for Markov words a positive answer
is equivalent to the uniqueness conjecture.

\begin{lemma}[The length of a directing word]
\label{lem:length-cf}
For a word $z$ put $M=e^{\mathsf T}\mu(z)e$ and $\rho=(\mu(z)e)_1$, which are the Gram label and
the Gram root when $z$ is a palindrome, and let $\sigma(z)$ be the sequence of integers obtained
from $z$ by replacing every letter $a$ by $1,1$ and every letter $b$ by $2,2$. Then
$\abs{F_a(z)}=\abs z_a+2\abs z_b+1$, and the regular continued fraction of $M/\rho$ is
$$
 \frac M\rho=[\,2;\ \sigma(z),\ 2\,].
$$
In particular the sum of the partial quotients of $M/\rho$ is $2\abs{F_a(z)}+2$, the directive
$\theta(F_a(z))$ has $2\abs{F_a(z)}$ letters, and $\abs{F_a(z)}=\abs{azb}_b+\abs{azb}-2$.
\end{lemma}

\begin{proof}
The morphism $\phi_a$ sends the letter $a$ to one letter and the letter $b$ to two, which gives the
first formula, and $\theta$ doubles every length. For a positive integer $d$ put
$D(d)=\left(\begin{smallmatrix}d&1\\1&0\end{smallmatrix}\right)$, so that $P_1=D(1)$ and $P=P_2=D(2)$,
with $A=P_1^2$ and $B=P_2^2$. Hence $P\mu(z)P=D(d_0)D(d_1)\cdots D(d_N)$, where
$(d_0,\dots,d_N)=(2,\sigma(z),2)$. Let $K$ denote the continuants, defined by $K()=1$, $K(d_0)=d_0$
and $K(d_0,\dots,d_N)=d_0K(d_1,\dots,d_N)+K(d_2,\dots,d_N)$. The first column of
$D(d_0)\cdots D(d_N)$ is $(K(d_0,\dots,d_N),K(d_1,\dots,d_N))^{\mathsf T}$: for a single factor it is
$(d_N,1)^{\mathsf T}=(K(d_N),K())^{\mathsf T}$, and multiplying on the left by $D(d_0)$ replaces a first
column $(x,y)^{\mathsf T}$ by $(d_0x+y,x)^{\mathsf T}$, which is the recurrence. On the other hand the
first column of $P\mu(z)P$ is $P\mu(z)e$, because $Pe_1=e$; with $\mu(z)e=(x,y)^{\mathsf T}$ it equals
$(2x+y,x)^{\mathsf T}=(M,\rho)^{\mathsf T}$. Therefore $M=K(2,\sigma(z),2)$ and $\rho=K(\sigma(z),2)$.
The quotient of these continuants is the continued fraction: $K(d_0)/K()=d_0$, and
$[d_0;d_1,\dots,d_N]=d_0+1/[d_1;\dots,d_N]=d_0+K(d_2,\dots,d_N)/K(d_1,\dots,d_N)$ by induction, which
is $K(d_0,\dots,d_N)/K(d_1,\dots,d_N)$ by the recurrence. All partial quotients are positive and the
last one equals two, so this is the regular continued fraction of $M/\rho$, which is unique. The pair
$1,1$ adds $2$ to the sum of the partial quotients and the pair $2,2$ adds $4$, so the sum is
$2+2\abs z_a+4\abs z_b+2=2\abs{F_a(z)}+2$. Finally $\abs{azb}_b=\abs z_b+1$ and $\abs{azb}=\abs z+2$.
\end{proof}

The length of the directing word is thus half the sum of the partial quotients of the label divided
by the root, minus one: it counts, with multiplicity, the steps of the Euclidean algorithm on the
pair $(M,\rho)$. For a central word $z$ the identity $M=K(2,\sigma(z),2)$ is the continued-fraction
description of Markov numbers given by Frobenius~\cite[formula~(6), p.~472]{Frobenius1913}; the
proof above does not use centrality. The value of a word in the matrices of Lagisquet, Pelantov\'a, Tavenas and
Vuillon~\cite{LagisquetEtAl2021} connects this length with their ordering of Markov numbers.

\begin{proposition}[Lengths of directing words of Markov words]
\label{prop:lengths-uniqueness}
For a word $w$ let $m(w)=(1,0)\,M^w\,(0,1)^{\mathsf T}$, where $M^w$ is the product along $w$ of
$A_L=\left(\begin{smallmatrix}1&1\\1&2\end{smallmatrix}\right)$ for the letter $a$ and
$B_L=\left(\begin{smallmatrix}2&1\\1&1\end{smallmatrix}\right)$ for the letter $b$; this is the
$m$-value of~\cite{LagisquetEtAl2021}. Then:
\begin{enumerate}[label=\textup{(\roman*)}]
\item $m(a\,F_a(z)\,b)=e^{\mathsf T}\mu(z)e$ for every word $z$;
\item if $z$ is central, then $aF_a(z)b$ is the lower Christoffel word with $p=\abs z_b+1$ letters $b$
 and $q=\abs z+2$ letters $a$, the label $M(z)$ is the Markov number of the Farey fraction $p/q$,
 and $\abs{F_a(z)}=p+q-2$;
\item the Frobenius--Markov uniqueness conjecture holds if and only if any two central words with
 the same label have directing words of the same length; the same holds when the length of the
 directing word is replaced by the number of letters $b$ of the word, or by its length.
\end{enumerate}
\end{proposition}

\begin{proof}
\emph{Part (i).} Put $X=\left(\begin{smallmatrix}1&-1\\0&1\end{smallmatrix}\right)$. Direct
multiplication gives
$$
 XAX^{-1}=\begin{pmatrix}1&1\\1&2\end{pmatrix}=A_L,\qquad
 XBX^{-1}=\begin{pmatrix}3&4\\2&3\end{pmatrix}=B_LA_L .
$$
Since $\phi_a(a)=a$ and $\phi_a(b)=ba$, this says $M^{\phi_a(c)}=X\mu(c)X^{-1}$ for both letters $c$,
and both sides are multiplicative in the word, so $M^{\phi_a(w)}=X\mu(w)X^{-1}$ for every word $w$.
As $aF_a(z)b=a\,a\,\phi_a(z)\,b$,
$$
 m(aF_a(z)b)=(1,0)\,A_L^2X\,\mu(z)\,X^{-1}B_L\binom01 .
$$
Now $(1,0)A_L^2=(2,3)$ and $(2,3)X=(2,1)=e^{\mathsf T}$, while $B_L(0,1)^{\mathsf T}=(1,1)^{\mathsf T}$
and $X^{-1}(1,1)^{\mathsf T}=(2,1)^{\mathsf T}=e$. This proves (i).

\emph{Part (ii).} If $z=\operatorname{Pal}(v)$, then
\eqref{eq:standard-palindrome-image} gives $F_a(z)=\operatorname{Pal}(av)$, a central word, so
$aF_a(z)b$ is a lower Christoffel word~\cite[Section~3]{BertheDeLucaReutenauer2008}. The morphism
$\phi_a$ keeps the number of letters $b$ and adds one letter $a$ for each letter of $z$; hence the
word $aF_a(z)b$ has $\abs z_b+1=p$ letters $b$ and $2+\abs z_a+\abs z_b=\abs z+2=q$ letters $a$, and
$p<q$. A lower Christoffel word is determined by its letter counts, which are coprime, so
$aF_a(z)b$ is the word written $c(p,q)$ in~\cite{LagisquetEtAl2021}, and by part (i) and
\cite[Proposition~3]{LagisquetEtAl2021} its value $M(z)$ is the Markov number of the Farey fraction
$p/q$; we write $m_{p/q}$ for that Markov number. By \cref{lem:length-cf}, $\abs{F_a(z)}=p+q-2$.
Conversely, every reduced fraction $0<p/q<1$ arises from exactly one central $z$. Write
$c(p,q)=a\,w\,b$ with $w$ central; $w$ is not empty, since it has $q-1\geqslant1$ letters $a$. If the
directive of $w$ began with $b$, then $w=F_b(w')=b\,\phi_b(w')$ with $\phi_b(a)=ab$, $\phi_b(b)=b$:
every letter $a$ of $w$ would be followed by a letter $b$, and $w$ would have more letters $b$ than
letters $a$. Since $w$ has $q-1>p-1$ letters $a$, its directive begins with $a$, so $w=F_a(z)$ for a
central $z$, unique because $F_a$ is injective.

\emph{Part (iii).} The uniqueness conjecture is equivalent to the injectivity of $p/q\mapsto m_{p/q}$
on the reduced fractions of $(0,1)$, since the Markov tree indexes by these fractions the
normalized Markov triples with distinct entries; by part (ii) this is the injectivity of the label
on central words, and it makes the three statements vacuously true. Conversely suppose that two
distinct central words have the same label, and let $p/q\neq p'/q'$ be their fractions. Write
$m_{n,d}$ for the least value of $m$ on the lattice paths from $(0,0)$ to $(d,n)$; by
\cite[Proposition~21]{LagisquetEtAl2021} it is attained at $c(n,d)$, so $m_{p,q}$ and
$m_{p',q'}$ are the two equal labels. By \cite[Theorem~2]{LagisquetEtAl2021}, for integers $j<k$,
$$
 m_{j,k}<m_{j,k+1},\qquad m_{j,k}<m_{j+1,k},\qquad m_{j+1,k}<m_{j,k+1}.
$$
If $p+q=p'+q'=s$, then $p\neq p'$, say $p'<p$; for $p'\leqslant j\leqslant p-1$ one has
$j<q-1<s-j-1$, so the third inequality applied to $j<s-j-1$ gives $m_{j+1,s-j-1}<m_{j,s-j}$, and
chaining from $j=p-1$ down to $j=p'$ gives $m_{p,q}<m_{p',q'}$, a contradiction. If $p=p'$, then
$q\neq q'$, say $q<q'$, and the first inequality applied to $p<k$ for $q\leqslant k<q'$ gives
$m_{p,q}<m_{p,q'}$. If $q=q'$, then $p\neq p'$, say $p<p'$, and the second inequality applied to
$j<q$ for $p\leqslant j<p'$ gives $m_{p,q}<m_{p',q}$. Each case contradicts the equality of the
labels, so a coincidence of labels among distinct central words changes the length of the
directing word, the number of letters $b$, and the length. Each of the three statements therefore
implies the uniqueness conjecture.
\end{proof}

Nothing in the proposition is deep, and most of it is known. Part~(i) is, in our normalization,
the identity $g(w)=m(aG(w)ab)$ of~\cite[Appendix~B]{LagisquetEtAl2021}, where $G(a)=a$ and $G(b)=ab$,
so that $aG(w)ab=aF_a(w)b$, and where $g$ is the function of Aigner's book, which is the label
$e^{\mathsf T}\mu(w)e$ written in another basis. Part~(iii) is an immediate consequence of the
ordering theorems, and Lee, Li, Rabideau and Schiffler prove more: $m_{p/q}=m_{p'/q'}$ with
$(p,q)\neq(p',q')$ forces $-5/4<(p-p')/(q-q')<-8/7$~\cite[Corollary~1.5]{LeeEtAl2023}. Since
$\abs{F_a(z)}=p+q-2$, a coincidence would change the length of the directing word by more than a
seventh and less than a quarter of the change of $q$, and never by zero. The point of the proposition
is the identification it makes: the directing word of the Thue--Morse coding is the central factor of
the Christoffel word of the fixed-sum theorem, so that the length question is an equivalent form of
the uniqueness conjecture.

Three consequences place the proposition. First, part~(ii) of \cref{conj:structure} implies that
twins have the same numbers of letters $a$ and $b$, because an exchange does not change them; so the
conjecture implies that twins have directing words of the same length, and by the proposition this
alone already implies uniqueness. Second, the equality of letter counts is a theorem on every
reflection axis: by \cref{cor:centre-classification}, two twins whose root midpoint is the ratio of
a palindrome are exchanged around it, so they have the same numbers of letters $a$ and $b$ and
directing words of the same length; in the enumeration up to $10^{26}$ every class of twins has a
single pair of letter counts. Third, in the language
of~\cite{LagisquetEtAl2021}, part~(ii) says that their Christoffel word $c(p,q)$ is $aF_a(z)b$: the
directing word of the Thue--Morse coding is its central factor, two letters shorter. Their fixed-sum
theorem says that two distinct lower Christoffel words with more letters $a$ than letters $b$ and
of the same length carry different Markov numbers; the length question asks for
the complementary statement, that a coincidence of Markov numbers keeps the length, and the two
statements together are the uniqueness conjecture.

\subsection{Off the axis}

Moving the midpoint breaks the reflection in an explicit way. The two matrix factors of
\cref{prop:displaced-equation} remain involutions, but they differ, and the left factor moves
the terminal column by a rank-one displacement. The integer defect and the parent-sensitive
bound of \cref{thm:displaced-ancestor-bound} give coordinates for this remaining problem: they
show that a balanced common ancestor of two distinct central twins must be small compared with
their label. They do not bound the two directive tails or prevent early branching. A proof of
uniqueness along these lines would have to show that central words cannot realize a displaced
midpoint, and centrality must enter there, since displaced twins exist among palindromes. The
displaced twins found by the enumeration are all diagonals of the squares of
\cref{cor:nested-square}: each pair is a composition of two exchanges, and its root is translated by
twice the label times the difference of the two centre ratios.
